\documentclass[11pt]{article}

\usepackage[margin=1in]{geometry}
\usepackage{amsmath,amssymb,amsthm}
\usepackage{array,booktabs}
\usepackage{enumitem}
\usepackage[hidelinks]{hyperref}

\theoremstyle{plain}
\newtheorem{proposition}{Proposition}
\newtheorem{lemma}{Lemma}
\newtheorem{corollary}{Corollary}
\newtheorem{conjecture}{Conjecture}
\theoremstyle{definition}
\newtheorem{remark}{Remark}

\newcommand{\F}{\mathbb{F}}
\newcommand{\Q}{\mathbb{Q}}
\newcommand{\N}{\mathbb{N}}
\newcommand{\ord}{\operatorname{ord}}
\newcommand{\Norm}{\operatorname{N}}
\newcommand{\Frob}{\operatorname{Frob}}
\newcommand{\nimmul}{\otimes}
\newcommand{\Qset}{\mathcal{Q}}
% claim-level tags (repo convention for this thread)
\newcommand{\PROVED}{\textbf{[\textsc{proved}]}}
\newcommand{\CERTIFIED}[1]{\textbf{[\textsc{certified} $#1$]}}
\newcommand{\CONSISTENT}{\textbf{[\textsc{consistent}]}}
\newcommand{\CONJECTURED}{\textbf{[\textsc{conjectured}]}}
\newcommand{\OPEN}{\textbf{[\textsc{open}]}}

\title{Draft: the finite excess in transfinite nim multiplication}
\author{a9lim}
\date{July 2026}

\begin{document}
\maketitle

\begin{abstract}
Conway's field $\mathrm{On}_2$ makes the ordinals below
$\omega^{\omega^{\omega}}$ an algebraic closure of $\F_2$ under
nim-arithmetic. Its multiplication is governed by a Kummer tower whose
carries are Lenstra's \emph{excess} values: for each odd prime $p$, the
carry is $\alpha_p=\kappa_{f(p)}+m_p$, where the transfinite part
$\kappa_{f(p)}$ has a known closed shape (Lenstra, DiMuro) but the finite
correction $m_p$ does not. This note consolidates the state of the excess
problem (\texttt{docs/OPEN.md} Problem~3 of \texttt{ogdoad}): the exact
finite-field reformulation of the root test as a multiplicative-order
criterion and its structural norm reduction; a complete analysis of the
$3$-power column via a half-angle splitting in the cyclotomic component
fields, certifying $m_r=1$ for all primes $r$ with
$f(r)\in\{3,9,\dots,3^6\}$ and forcing $m_p\geq4$ unconditionally on the
$f(p)=2\cdot3^k$ column --- now sharpened to $m_p=4$ \emph{exactly} at
every prime current factor tables reach (universally for $k\leq6$), via
a corrected compositum norm
$(\kappa+4)(\kappa+5)=\kappa^2+\kappa+\omega$ that collapses the $m=4$
test into the same trinomial field, together with a correction of the
formerly cumulative $D_k$ to the exact current factor $D'_k$ and an
all-level Capelli/resultant formulation, an exact-degree sparse
projection, and a dominant-factor sieve; the empirically exact $0/1/4$
candidate rule and the proof that the Lenstra component set alone cannot
decide the excess; an exact transverse-norm formulation of every
singleton-odd row; two embedded maximal-order problems (the Conway
Fermat tower on the even side and explicit Gaussian-period primitivity on
the $3$-power side), with the latter now reduced to primitivity of the
irreducible Dickson iterate $D_{3^k}(X,1)+1$ and to an explicit recursive
Singer-generator problem in a cubic norm-one torus, proved analytically
through the first three prime-order kernel steps, equipped here with an
unconditional partition-order bound exponential in the square root of
the level and an exact short cyclotomic-period test, and reduced thereafter
to character nonvanishing for a norm-coherent cyclotomic unit; a new
Fibonacci-polynomial stratification of the Conway--Fermat arm, including
universal countermodels proving that degree, trace, and norm cannot
force maximal quotient order, and an exact resultant recursion for the
remaining distinguished-factor conjecture, now realized as a Kummer
symbol at a distinguished point of a fixed supersingular elliptic
correspondence, together with an exact full-stratum-to-proper-stratum
counterexample proving that one-step resultant induction does not preserve
maximality for arbitrary predecessor factors, and a complete trace/norm
ancestry criterion whose exact power trace-fibres collapse back to the same
fibotomic test; and the fact that
generic bounded-translate arguments cannot prove the conjecture, by a
finite-window Kummer-cover theorem producing arbitrarily large exact-degree
elements whose entire fixed translate window consists of $p$-th powers.
The three distinguished arms are sharpened further: cubic failure saturates
every cyclotomic secant although the natural Kummer cover is rational, while
Stickelberger and Hilbert reciprocity are proved to vanish on precisely the
simple exceptional eigenspace carrying its local Artin value; the
Fermat obstruction is the ramified toric projection of a Miller unit in a
three-point generalized Jacobian, whose complete extension calculation now
has explicit Rosenlicht/Miller and rational five-torsion coordinates while
proving their evident reciprocity collapse circular; and $D'_k$ is equivalent both to an explicit
translated-binomial irreducibility statement over $\F_{16}$ and to a relative
Dickson-composition criterion with a proved degree ceiling, while one bad
translate saturates full six- and eight-dimensional additive spaces and the
complete semiprimitive Gauss--Jacobi/Fourier table is proved compatible with
that failure, now sharpened to one extremal second Dirichlet coefficient
whose functional equation, trace dual, and exact subspace polynomial are
all proved compatible with the bad value.  Finally,
proving the candidate would settle
Lenstra's open boundedness question, while already requiring those two
hard maximal-order inputs.
Feasibility of the first open row, $p=719$, is assessed: a
Frobenius-orbit norm recurrence reduces it to arithmetic in
$\F_{2^{359}}$, with the cost concentrated in tower-aware Frobenius
composition. Each statement carries an explicit claim-level tag.
\end{abstract}

\paragraph{Claim levels.}
\PROVED{} -- proved here, machine-checked, or classically verified;
\CERTIFIED{4\leq k\leq6} -- proved for the stated range, blocked beyond it only
by unfactored cofactors; \CONSISTENT{} -- no counterexample found, not
proved; \CONJECTURED{} -- explicitly open in the literature or posed in
this thread; \OPEN{} -- neither confirmed nor refuted, no feasible attack
on record.

\section{Notation, implemented state, external data}\label{sec:notation}

\subsection{Notation}

For an odd prime $p$:
\begin{itemize}[leftmargin=*,itemsep=0pt]
\item $f(p)=\ord_p(2)$, the multiplicative order of $2$ modulo $p$;
\item $\kappa_h$ is the Lenstra/DiMuro tower element indexed by $h$;
\item $\Qset(h)$ is Lenstra's set of prime-power components appearing in
$\kappa_h$;
\item the \emph{excess} $m_p$ is the least finite $m$ such that
$\kappa_{f(p)}+m$ has no $p$-th root in the relevant finite component
field;
\item the Kummer carry is $\alpha_p=\kappa_{f(p)}+m_p$.
\end{itemize}
For the ordinal tower, a row with $\Qset(f(p))=\{q\}$ and finite excess
$m$ gives the ordinal sum corresponding to $\kappa_q+m$. Addition of
nimbers is XOR; $\nimmul$ is the nim product.

\subsection{Implemented tower state}

In the Rust tower (\path{src/scalar/big/ordinal/tower.rs}), as of 2026-06-13:
\begin{itemize}[leftmargin=*,itemsep=0pt]
\item the finite excess rows $m_p$ are now sourced from OEIS
A380496~\cite{OEIS}, the b-file's $126$ known rows (odd primes
$3 \le p \le 709$); the first $14$ rows reproduce DiMuro Table~1 + the
formerly locally verified $\alpha_{47}=\omega^{(\omega^7)}+1$;
\item the operational boundary is $\alpha_{719}$, the first OEIS-unknown
row: a carry needing $\alpha_{719}$ or beyond returns \texttt{None};
\item the table extends \emph{reach}, not \emph{feasibility}: each
$\alpha_p$ is reconstructed in code from $\ord_p(2)$, $\Qset(f(p))$, and
$m_p$, but for large $p$ the $\Qset$/finite-subfield reconstruction lives
in the degree-$e_p$ component field ($e_p$ in the millions; see the
snapshot below), so only the small-$e_p$ rows are materializable
end-to-end today.
\end{itemize}

\subsection{External data snapshot (2026-06-09)}

\begin{itemize}[leftmargin=*,itemsep=0pt]
\item OEIS A380496~\cite{OEIS} has $1417$ extended rows: $799$ known and
$618$ unknown; the b-file has $126$ initial known rows.
\item The first OEIS unknown is row $n=127$, the $127$th odd prime
$p=719$. For $p=719$: $f(719)=\ord_{719}(2)=359$ and
$\Qset(359)=\{359\}$.
\item The transfinite-nim-calculator logs~\cite{Peeters} record the direct
component exponent as $e_{719}=1{,}258{,}230{,}380$, the practical wall
for direct exponentiation.
\end{itemize}

\section{The order criterion and structural norm reduction}
\label{sec:norm}

\subsection{The exact criterion}

Throughout, $E$ denotes the degree of the component field containing
$\beta=\kappa_{f(p)}+m$, so $f=f(p)\mid E$ and hence
$p\mid 2^{E}-1$.

\begin{proposition}[power criterion]\label{prop:power}
\PROVED{} $\beta\in\F_{2^E}^{\times}$ has no $p$-th root in $\F_{2^E}$
iff $\beta^{(2^E-1)/p}\neq1$. Hence
\[
  m_p=\min\bigl\{m\geq0:\ (\kappa_{f(p)}+m)^{(2^E-1)/p}\neq1\bigr\}.
\]
\end{proposition}

\begin{proof}
$\F_{2^E}^{\times}$ is cyclic of order $2^E-1$ and $p\mid 2^E-1$; in a
cyclic group of order $N$, $\beta$ is a $p$-th power iff
$\beta^{N/p}=1$.
\end{proof}

\begin{corollary}[order form]\label{cor:order}
\PROVED{} If $v_p(2^E-1)=1$, then $\beta$ has no $p$-th root in
$\F_{2^E}$ iff $p\mid\ord(\beta)$, so
$m_p$ is the least $m$ with $p\mid\ord(\kappa_{f(p)}+m)$.
\end{corollary}

\begin{remark}[Wieferich caveat]\label{rem:wieferich}
In general ``no $p$-th root'' is equivalent to
$v_p(\ord\beta)=v_p(2^E-1)$, which collapses to
$p\mid\ord(\beta)$ only when $v_p(2^E-1)=1$. By lifting the exponent,
\[
v_p(2^E-1)=v_p(2^{f(p)}-1)+v_p(E/f(p)).
\]
Thus the hypothesis can fail either at base-$2$ Wieferich-type primes
with $p^2\mid 2^{f(p)}-1$ or when $p\mid E/f(p)$.  The only known
base-$2$ Wieferich primes are $p=1093$ and $p=3511$, both of which lie
inside the extended A380496 range. For those rows, and for any row with
$p\mid E/f(p)$, the exact test of Proposition~\ref{prop:power} (or the
full-$p$-part order condition) must be used. The certifications of
Section~\ref{sec:family} compute orders exactly, with squarefreeness
checks, and are unaffected. The research-note predecessor of this
document stated only the order form; the power criterion is the
correct general statement.
\end{remark}

\subsection{Norm reduction}

\begin{proposition}[norm reduction]\label{prop:norm}
\PROVED{} For $\beta\in\F_{2^E}^{\times}$ with $f=f(p)\mid E$,
\[
  \beta^{(2^E-1)/p}
  =
  \Norm_{\F_{2^E}/\F_{2^f}}(\beta)^{(2^f-1)/p}.
\]
\end{proposition}

\begin{proof}
$\Norm_{\F_{2^E}/\F_{2^f}}(\beta)=\beta^{(2^E-1)/(2^f-1)}$, and
$\bigl((2^E-1)/(2^f-1)\bigr)\cdot\bigl((2^f-1)/p\bigr)=(2^E-1)/p$.
\end{proof}

So the $p=719$ test reduces to: (1) compute
$N=\Norm_{\F_{2^E}/\F_{2^{359}}}(\kappa_{359}+1)$ in $\F_{2^{359}}$;
(2) test whether $719\mid\ord(N)$ in $\F_{2^{359}}^{\times}$, a field of
degree $359$ in which $719\mid2^{359}-1$.

\begin{proposition}[Galois-conjugacy invariance]\label{prop:galois}
\PROVED{} The order of $\kappa_{f(p)}+m$ is constant on its Frobenius
orbit in $\F_{2^E}$, and the norm
$N=\prod_{i=0}^{E/f-1}\Frob^{if}(\kappa_{f}+m)$ is determined
structurally by the tower equations without materialising the full orbit.
\end{proposition}

\begin{proposition}[tower pinning]\label{prop:pinning}
\PROVED{} The tower equations pin $\kappa_{359}$ up to Frobenius
conjugacy: each composed minimal polynomial along the component chain
$359\to179\to89\to11\to5\to\text{finite}$ is irreducible (the
non-$q$-th-power property defining the lower excess guarantees this), so
$\kappa_{359}$ is a well-defined conjugacy class and $N$ is a well-defined
element of $\F_{2^{359}}$.
\end{proposition}

\subsection{Exact support and the general singleton-odd obstruction}

For $h>1$, write
\[
  \Pi_h:=\prod_{\substack{\ell\ \mathrm{odd\ prime}\\
                    \ord_\ell(2)=h}}
          \ell^{v_\ell(2^h-1)}.
\]
This is the full prime-power part of $2^h-1$ supported at primes whose
order of $2$ is exactly $h$; using the full valuations, rather than only
the radical, incorporates the Wieferich cases.

\begin{proposition}[exact-order support]\label{prop:support}
\PROVED{} Let $x=\kappa_h$, let $E=d(x)$, and put
\[
  T_h=\Norm_{\F_{2^E}/\F_{2^h}}(x).
\]
Then $m_p=0$ for every odd prime $p$ with $f(p)=h$ if and only if
\[
  \Pi_h\mid\ord(T_h).
\]
For an individual such $p$, $m_p=0$ iff
$v_p(\ord T_h)=v_p(2^h-1)$.
\end{proposition}

\begin{proof}
The definition of $\kappa_h$ gives $h\mid E$.  Proposition~\ref{prop:norm}
with $\beta=x$ identifies the $p$-power test in $\F_{2^E}$ with the
same test on $T_h$ in $\F_{2^h}$.  In the cyclic group
$\F_{2^h}^{\times}$, failure to be a $p$-th power is exactly the presence
of the full $p$-primary part of the group order in $\ord(T_h)$.  Taking
the product over the distinct eligible primes proves the claim.  No
hypothesis on $p\nmid E/h$ is hidden here: if
$c=v_p((2^E-1)/(2^h-1))$, the norm exponent gives
$v_p(\ord T_h)=\max\{0,v_p(\ord x)-c\}$, so equality with
$v_p(2^h-1)$ is equivalent to $v_p(\ord x)=v_p(2^E-1)$.
\end{proof}

The singleton-odd case admits a sharper field square.  Let $r$ be an odd
prime, $q=r^a$, and set $b_r=d(\alpha_r)$, where
$\kappa_q^q=\alpha_r$.  Since $\alpha_r<\kappa_r$, Lenstra's field
description gives $r\nmid b_r$, while the successive Kummer equations
have degree $r$~\cite[Proposition~1.8]{Lenstra77}; hence
\[
  d(\kappa_q)=qb_r,\qquad \gcd(q,b_r)=1.
\]

\begin{proposition}[singleton transverse-norm criterion]
\label{prop:singleton-transverse}
\PROVED{} Suppose $\Qset(h)=\{q\}$ with $q=r^a$ odd.  Then
$h=qs$ for a divisor $s\mid b_r$.  In the field square
\[
  F=\F_{2^{qb_r}},\quad B=\F_{2^{b_r}},\quad
  K=\F_{2^{qs}},\quad C=\F_{2^s},
\]
one has $B\cap K=C$ and $F=BK$.  Put $x=\kappa_q$ and
\[
  \Theta_{q,s}:=\Norm_{F/K}(x+1)\in K^{\times}.
\]
For every prime $p$ with $f(p)=h$:
\begin{enumerate}[label=(\roman*),leftmargin=*]
\item $x$ is a $p$-th power in $F$, so $m_p\geq1$;
\item $m_p=1$ iff
\[
  v_p(\ord\Theta_{q,s})=v_p(2^{qs}-1);
\]
\item the easy Kummer norm is
\[
  \Norm_{F/B}(x+1)=1+\alpha_r,
\]
but $p\nmid |B^{\times}|$, so this norm cannot see the required
$p$-primary obstruction.  The norm to $K$ is genuinely transverse.
\end{enumerate}
\end{proposition}

\begin{proof}
Lenstra's component theorem gives
$\kappa_h=\sum_{t\in\Qset(h)}\kappa_t$ and, for each component $t$,
$t\mid h$ and $\gcd(t,h/t)=1$
\cite[Corollary~2.2]{Lenstra77}.  Thus the singleton assumption identifies
$\kappa_h=\kappa_q$.  Since $h\mid d(\kappa_h)=d(\kappa_q)=qb_r$, one
has $h=qs$ with $s\mid b_r$.  The
intersection and compositum statements follow from the gcd/lcm rules for
finite subfields and $\gcd(q,b_r)=1$.

Since $p>h\geq r$, one has $p\nmid q$.  Also $f(p)=qs\nmid b_r$, so
$p\nmid2^{b_r}-1$.  The relation $x^q=\alpha_r\in B$ therefore gives
$p\nmid\ord(x)$, proving (i), including the full $p$-primary statement.
Part (ii) is Proposition~\ref{prop:norm} applied to $x+1$.  Finally,
$X^q+\alpha_r$ is the minimal polynomial of $x$ over $B$, and evaluating
it at $1$ gives (iii); the invisibility follows again from
$p\nmid2^{b_r}-1$.
\end{proof}

Thus the singleton-one arm is not an unspecified root search.  It is the
uniform assertion that the explicitly defined transverse norm
$\Theta_{q,s}$ contains the full $p$-primary part for every prime
$p$ of exact order $qs$.  The $C_k$ norms below are the special case
$q=3^k$, $b_r=2$, $s=1$; the $m_{359}$ obstruction is the same square at
$q=179$, $b_r=19580$, $s=1$.
Lenstra's use of Carlitz's primitive-root distribution theorem gives an
effective bound at which \emph{some} finite translate works
\cite{Carlitz53,Lenstra77}; it does not evaluate this fixed first
transverse norm.

There is a sharp reason that this effective argument cannot be replaced,
for arbitrary finite-field elements, by any fixed initial translate
window.

\begin{proposition}[finite-window Kummer no-go]
\label{prop:translate-no-go}
\PROVED{} Let $p$ be an odd prime and let

\[
 A=\{b_1,\ldots,b_r\}\subset\overline{\F}_2
\]

be a finite set of distinct elements.  Choose
$k_0=\F_{2^d}$ containing both $A$ and $\mu_p$.  There are
arbitrarily large integers $n$ and elements
$x_n\in\F_{2^{dn}}$ of exact degree $dn$ over $\F_2$ such that

\[
 x_n+b_i\in(\F_{2^{dn}}^\times)^p
 \qquad(1\leq i\leq r).
\]

Consequently no absolute truncation of Lenstra's translate-generation
lemma, nor any argument using only full degree, an order-$p$
character, and a fixed affine translate set, can prove an absolute
bound for arbitrary elements.  This does not decide the distinguished
elements \(\kappa_{f(p)}\): any uniform excess theorem must use their
special Conway-tower identities.
\end{proposition}

\begin{proof}
Over $k_0(X)$ form the Kummer extension

\[
 \mathcal K=k_0(X)(Y_1,\ldots,Y_r),
 \qquad Y_i^p=X+b_i.
\]

The classes of the $X+b_i$ are independent in
$\overline{k}_0(X)^\times/
 \overline{k}_0(X)^{\times p}$: in a relation
$\prod_i(X+b_i)^{e_i}=G^p$, $0\leq e_i<p$, the valuation at
$X=-b_i$ forces $e_i=0$.  Thus the smooth projective curve $C$
with function field \(\mathcal K\) is geometrically integral and

\[
 [\mathcal K:k_0(X)]=p^r,
 \qquad \operatorname{Gal}(\mathcal K/k_0(X))
       \cong(\mathbb Z/p\mathbb Z)^r.
\]

The cover $C\to\mathbb P^1_X$ is ramified at the $r$ points
$-b_i$ and at infinity, each with inertia order $p$.  Hence
Riemann--Hurwitz gives

\[
 2g(C)-2
 =-2p^r+(r+1)p^r\left(1-\frac1p\right).
\]

Put $Q=2^d$.  Hasse--Weil, followed by deletion of the finitely many
branch points, supplies

\[
 Q^n+O_{p,r}(Q^{n/2})
\]

unramified $\F_{Q^n}$-points.  Since the map to the $X$-line has
degree $p^r$, the number of resulting $X$-coordinates is at least

\[
 \frac{2^{dn}-O_{p,r}(2^{dn/2})}{p^r}.
\]

The union of all proper subfields of \(\F_{2^{dn}}\) has at most

\[
 \sum_{\substack{e\mid dn\\e<dn}}2^e
 \leq dn\,2^{dn/2}
\]

elements.  For all sufficiently large $n$ the former lower bound is
larger.  Choose an unramified $X$-coordinate $x_n$ outside that
union.  It has exact degree $dn$, and its lifted coordinates satisfy
$x_n+b_i=Y_i^p\ne0$ for every $i$, as required.
\end{proof}

\subsection{Feasibility for $p=719$}

The component field degree is
$E=2\cdot2\cdot5\cdot11\cdot89\cdot179\cdot359=1{,}258{,}230{,}380$; the
norm requires $E/f=3{,}504{,}820$ Frobenius steps.

\CONSISTENT{} The norm can in principle be computed by the
binary-splitting Frobenius recurrence
$N_{2k}=N_k\cdot\Frob^{fk}(N_k)$ (iterated-Frobenius technique, von zur
Gathen--Shoup~\cite{vzGS92}, Kaltofen--Shoup~\cite{KS98}), requiring about
$\log_2(3{,}504{,}820)\approx22$ squarings-with-Frobenius-composition at
degree $E$ over $\F_2$. Each Frobenius composition is a modular
composition in $\F_2[x]$ modulo a degree-$E$ polynomial; with dense
polynomial arithmetic at $E\approx1.26\times10^9$ this is not locally
feasible.

\CONJECTURED{} (as a cost model, not a theorem) Kummer-tower-aware
Frobenius -- standard-lattice arithmetic for compatibly embedded finite
fields in the style of De~Feo--Randriam--Rousseau~\cite{DFRR19}, where
$q$-power Frobenius acts near-monomially per tower level -- is the likely
$10$--$100\times$ lever: the sparse factored structure
$E=2\cdot2\cdot5\cdot11\cdot89\cdot179\cdot359$ admits a stepwise
Frobenius representation in which each level contributes a small-degree
modular composition. A \emph{measured} cost model, not a run, is the
correct first deliverable.

\OPEN{} Whether the Kummer-tower Frobenius model makes the $m_{719}$
computation feasible in wall-clock hours rather than months on available
hardware (a single consumer GPU; GF(2) arithmetic is XOR-sliced and
GPU-friendly).

\PROVED{} (2026-06-14 local rehearsal.) The first two dependency rows for
the $p=719$ norm path are locally certified by
\path{experiments/ordinal_excess_probe.py}: $m_{89}=1$ in the
$E=220$ component field, and $m_{179}=1$ in the $E=19{,}580$ component
field via the fixed-base power path (\texttt{--deep}, about one minute
locally). The remaining local rehearsal row is $m_{359}=1$, with
$E=3{,}504{,}820$. This row is already source-pinned by the imported
A380496 table ($359\leq709$); what remains open here is an independent
certificate on the same code path needed for the genuinely unknown
$p=719$ row.

\subsection{\texorpdfstring{The $m_{359}$ rehearsal obstruction}{The m359 rehearsal obstruction}}
\label{sec:m359-obstruction}

Set
\[
  d=19{,}580,\qquad E=179d=3{,}504{,}820,\qquad
  B=\F_{2^d},\qquad F=\F_{2^E},\qquad L=\F_{2^{179}}.
\]
For the $m_{359}=1$ rehearsal, $f(359)=179$,
$\Qset(f(359))=\Qset(179)=\{179\}$, and the candidate element is
$\beta=\kappa_{179}+1$. To represent $\kappa_{179}$, the tower uses the
lower prime row $f(179)=178$, $\Qset(178)=\{89\}$, $m_{179}=1$.
Thus in the tower model $F=B(x)$ with $x=\kappa_{179}$ and
\[
  x^{179}=\kappa_{89}+1\in B,
\]
because the lower row $m_{179}=1$ has already been certified.

\PROVED{} The cheap Kummer norm down the top tower step is not the
$359$-test norm. Indeed,
\[
  \Norm_{F/B}(x+1)=1+(\kappa_{89}+1)=\kappa_{89}
\]
(evaluate the monic polynomial $T^{179}+(\kappa_{89}+1)$ at $T=1$, in
characteristic $2$). But $359\nmid |B^\times|=2^{19580}-1$, since
$\ord_{359}(2)=179$ and $179\nmid19580$. Thus the structurally cheap
norm lands in exactly the field where the $359$-primary obstruction is
invisible.

\PROVED{} The norm forced by Proposition~\ref{prop:norm} is the
transverse one
\[
  \Norm_{F/L}(\beta)=\prod_{i=0}^{19579}\Frob^{179i}(\beta)\in L.
\]
Here $\gcd(179,19580)=1$, so $B\cap L=\F_2$ and
$F=BL$. The Frobenius $\Frob^{179}$ preserves $B$ as a set but acts
nontrivially on it; it is not the relative Kummer action
$x\mapsto \zeta x$ over $B$. Consequently the known lower certificates
for $m_{89}$ and $m_{179}$ do not propagate to $m_{359}$ by taking the
easy top-step norm.

\PROVED{} The Wieferich caveat is absent at the next two pressure
points: direct modular checks give
$2^{179}\not\equiv1\pmod{359^2}$ and
$2^{359}\not\equiv1\pmod{719^2}$. The order form of
Corollary~\ref{cor:order} is therefore equivalent to the full power
criterion for the $m_{359}$ rehearsal and the proposed $m_{719}$ test.

\CONSISTENT{} (2026-06-16 diagnostic.) Re-running the maintained probe
confirmed the quick certificate table and the \texttt{--deep}
$m_{179}=1$ certificate. A side norm experiment then computed the
correct transverse norm in the current term basis for the feasible
analogues: for $p=89$ the norm to $\F_{2^{11}}$ has support
$111/220$, and for $p=179$ the norm to $\F_{2^{178}}$ has support
$9691/19580$. Thus the target-subfield element is essentially half-dense
in this representation; using it as the new fixed base makes the final
root-test exponent slower than the direct fixed-base certificate. This
is a representation-level diagnostic, not a theorem that no sparse
description exists.

\OPEN{} The current pure-Python term algebra has no feasible route to a
local $m_{359}$ certificate. The direct fixed-base path would have to
precompute multiplication by $\beta$ for $E=3{,}504{,}820$ basis terms
and then run an $E$-bit exponent. The norm path lowers the mathematical
problem to a degree-$179$ target field, but computing that norm still
requires the $19580$-term transverse Frobenius orbit through
$F_{2^{3504820}}$. A practical certificate now appears to require either
(i) dense/sliced GF(2) arithmetic such as \texttt{gf2x}/NTL, or
(ii) a tower-aware Frobenius representation that makes the transverse
orbit cheap. This is the obstruction hit in the present pass.

\begin{remark}[shared abstraction]\label{rem:bridgek}
The structural primitive needed here --
$\Norm_{E/K}(\beta)=\prod_i\Frob^{i}(\beta)$ computed by a
Frobenius-orbit recurrence -- is the same primitive that the cyclic-algebra
Brauer-invariant bridge (Bridge~K, \path{docs/TBD.md}) requires for
reduced norms over a Frobenius-generated cyclic extension. The existing
\texttt{FieldExtension::relative\_norm} does not apply to the
term-algebra/transfinite setting, but the shape is identical; factoring
out a reusable \path{relative_norm_over_frobenius_orbit} is an
engineering opportunity. This is not a claim that the bounded
\texttt{Fpn} norm certifies $m_{719}$.
\end{remark}

\section{The $3^k$ family}\label{sec:family}

Write $\zeta:=\kappa_{3^k}$ and $h:=3^k$ throughout this section.

\subsection{Structure}

\begin{proposition}[cyclotomic recognition]\label{prop:cyclo}
\PROVED{} $\zeta^{3^k}=\kappa_2$, which has order $3$; hence $\zeta$ is a
primitive $3^{k+1}$-st root of unity. Since $2$ is a primitive root
modulo $3^{k+1}$, the cyclotomic polynomial
$\Phi_{3^{k+1}}(x)=x^{2h}+x^h+1$ is irreducible over $\F_2$, the
component field is $\F_2(\zeta)=\F_{2^{2h}}$, and it carries the index-$2$
subfield $L=\F_{2^h}$.
\end{proposition}

\begin{proposition}[half-angle splitting]\label{prop:halfangle}
\PROVED{} With $s=(3^{k+1}+1)/2$ (the inverse of $2$ modulo $3^{k+1}$),
\[
  \kappa_{3^k}+1=\zeta^{s}\,\bigl(\zeta^{s}+\zeta^{-s}\bigr),
\]
where $\zeta^{s}$ lies in the norm-one circle $U$ (of order exactly
$3^{k+1}$ within the subgroup of order $2^h+1$) and
$(\zeta^{s}+\zeta^{-s})^2=\zeta+\zeta^{-1}=:\gamma_k\in L^{\times}$.
Since $\F_{2^{2h}}^{\times}=L^{\times}\times U$ with coprime orders,
\[
  \ord(\kappa_{3^k}+1)=3^{k+1}\cdot\ord(\gamma_k),
  \qquad
  \ord(\gamma_k)\mid 2^{3^k}-1.
\]
Machine-checked corollaries: $(\kappa+1)^{2^h-1}=\zeta^{-1}$ and
$\Norm_{F/L}(\kappa+1)=\gamma_k$.
\end{proposition}

\begin{proposition}[translates $m=2,3$]\label{prop:translates}
\PROVED{} $\kappa_{3^k}+2$ and $\kappa_{3^k}+3$ split as in
Proposition~\ref{prop:halfangle} with $\gamma_k$ replaced by a Galois
conjugate. Therefore
$\ord(\kappa_{3^k}+m)=3^{k+1}\cdot\ord(\gamma_k)$ for $m=1,2,3$, and
$\ord(\kappa_{3^k})=3^{k+1}$.
\end{proposition}

\begin{proposition}[the $2\cdot3^k$ exception, unconditionally]
\label{prop:exception}
\PROVED{} Any prime $p$ with $f(p)=2\cdot3^k$ satisfies
$p\mid 2^{3^k}+1$, hence $p$ divides neither $3^{k+1}$ nor $2^{3^k}-1$,
hence $p$ never divides $\ord(\kappa_{3^k}+m)$ for $m\in\{0,1,2,3\}$. By
the order criterion, $m_p\geq4$. The full primitivity conjecture $C_k$
below is not needed; the splitting alone forces the bound. The argument
applies beyond the current tables, e.g.\ $p=87211$, $f=54$.
\end{proposition}

\begin{proposition}[the $3$-power column]\label{prop:column}
\PROVED{} For a prime $r$ with $f(r)=3^k$, the values $m\in\{0,2,3\}$ are
impossible: $m_r=0$ would require $r\mid\ord(\kappa_{3^k})=3^{k+1}$, a
power of $3$, while $r\equiv1\pmod{3^k}$ is a prime distinct from $3$;
and $m=2,3$ share $m=1$'s order class by
Proposition~\ref{prop:translates}. Any failure of $C_k$ therefore gives
$m_r\geq4$ directly.
\end{proposition}

\begin{proposition}[norm tower]\label{prop:tower}
\PROVED{} $\gamma_{k-1}=\gamma_k^3+\gamma_k
=\Norm_{L_k/L_{k-1}}(\gamma_k)$ (minimal polynomial
$X^3+X+\gamma_{k-1}$), so $\ord(\gamma_{k-1})\mid\ord(\gamma_k)$. With
lifting-the-exponent ($v_r(2^{3^k}-1)=v_r(2^{f(r)}-1)$ for old primes
$r\neq3$), full-order parts propagate upward, and conjecture $C_k$
reduces to $C_{k-1}$ plus primitivity at the new primes
$r\mid\Phi_{3^k}(2)$.
\end{proposition}

\begin{proposition}[the norm-kernel seam]\label{prop:norm-kernel-seam}
\PROVED{} Put $q=2^{3^{k-1}}$, so
$L_k=\F_{q^3}$ and $L_{k-1}=\F_q$.  There is a canonical direct product
\[
  L_k^{\times}=L_{k-1}^{\times}\times U_k,
  \qquad |U_k|=q^2+q+1=\Phi_{3^k}(2),
\]
where $U_k=\ker\Norm_{L_k/L_{k-1}}$.  If
$\gamma_k=a_kb_k$ in this product, then
\[
  \gamma_{k-1}=\Norm(\gamma_k)=a_k^3,
\]
and cubing is an automorphism of $L_{k-1}^{\times}$.  Consequently
\[
  C_k\iff C_{k-1}\ \text{and}\ b_k\ \text{generates }U_k.
\]
In particular, the norm recursion determines the old component $a_k$
exactly but supplies no order information about the new component $b_k$.
\end{proposition}

\begin{proof}
The two displayed factor orders are coprime:
\[
 \gcd(q-1,q^2+q+1)=\gcd(q-1,3)=1,
\]
because $q=2^{\text{odd}}\equiv2\pmod3$.  The cyclic group
$L_k^{\times}$ therefore splits as asserted.  The relative norm is
$z\mapsto z^{1+q+q^2}$; it kills $U_k$ and restricts on
$L_{k-1}^{\times}$ to $a\mapsto a^3$, which is bijective because
$3\nmid q-1$.  The final equivalence follows from coprimality of the two
factor orders.
\end{proof}

\begin{proposition}[primitive Dickson iterate]\label{prop:dickson-iterate}
\PROVED{} Let $D_n(X,1)$ be the Dickson polynomial characterized by
$D_n(u+u^{-1},1)=u^n+u^{-n}$, and put $f(X)=X^3+X$.  Then
\[
 P_k(X):=D_{3^k}(X,1)+1=f^{\circ k}(X)+1
\]
is the minimal polynomial of $\gamma_k$ over $\F_2$.  In particular it
is irreducible of degree $3^k$, and
\[
 C_k\iff P_k(X)\text{ is a primitive polynomial over }\F_2.
\]
The polynomials also obey the coefficient recursion
\[
 P_{k+1}=P_k^3+P_k^2+1,
 \qquad P_1=X^3+X+1.
\]
\end{proposition}

\begin{proof}
The Dickson composition law gives
$D_{3^k}=D_3^{\circ k}=f^{\circ k}$ in characteristic two.  Since
$\zeta^{3^k}$ is a primitive cube root of unity,
\[
 D_{3^k}(\gamma_k,1)
 =\zeta^{3^k}+\zeta^{-3^k}=1.
\]
On the other hand, the Frobenius orbit of $\zeta+\zeta^{-1}$ has length
$3^k$: the least $j>0$ with $2^j\equiv\pm1\pmod{3^{k+1}}$ is $3^k$.
Thus the displayed monic polynomial, which has degree $3^k$, is the
minimal polynomial.  Its root is primitive precisely when the
polynomial is primitive.  Finally, writing $D_{3^k}=P_k+1$ and using
$D_3(Y)=Y^3+Y$ gives
$P_{k+1}=(P_k+1)^3+(P_k+1)+1=P_k^3+P_k^2+1$.
\end{proof}

\begin{proposition}[shifted cyclotomic-unit form]\label{prop:shifted-unit}
\PROVED{} Put $x_k=1+\zeta\in\F_{2^{2h}}^\times$, where $h=3^k$ and
$N=3^{k+1}$.  Then
\[
 x_k^{2^h-1}=\zeta^{-1},\qquad
 x_k^{2^h+1}=\gamma_k,\qquad
 \ord(x_k)=N\,\ord(\gamma_k).
\]
Consequently
\[
 C_k\iff \ord(1+\zeta)=3^{k+1}(2^{3^k}-1).
\]
Thus the cubic arm is equivalently maximal order for the reduction at
$2$ of the cyclotomic unit $1+\zeta_{3^{k+1}}$.
\end{proposition}

\begin{proof}
Because $2^h\equiv-1\pmod N$,
$x_k^{2^h}=1+\zeta^{-1}=\zeta^{-1}x_k$.  Dividing and multiplying by
$x_k$ gives the first two identities.  The factors $2^h-1$ and
$2^h+1$ are coprime.  In the corresponding direct-product
decomposition of $\F_{2^{2h}}^\times$, the first displayed power has
exact order $N$ and the second has order $\ord(\gamma_k)$, so the order
of $x_k$ is their product.
\end{proof}

\begin{proposition}[norm-coherent cyclotomic lift]\label{prop:unit-lift}
\PROVED{} Let $\xi_k=\zeta_{3^{k+1}}$ in characteristic zero and
$u_k=1+\xi_k$.  Then $2$ is inert in $\Q(\xi_k)$ and
\[
 \Norm_{\Q(\xi_k)/\Q(\xi_{k-1})}(u_k)=u_{k-1}.
\]
The norm to the maximal real subfield is
\[
 \Norm_{\Q(\xi_k)/\Q(\xi_k+\xi_k^{-1})}(u_k)
 =2+\xi_k+\xi_k^{-1},
\]
whose reduction modulo the unique prime above $2$ is $\gamma_k$.
Thus Proposition~\ref{prop:shifted-unit} is the residue-field shadow of
a norm-coherent cyclotomic-unit system in the cyclotomic
$\mathbb Z_3$-tower.
Conditional on $C_{k-1}$, $C_k$ is exactly the assertion that, for every
prime $\ell\mid\Phi_{3^k}(2)$, the reduction of $u_k$ is not an
$\ell$-th power.
\end{proposition}

\begin{proof}
The residue degree of $2$ in $\Q(\xi_k)$ is
$\ord_{3^{k+1}}(2)=2\cdot3^k=\varphi(3^{k+1})$, proving inertness.  The
three relative conjugates of $\xi_k$ are $\xi_k,\xi_k\omega$, and
$\xi_k\omega^2$, so
\[
 \prod_{j=0}^2(1+\xi_k\omega^j)=1+\xi_k^3=1+\xi_{k-1}.
\]
Multiplication by the complex-conjugate factor gives the real norm.
The final assertion follows from Proposition~\ref{prop:shifted-unit}:
in a cyclic group an element has the full $\ell$-primary part of the
group order precisely when it is not an $\ell$-th power.
\end{proof}

\begin{proposition}[circular-unit and ray-class form]
\label{prop:circular-ray}
\PROVED{} Let
\[
 F_k=\Q(\xi_k+\xi_k^{-1}),\qquad
 \lambda_k=(1-\xi_k)(1-\xi_k^{-1}),
\]
let $\mathfrak p_k=(2)$ be the unique prime of $F_k$ above $2$, and
let $\mathcal C_k\subset\mathcal O_{F_k}^{\times}$ be the real
circular-unit group.  If $\sigma(\xi_k)=\xi_k^2$, then
\[
 c_k:=\frac{\sigma(\lambda_k)}{\lambda_k}
     =(1+\xi_k)(1+\xi_k^{-1})
     =2+\xi_k+\xi_k^{-1}
\]
is a circular unit, its Galois conjugates generate a subgroup of
$2$-power index in $\mathcal C_k$, and its reduction modulo
$\mathfrak p_k$ is $\gamma_k$.  Since the residue-field unit group has
odd order, that $2$-power index disappears on reduction.  Consequently
\[
 \overline{\mathcal C_k}=\langle\gamma_k\rangle
 \subset (\mathcal O_{F_k}/\mathfrak p_k)^{\times}
       \simeq\F_{2^{3^k}}^{\times},                         \tag{\(\clubsuit\)}
\]
and hence
\[
 [\F_{2^{3^k}}^{\times}:\langle\gamma_k\rangle]
 =[(\mathcal O_{F_k}/\mathfrak p_k)^{\times}:
     \overline{\mathcal C_k}].                              \tag{\(\spadesuit\)}
\]
In particular, $C_k$ is exactly surjectivity modulo $2$ of the real
circular units of conductor $3^{k+1}$.

More precisely, put $E_k=\mathcal O_{F_k}^{\times}$ and let
$\operatorname{Cl}_{\mathfrak p_k}(F_k)$ be the ray class group of finite modulus
$\mathfrak p_k$.  The ray class exact sequence gives
\[
 [\F_{2^{3^k}}^{\times}:\langle\gamma_k\rangle]
 =\frac{|\operatorname{Cl}_{\mathfrak p_k}(F_k)|}
        {|\operatorname{Cl}(F_k)|}
   [\overline{E_k}:\overline{\mathcal C_k}],                \tag{\(\heartsuit\)}
\]
where
$[\overline{E_k}:\overline{\mathcal C_k}]\mid
 [E_k:\mathcal C_k]$.  The circular-unit index theorem identifies
the odd part of the latter index with the odd part of
$|\operatorname{Cl}(F_k)|$.
Thus, for an odd prime $\ell\nmid|\operatorname{Cl}(F_k)|$,
\[
 v_{\ell}\!\left(
 [\F_{2^{3^k}}^{\times}:\langle\gamma_k\rangle]\right)
 =v_{\ell}\!\left(
 \frac{|\operatorname{Cl}_{\mathfrak p_k}(F_k)|}
      {|\operatorname{Cl}(F_k)|}\right).                    \tag{\(\diamondsuit\)}
\]
Equation~\((\heartsuit)\) names the obstruction but does not reduce its
difficulty: the ray class group is defined through this residue-unit
quotient, and the class-number exclusion in \((\diamondsuit)\) is itself
not known uniformly in $k$.
Finally, with $h=3^k$, $q=2^{h/3}$, and
$\tau=\sigma^{h/3}$, the relative circular unit
\[
 \varepsilon_k=\frac{\tau(c_k)}{c_k},\qquad
 \Norm_{F_k/F_{k-1}}(\varepsilon_k)=1,
\]
reduces to $\eta_k$.  Thus the exact character obstruction of
Proposition~\ref{prop:character-test} is the new relative part of the
same ray-class obstruction.
\end{proposition}

\begin{proof}
The class of $2$ generates
$(\mathbb Z/3^{k+1}\mathbb Z)^{\times}/\{\pm1\}$, so $\sigma$ generates
$\operatorname{Gal}(F_k/\Q)$.  The squares of the standard real
circular units are generated, up to $-1$, by the quotients
$\sigma^j(\lambda_k)/\lambda_k$.  Telescoping gives
\[
 \frac{\sigma^j(\lambda_k)}{\lambda_k}
 =\prod_{i=0}^{j-1}\sigma^i(c_k),
\]
so the conjugates of $c_k$ generate this square subgroup.  Squaring is
an automorphism of the odd-order residue-field unit group, so the full
circular-unit group and this subgroup have the same reduction image.
Reduction commutes with $\sigma$, which becomes Frobenius at the inert
prime $2$, and
\[
 \overline{c_k}
 =(1+\zeta)(1+\zeta^{-1})=\zeta+\zeta^{-1}=\gamma_k.
\]
The reductions of the conjugates are therefore
$\gamma_k,\gamma_k^2,\gamma_k^{2^2},\ldots$, and they generate the
single multiplicative subgroup $\langle\gamma_k\rangle$.  This proves
\((\clubsuit)\) and \((\spadesuit)\).

The exact sequence
\[
 E_k\longrightarrow
 (\mathcal O_{F_k}/\mathfrak p_k)^{\times}\longrightarrow
 \operatorname{Cl}_{\mathfrak p_k}(F_k)
 \longrightarrow\operatorname{Cl}(F_k)
 \longrightarrow1
\]
proves \((\heartsuit)\).  The divisibility of reduction indices is
immediate from $\mathcal C_k\subset E_k$, and the circular-unit index
theorem gives \((\diamondsuit)\); see Washington~\cite[Ch.~8]{Washington97}.
Finally $\tau$ generates $\operatorname{Gal}(F_k/F_{k-1})$, so the norm of its
coboundary $\tau(c_k)/c_k$ telescopes to $1$.  Modulo $2$ it is
$\gamma_k^q/\gamma_k=\eta_k$.
\end{proof}

\begin{proposition}[exceptional eigenspace and Hilbert-reciprocity no-go]
\label{prop:cubic-reciprocity-no-go}
\PROVED{} Put \(n=3^k\), \(m=3n\), let
\[
 F=\mathbb Q(\xi_m+\xi_m^{-1}),\qquad
 \mathfrak p=(2),
\]
and write \(\Delta=\operatorname{Gal}(F/\mathbb Q)=\langle\sigma\rangle\),
where \(\sigma(\xi_m)=\xi_m^2\).  Let
\(\ell\mid\Phi_n(2)\), and set
\[
 V_\ell=(\mathcal O_F/\mathfrak p)^\times/
          ((\mathcal O_F/\mathfrak p)^\times)^\ell .
\]
Then \(V_\ell\) is one-dimensional over \(\mathbb F_\ell\), and
\(\sigma\) acts on it as the scalar \(2\).  Let
\[
 \operatorname{ev}_2:\mathbb F_\ell[\Delta]\longrightarrow\mathbb F_\ell,
 \qquad \sum a_j\sigma^j\longmapsto\sum a_j2^j.
\]
The following reciprocity operators vanish on \(V_\ell\):
\[
 \operatorname{ev}_2(1+\sigma+\cdots+\sigma^{n-1})=0,
 \qquad \operatorname{ev}_2(2-\sigma)=0.                 \tag{\(\mathsf{ER}_1\)}
\]
If \(q=2^{n/3}\) and \(\tau=\sigma^{n/3}\), then
\[
 \operatorname{ev}_2(1+\tau+\tau^2)=1+q+q^2=0,
 \qquad \operatorname{ev}_2(\tau-1)=q-1\ne0.             \tag{\(\mathsf{ER}_2\)}
\]
Moreover the plus projection of the classical Stickelberger element of
\(\mathbb Q(\xi_m)\) acts as zero on \(V_\ell\).  Its zero is simple:
for
\[
 \nu_n(X)=1+X+\cdots+X^{n-1}
\]
one has
\[
 \nu_n(2)=0,\qquad
 \nu_n'(2)=n2^{n-1}\ne0\pmod\ell.                       \tag{\(\mathsf{ER}_3\)}
\]

Let
\[
 c=(1+\xi_m)(1+\xi_m^{-1}),\qquad
 \gamma=\overline c\in\mathbb F_{2^n}^\times,\qquad
 \mathcal A_{k,\ell}=\gamma^{(2^n-1)/\ell}\in\mu_\ell.
                                                               \tag{\(\mathsf{ER}_4\)}
\]
Then cubic failure at \(\ell\) is exactly
\(\mathcal A_{k,\ell}=1\).  Although
\(\mathcal A_{k,\ell}\) is a local Kummer/Artin symbol, global Hilbert
reciprocity imposes no condition on this one distinguished value.
More precisely, put \(E=F(\mu_\ell)\).  The prime \(\mathfrak p\)
splits into \(\ell-1\) primes in \(E/F\).  If
\(\mathfrak P_0\mid\mathfrak p\) is chosen so that
\[
 (c,2)_{\mathfrak P_0,\ell}=\mathcal A_{k,\ell}
\]
for the reciprocity convention sending the uniformizer \(2\) to
arithmetic Frobenius, then the \(\ell-1\) symbols above two are
\[
 (c,2)_{a\mathfrak P_0,\ell}=\mathcal A_{k,\ell}^{\,a},
 \qquad a\in\mathbb F_\ell^\times.                       \tag{\(\mathsf{ER}_5\)}
\]
Every local symbol away from two is one, and Hilbert reciprocity
reduces to
\[
 \prod_{a\in\mathbb F_\ell^\times}\mathcal A_{k,\ell}^{\,a}
 =\mathcal A_{k,\ell}^{\ell(\ell-1)/2}=1,                \tag{\(\mathsf{ER}_6\)}
\]
which holds for every possible value of
\(\mathcal A_{k,\ell}\).  Thus norm coherence, the classical
Stickelberger ideal, and global Kummer reciprocity cannot distinguish
failure from nonfailure.  A completion must compute a first-order unit
regulator in the simple exceptional eigenspace, rather than another
ideal factorization or product formula.
\end{proposition}

\begin{proof}
The standard cyclotomic-divisor criterion gives
\(\operatorname{ord}_\ell(2)=n\).  Since
\((\mathcal O_F/\mathfrak p)^\times\simeq\mathbb F_{2^n}^\times\),
its quotient by \(\ell\)-th powers has order \(\ell\), and Frobenius
acts by \([x]\mapsto[x^2]=2[x]\).  This proves the module assertion.
Equations \((\mathsf{ER}_1)\) follow from
\(2^n=1\pmod\ell\).  The element \(q=2^{n/3}\) has exact order three
modulo \(\ell\), proving \((\mathsf{ER}_2)\).

Let \(G=\operatorname{Gal}(\mathbb Q(\xi_m)/\mathbb Q)\), let \(j\)
be complex conjugation, and write
\[
 \theta_m=\frac1m
 \sum_{\substack{1\le a<m\\(a,m)=1}}a\,\sigma_a^{-1}
\]
for the Stickelberger element.  Pairing \(a\) with \(m-a\) gives
\[
 (1+j)\theta_m=\sum_{g\in G}g.
\]
On an even representation, the right side is twice the norm element
of \(\Delta\), whose evaluation at two is zero by
\((\mathsf{ER}_1)\).  Since \(\ell\ne2,3\), division by two and by
\(m\) is valid modulo \(\ell\).  Hence \(\theta_m\), and therefore
every classical Stickelberger-ideal multiple controlling
Gauss--Jacobi-sum ideals, acts trivially on \(V_\ell\).  The
particularly relevant integrality factor \(2-\sigma\) vanishes
separately.  Differentiating
\[
 (X-1)\nu_n(X)=X^n-1
\]
at \(X=2\) gives \((\mathsf{ER}_3)\), since \(\ell\nmid2n\).

The reduction of \(c\) is \(\gamma\).  In the local field
\(F_{\mathfrak p}\), which contains \(\mu_\ell\), adjoining an
\(\ell\)-th root \(z\) of \(c\) is the unramified Kummer extension
selected by the residue of \(c\).  Arithmetic Frobenius satisfies
\[
 \frac{\operatorname{Frob}(z)}z
 =\overline z^{\,2^n-1}
 =\gamma^{(2^n-1)/\ell},
\]
proving \((\mathsf{ER}_4)\) and the failure criterion.  Also, for the
relative circular unit \(\varepsilon=\tau(c)/c\),
\[
 (\varepsilon,2)_{\mathfrak p,\ell}
 =\mathcal A_{k,\ell}^{q-1};
\]
because \(q-1\ne0\pmod\ell\), the relative coboundary carries exactly
the same unknown symbol and supplies no new relation.

It remains to audit global reciprocity.  The completion
\(F_{\mathfrak p}\) is the unramified degree-\(n\) extension of
\(\mathbb Q_2\).  Since
\(\mathbb Q_2(\mu_\ell)/\mathbb Q_2\) is also unramified of degree
\(\operatorname{ord}_\ell(2)=n\), it is already contained in
\(F_{\mathfrak p}\).  The conductors \(m\) and \(\ell\) are coprime,
so \(F\cap\mathbb Q(\mu_\ell)=\mathbb Q\); consequently
\([E:F]=\ell-1\), and \(\mathfrak p\) splits completely in \(E/F\).
Galois functoriality of the Hilbert symbol gives
\((\mathsf{ER}_5)\).

We now check every place above \(\ell\), the only point at which wild
local reciprocity might have supplied extra information.  Let
\(K_0\) be the maximal unramified subfield of an \(\ell\)-adic
completion of \(E\).  That completion is
\(K=K_0(\mu_\ell)\), with absolute ramification index
\(e=\ell-1\), and both \(c\) and \(2\) belong to \(K_0^\times\).
For any unit \(a\in K_0^\times\), write
\[
 a=[\overline a]\,\langle a\rangle,\qquad
 \langle a\rangle\in1+\ell\mathcal O_{K_0}.
\]
The Teichm\"uller factor has order prime to \(\ell\), hence is an
\(\ell\)-th power.  If \(U_i=1+\mathfrak m_K^i\), then
\[
 1+\ell\mathcal O_{K_0}\subset U_{\ell-1}.
\]
The standard Hilbert-filtration orthogonality for a field containing
\(\mu_\ell\) is~\cite{Iwasawa68}
\[
 (U_i,U_j)_\ell=1\qquad\text{when}\qquad
 i+j>\frac{\ell e}{\ell-1}=\ell.                       \tag{\(\mathsf{ER}_7\)}
\]
Taking \(i=j=\ell-1\) proves
\((c,2)_{\lambda,\ell}=1\) at every \(\lambda\mid\ell\).
At every finite place away from \(2\ell\), both entries are units and
the tame Hilbert-symbol formula is one; archimedean places contribute
nothing because \(\ell\) is odd.  The global product formula is
therefore precisely \((\mathsf{ER}_6)\).  Since
\(\sum_{a=1}^{\ell-1}a\equiv0\pmod\ell\), this relation is
tautological and leaves the single symbol at \(\mathfrak P_0\)
arbitrary.
\end{proof}

\begin{proposition}[semidirect ray extension forced by cubic failure]
\label{prop:cubic-ray-extension}
\PROVED{} Retain Proposition~\ref{prop:circular-ray}, put
\(n=3^k\), and let \(\ell\mid\Phi_{3^k}(2)\).  Assume
\(\ell\nmid|\operatorname{Cl}(F_k)|\).  Then
\[
 \gamma_k\in(\F_{2^n}^\times)^\ell                  \tag{\(\mathsf{RF}\)}
\]
if and only if there exists a cyclic extension \(M/F_k\) of degree
\(\ell\), conductor exactly \(\mathfrak p_k=(2)\), totally tamely
ramified there and unramified at every other finite prime.  This
extension is Galois over \(\Q\), with
\[
 \operatorname{Gal}(M/\Q)\cong
 C_\ell\rtimes_2 C_n,
 \qquad \sigma\tau\sigma^{-1}=\tau^2.               \tag{\(\mathsf{RF}_1\)}
\]
Moreover, the unique prime of \(F_k\) above three splits completely in
\(M/F_k\).  Both prescribed local behaviors are admissible: at two this
is the standard tame extension with inertia \(C_\ell\), residue degree
\(n\), and Frobenius action \(\tau\mapsto\tau^2\); at three the
\(C_\ell\)-layer is split.  Thus secant saturation and circular-unit
reciprocity give no local contradiction.  Without the class-number
hypothesis, cubic failure yields the sharp alternative
\[
 \ell\mid|\operatorname{Cl}(F_k)|
 \quad\hbox{or the extension in \((\mathsf{RF}_1)\) exists}.  \tag{\(\mathsf{RF}_2\)}
\]
\end{proposition}

\begin{proof}
Write \(E_k=\mathcal O_{F_k}^\times\) and
\(\mathcal C_k\) for the circular units.  Proposition~\ref{prop:circular-ray}
identifies the reduction of \(\mathcal C_k\) with
\(\langle\gamma_k\rangle\).  If \((\mathsf{RF})\) holds, then
\(\mathcal C_k\) maps trivially to
\[
 V:=\F_{2^n}^\times/(\F_{2^n}^\times)^\ell\cong C_\ell.
\]
The odd part of \([E_k:\mathcal C_k]\) is the odd part of the class
number.  Under the stated hypothesis its \(\ell\)-part vanishes, so
\(E_k\) also maps trivially to \(V\).  The ray sequence then supplies
an order-\(\ell\) quotient of
\(\operatorname{Cl}_{\mathfrak p_k}(F_k)\).  Because the ordinary class
group has no \(\ell\)-part, the quotient lies entirely in tame inertia
at \(\mathfrak p_k\).  Class field theory gives the asserted
\(M/F_k\).

The quotient is Galois-stable.  The generator \(\sigma\) of
\(\operatorname{Gal}(F_k/\Q)\) acts on the residue field as
\(x\mapsto x^2\), hence on \(V\) by exponent two.  Since
\(\ord_\ell(2)=n\), this action is faithful; Schur--Zassenhaus, or the
coprimality of \(n\) and \(\ell\), gives \((\mathsf{RF}_1)\).
Conversely, tame inertia in any such extension gives the unique quotient
\(\F_{2^n}^\times\twoheadrightarrow C_\ell\).  Global units lie in its
kernel, so in particular \(\gamma_k\) is an \(\ell\)-th power.

Finally let
\[
 \lambda_k=(1-\xi_k)(1-\xi_k^{-1}),\qquad
 c_k=(1+\xi_k)(1+\xi_k^{-1}).
\]
The element \(\lambda_k\) generates the unique prime
\(\mathfrak r_k\) above three, and
\(\lambda_k/c_k\equiv1\pmod{\mathfrak p_k}\), because both reduce to
\((1+\zeta)(1+\zeta^{-1})=\gamma_k\).  Hence
\(\mathfrak r_k=(\lambda_k/c_k)\) is trivial in the full ray class
group and splits in \(M/F_k\).  The local statements now follow from
\(\ell\mid2^n-1\) and the displayed semidirect action.  If the
class-number hypothesis fails, the first alternative in
\((\mathsf{RF}_2)\) holds; otherwise the preceding construction proves
the second.
\end{proof}

\begin{proposition}[explicit cubic kernel tower]\label{prop:eta-tower}
\PROVED{} Define $\eta_0=0$.  For $k\geq1$ put
\[
 q_k=2^{3^{k-1}},\qquad
 \eta_k=\gamma_k^{q_k-1},\qquad
 U_k=\ker\!\left(\Norm_{\F_{q_k^3}/\F_{q_k}}\right).
\]
Then $\eta_k\in U_k$, $|U_k|=q_k^2+q_k+1=\Phi_{3^k}(2)$, and, with
$\omega=\zeta^{3^k}\in\F_4$,
\[
 q_k\equiv3^k-1\pmod{3^{k+1}},\qquad
 \eta_k=\frac{\omega+\omega^2\zeta^2}{1+\zeta^2},\qquad
 \zeta^2=\frac{\eta_k+\omega}{\eta_k+\omega^2}.
\]
Moreover $\eta_k$ has the irreducible cubic equation over
$\F_{q_k}=\F_{2^{3^{k-1}}}$
\[
 X^3+\eta_{k-1}X^2+(\eta_{k-1}+1)X+1=0.       \tag{\(*\)}
\]
Equivalently,
\[
 \Norm(\eta_k)=1,\qquad
 \operatorname{Tr}(\eta_k)=\eta_{k-1},\qquad
 e_2(\eta_k)=\eta_{k-1}+1.
\]
It also satisfies
\[
 \eta_k^{q_k+1}+\eta_k+1=0,
 \qquad \eta_k+1=\eta_k^{q_k+1},
 \qquad \ord(\eta_k+1)=\ord(\eta_k).          \tag{\(\dagger\)}
\]
Thus $\eta_k$ lies in the planar Singer difference set
\[
 \mathcal H_k=\{y\in U_k:y^{q_k+1}+y+1=0\}
\]
of size $q_k+1$ inside $U_k$~\cite{MRS20}.
Finally,
\[
 C_k\iff C_{k-1}\ \text{and}\
 \ord(\eta_k)=q_k^2+q_k+1.                    \tag{\(**\)}
\]
(Here $C_0$ is the tautological assertion that $1$ generates
$\F_2^\times$.)
Thus the new-prime obstruction is the single assertion that the
distinguished root of \((*)\) is a generator of the cubic norm-one
torus (equivalently, multiplication by it is a Singer cycle).
\end{proposition}

\begin{proof}
The congruence starts with $2\equiv3-1\pmod9$ and lifts by cubing:
if $a\equiv-1+3^k\pmod{3^{k+1}}$, then
$a^3\equiv-1+3^{k+1}\pmod{3^{k+2}}$.  Hence
$\zeta^{q_k}=\omega\zeta^{-1}$.  Dividing
$\gamma_k^{q_k}$ by $\gamma_k$ now gives the stated M\"obius formula;
solving it gives the inverse formula.

Put $t=\zeta^2$ and
$M(T)=(\omega+\omega^2T)/(1+T)$.  The same formula one level down says
$\eta_{k-1}=M(t^3)$ (also for $k=1$, when it gives $\eta_0=0$).  If
$a=\eta_{k-1}$ and $X=M(t)$, then
\[
 \left(\frac{X+\omega}{X+\omega^2}\right)^3
 =\frac{a+\omega}{a+\omega^2}.
\]
Cross-multiplication and $\omega^2+\omega+1=0$ reduce exactly to
\((*)\).  Also
$\eta_k^{1+q_k+q_k^2}=\gamma_k^{q_k^3-1}=1$, so $\eta_k\in U_k$.
It cannot lie in $\F_{q_k}$: otherwise its norm-one condition would give
$\eta_k^3=1$, hence $\eta_k=1$ because $3\nmid q_k-1$; the inverse
M\"obius formula would then give $t=\omega$, contradicting
$\ord(t)=3^{k+1}>3$.  Thus \((*)\) is irreducible, and its coefficients
give the three symmetric identities.

The earlier minimal polynomial
$X^3+X+\gamma_{k-1}$ gives
$\operatorname{Tr}_{\F_{q_k^3}/\F_{q_k}}(\gamma_k)=0$.  Dividing
$\gamma_k+\gamma_k^{q_k}+\gamma_k^{q_k^2}=0$ by $\gamma_k$ yields
$1+\eta_k+\eta_k^{q_k+1}=0$, proving the first two parts of
\((\dagger)\).  Since
$\gcd(q_k+1,q_k^2+q_k+1)=1$, exponentiation by $q_k+1$ is an
automorphism of $U_k$, proving the equality of orders.  The cited
difference-set identification is the norm-one realization of the planar
Singer set.

In the direct product of Proposition~\ref{prop:norm-kernel-seam},
raising to $q_k-1$ kills the old-field component and is an automorphism
on $U_k$, since $\gcd(q_k-1,q_k^2+q_k+1)=1$.  Therefore the kernel
component of $\gamma_k$ generates $U_k$ exactly when $\eta_k$ does;
Proposition~\ref{prop:norm-kernel-seam} gives \((**)\).
\end{proof}

\begin{proposition}[Singer-orbit sieve]\label{prop:singer-sieve}
\PROVED{} Let $k\geq1$.  Put $n=3^k$, $q=2^{n/3}$, and
$N=q^2+q+1=|U_k|$.  If a prime $\ell\mid N$ satisfies
\[
  \ell>\frac{N}{n(n-1)+1},
\]
then $\eta_k\notin U_k^\ell$.  Consequently, if
$\ell^a\Vert N$, then
\[
  v_\ell(\ord\eta_k)=a.
\]
Unconditionally,
\[
  \ord(\eta_k)\geq n(n-1)+1,
  \qquad
  [U_k:\langle\eta_k\rangle]
  \leq\frac{N}{n(n-1)+1}.
\]
\end{proposition}

\begin{proof}
Let
\[
  D=\{y\in U_k:y^{q+1}+y+1=0\}.
\]
This is the cyclic $(N,q+1,1)$ Singer difference set: every
nonidentity element of $U_k$ has a unique expression $xy^{-1}$ with
distinct $x,y\in D$.  Since the minimal polynomial of $\eta_k$ over
$\F_2$ has degree $n$, its Frobenius orbit
\[
  A=\{\eta_k^{2^i}:0\leq i<n\}
\]
has $n$ elements.  The equation defining $D$ is Frobenius-stable, so
$A\subset D$.

Suppose $\eta_k\in H:=U_k^\ell$.  The subgroup $H$ is characteristic
in the cyclic group $U_k$, hence Frobenius-stable, and $A\subset H$.
The $n(n-1)$ quotients $xy^{-1}$ for ordered distinct pairs
$(x,y)\in A^2$ are nonidentity elements of $H$ and are all distinct by
the Singer difference-set property.  Therefore
\[
  |H|-1\geq n(n-1).
\]
Since $|H|=N/\ell$, this contradicts the displayed strict inequality.
In a cyclic group of order $\ell^a m$, $\ell\nmid m$, lying outside
the subgroup of $\ell$-th powers is equivalent to having the full
$\ell^a$ part in one's order, proving the first assertion.

Finally the same distinct quotients all lie in
$\langle\eta_k\rangle$, whether or not an $\ell$ is fixed.  Thus
$\ord(\eta_k)-1\geq n(n-1)$, which gives the last two inequalities.
\end{proof}

\begin{remark}[why the translated orbit does not double the sieve]
The identity $y+1=y^{q+1}$ puts the translated Frobenius orbit
$A+1=A^{q+1}$ in the same subgroup as $A$.  It supplies no second
quotient family: modulo $N$ one has $(q+1)^2\equiv q$ and
$q(q+1)\equiv-1$, while the quotient set of $A$ is stable under both
$q$-th powers and inversion.  Hence the internal quotient sets of
$A$ and $A+1$ coincide.  Cross quotients already collide, for example
$(y+1)/y=y^q$.  Improving Proposition~\ref{prop:singer-sieve} therefore
requires cross-correlation or character information beyond the
$(N,q+1,1)$ difference-set axiom.
\end{remark}

\begin{proposition}[partition-order bound]\label{prop:cubic-partitions}
\PROVED{} Retain Proposition~\ref{prop:singer-sieve}.  Let
\(B(n)\) be the number of partitions into distinct positive parts not
divisible by \(3\) whose total is at most \(n=3^k\).  Then
\[
  \ord(\eta_k)\geq B(n).
\]
In particular, if \(p^{\rm dist}_{\not\,3}(n)\) counts such partitions
of \(n\) itself, then
\[
 \ord(\eta_k)\geq p^{\rm dist}_{\not\,3}(n)
 =\exp\!\left(\left(\frac{\sqrt2\,\pi}{3}+o(1)\right)\sqrt n\right).
                                                               \tag{\(\mathsf B\)}
\]
If a prime \(\ell\mid N\) satisfies
\[
  \ell>\frac{N}{B(n)},
\]
then \(\eta_k\notin U_k^\ell\), so the order of \(\eta_k\) contains
the full \(\ell\)-primary part of \(N\).
\end{proposition}

\begin{proof}
Put \(r=3n\), \(t=\zeta^2\), and retain
\(\omega=\zeta^n\), so \(t^n=\omega^2\) and
\[
 \eta_k=M(t),\qquad
 M(X)=\frac{\omega+\omega^2X}{1+X}.
\]
The element \(2\) generates
\((\mathbb Z/r\mathbb Z)^\times\), of order \(2n\).  For every
positive \(j\) not divisible by \(3\), choose \(z_j\) with
\(2^{z_j}\equiv j\pmod r\).  Frobenius gives
\[
 \eta_k^{2^{z_j}}
 =\omega^j\frac{1+\omega^jt^j}{1+t^j}.                 \tag{\(6\)}
\]
For a finite set \(I\) of such \(j\), put
\[
 s_I=\sum_{j\in I}j,\quad E_I=\sum_{j\in I}2^{z_j},\quad
 A_I(X)=\prod_{j\in I}(1+\omega^jX^j),\quad
 B_I(X)=\prod_{j\in I}(1+X^j).
\]
Thus
\[
 \eta_k^{E_I}=\omega^{s_I}\frac{A_I(t)}{B_I(t)}.       \tag{\(7\)}
\]
All denominators are nonzero when \(s_I\leq n<r\).

Suppose \(s_I,s_J\leq n\) and
\(\eta_k^{E_I}=\eta_k^{E_J}\).  Cross-multiplication gives
\[
 F(X):=\omega^{s_I}A_I(X)B_J(X)
       +\omega^{s_J}A_J(X)B_I(X),\qquad F(t)=0.         \tag{\(8\)}
\]
Now \(M(t^{-1})=M(t)+1=\eta_k+1\), and the Singer equation
\(\eta_k+1=\eta_k^{q+1}\) shows that the assumed equality also holds
after replacing \(\eta_k\) by \(M(t^{-1})\).  Hence \(F(t^{-1})=0\).
Over \(\F_4\), the two elements \(t,t^{-1}\) have the distinct
degree-\(n\) minimal polynomials
\[
 X^n+\omega^2,\qquad X^n+\omega,
\]
because \(\ord_r(4)=n\).  Therefore
\[
 X^{2n}+X^n+1\mid F(X).                                \tag{\(9\)}
\]
But \(\deg F\leq s_I+s_J\leq2n\).  In the only possible equality
case \(s_I=s_J=n\), the two leading coefficients are both
\(\omega^{2n}\) and cancel in characteristic two.  Thus
\(\deg F<2n\), and (9) forces \(F=0\).

The constant coefficient first gives
\(\omega^{s_I}=\omega^{s_J}\).  Cancel that scalar and all factors
belonging to \(I\cap J\).  If \(j_0\) is the least remaining part,
comparison of the \(X^{j_0}\)-coefficients gives
\(\omega^{j_0}=1\): on one side \(j_0\) occurs in an \(A\)-factor,
with coefficient \(\omega^{j_0}\), and on the other in a \(B\)-factor,
with coefficient \(1\); larger parts cannot contribute.  This is
impossible because \(3\nmid j_0\).  Hence \(I=J\).

The \(B(n)\) powers in (7) are therefore pairwise distinct, proving
the order bound.  Meinardus's partition theorem~\cite{Meinardus54},
applied to the Euler product
\[
 G(x):=\prod_{3\nmid j}(1+x^j)
 =\frac{\prod_{j\geq1}(1+x^j)}{\prod_{j\geq1}(1+x^{3j})},
 \qquad
 \log G(e^{-t})\sim\frac{\pi^2}{18t},
\]
gives \((\mathsf B)\).  Finally, if
\(\eta_k\in U_k^\ell\), all \(B(n)\) displayed powers lie in the
subgroup of order \(N/\ell\).  The strict inequality rules this out;
the cyclic-group argument of Proposition~\ref{prop:singer-sieve}
then gives the full \(\ell\)-primary assertion.
\end{proof}

\begin{remark}[ceiling of the height-budgeted amplification]
The preceding construction is scale-optimal within the following
precise class of degree arguments.  For
\(J=(\mathbb Z/3n\mathbb Z)^\times/\{\pm1\}\), assign an atom indexed
by \(j\in J\) the height
\[
 h(j):=\min\{|b|:b\equiv\pm j\pmod{3n}\}.
\]
Suppose a family of integer exponent vectors is certified pairwise
distinct only by bounding the Laurent span after cross-multiplication
with the triangle estimate
\[
 w(e-e'):=\sum_{j\in J}h(j)|e_j-e'_j|\leq2n
\]
and then invoking the two degree-\(n\) minimal polynomials as above
(including the verified leading-term cancellation at the boundary).
Then the family has cardinality \(\exp(O(\sqrt n))\).

Indeed, fix one vector \(e_0\).  Every difference \(e-e_0\) lies in
the displayed weighted \(\ell^1\)-ball.  Each positive height occurs
only \(O(1)\) times, so the number of such lattice points is bounded
by the coefficients through degree \(2n\) of a fixed power of
\[
 \prod_{m\geq1}
 \bigl(1+2x^m+2x^{2m}+\cdots\bigr)
 \preccurlyeq
 \prod_{m\geq1}(1-x^m)^{-2}.
\]
Colored-partition bounds give \(\exp(O(\sqrt n))\).  Polynomially many
Frobenius or cyclic shifts do not change that scale, whereas
\(|U_k|=\exp((2\log2/3+o(1))n)\).  Thus more triangle-budgeted subset
encodings cannot prove primitivity.  This does \emph{not} rule out
systematic factor cancellations or exact reduction modulo
\(\Phi_{3n}\); those require new arithmetic input rather than the
generic degree budget.
\end{remark}

\begin{proposition}[no transfer and the exact character test]
\label{prop:character-test}
\PROVED{} The kernel orders are pairwise coprime: for $1\leq i<k$,
\[
 \gcd\!\left(\Phi_{3^i}(2),\Phi_{3^k}(2)\right)=1.
\]
For each prime $\ell\mid |U_k|$, let
$\chi_{k,\ell}:U_k\to\mu_\ell$ be a character of exact order $\ell$.
Then
\[
 \operatorname{Hom}\!\left(
 \F_{2^{3^{k-1}}}^{\times},\mu_\ell\right)=1; \tag{\(\ddagger_0\)}
\]
no multiplicative character or power-residue datum from a previous
level can transfer an $\ell$-part to level $k$.  Moreover,
\[
 C_k\iff C_{k-1}\ \text{and}\
 \chi_{k,\ell}(\eta_k)\neq1
 \quad\text{for every prime }\ell\mid |U_k|.       \tag{\(\ddagger\)}
\]
Equivalently, with
\[
 \rho_{k,\ell}(A)=\chi_{k,\ell}(A^{q_k-1})
 \quad(A\in\F_{q_k^3}^{\times}),
\]
where $\rho_{k,\ell}$ is trivial on $\F_{q_k}^{\times}$, the missing
assertion is $\rho_{k,\ell}(\gamma_k)\neq1$ for every such $\ell$.
This is the smallest power-residue form of the cubic obstruction.
\end{proposition}

\begin{proof}
If a prime $r$ divided two displayed cyclotomic values, then
$r\neq3$ and the standard cyclotomic-divisor criterion would give both
$\ord_r(2)=3^i$ and $\ord_r(2)=3^k$, a contradiction.  Thus every
$\ell\mid |U_k|$ is coprime to
$|\F_{2^{3^{k-1}}}^{\times}|=\prod_{j<k}\Phi_{3^j}(2)$, proving
\((\ddagger_0)\).  In the cyclic group $U_k$, the kernel of an exact
order-$\ell$ character is $U_k^\ell$.  If
$|U_k|=\ell^a m$ with $\ell\nmid m$, then an element has the full
$\ell^a$ in its order exactly when it is not in $U_k^\ell$.  Applying
this for every prime divisor and using \((**)\) proves
\((\ddagger)\).  The formula for $\rho_{k,\ell}$ and
$\eta_k=\gamma_k^{q_k-1}$ give the final equivalence.
\end{proof}

\begin{proposition}[short cyclotomic-period form of the cubic obstruction]
\label{prop:cubic-short-period}
\PROVED{} Let \(k\geq1\), put $n=3^k$, $q=2^{n/3}$, and let
$\ell$ be a prime divisor of $\Phi_{3^k}(2)$.  Let
$\rho=\rho_{k,\ell}$ be the exact
$\ell$-character from Proposition~\ref{prop:character-test}, realized
with values in the complex roots of unity.  In
$\F_{2^{2n}}$ let $\mu_{3n}$ denote the roots of unity of order
dividing $3n=3^{k+1}$, and set
\[
 \mathcal S_{k,\ell}
 :=\sum_{\substack{t\in\mu_{3n}\\t\ne1}}
       \rho(t+t^{-1}).                                    \tag{\(\mathsf S\)}
\]
Every argument of $\rho$ in this sum is a nonzero element of
$\F_{2^n}$.  If $\rho(\gamma_k)=\zeta_\ell^a$, then
\[
 \mathcal S_{k,\ell}
 =(n-1)+2\sum_{i=0}^{n-1}\zeta_\ell^{a2^i}.              \tag{\(\mathsf P\)}
\]
Consequently
\[
 \chi_{k,\ell}(\eta_k)=1
 \iff \mathcal S_{k,\ell}=3n-1.                          \tag{\(\mathsf C\)}
\]
Thus the surviving datum is exactly a short Gauss period over the
order-$n$ subgroup generated by $2$ in
$(\mathbb Z/\ell\mathbb Z)^\times$.
\end{proposition}

\begin{proof}
Since $2^n\equiv-1\pmod{3n}$, every $t+t^{-1}$ in
\((\mathsf S)\) is fixed by $2^n$ and lies in $\F_{2^n}$; it is
nonzero unless $t=1$.  The $n-1$ nonidentity elements of
$\mu_n$ give real cyclotomic elements in the lower field
$\F_{2^{n/3}}^\times$, on which $\rho$ is trivial.

The remaining $2n$ elements are the primitive $3n$-th roots.  Pairing
them by inversion leaves the $n$ Frobenius conjugates
$\gamma_k^{2^i}$, each twice.  This proves \((\mathsf P)\).  Since
$q-1$ is invertible modulo $\ell$,
$\chi_{k,\ell}(\eta_k)=1$ exactly when
$\rho(\gamma_k)=1$, which gives the forward implication in
\((\mathsf C)\).  Conversely every summand in \((\mathsf S)\) has
complex absolute value one.  Equality with the number of summands
$3n-1$ forces equality in the triangle inequality and hence forces
every summand to be $1$, in particular $\rho(\gamma_k)=1$.
\end{proof}

\begin{remark}[why a full Singer character sum is too coarse]
The Singer difference-set character sum controls a set of size $q+1$
and has the usual square-root magnitude.  Equation~\((\mathsf P)\)
instead localizes to only $3n-1$ cyclotomic terms, with the unknown
part supported on the order-$n$ subgroup $\langle2\rangle$ modulo
$\ell$.  Generic Weil or difference-set bounds do not evaluate this
localized scalar; an all-level proof must rule out its extremal value
in \((\mathsf C)\).
\end{remark}

\begin{proposition}[cyclotomic-secant saturation]
\label{prop:cubic-secant-saturation}
\PROVED{} Retain $n=3^k$, put $m=3n$, and let
$K=\F_{2^{2n}}$ contain a primitive $m$-th root \(\zeta\).  For
every prime \(\ell\mid\Phi_{3^k}(2)\), the following are equivalent:

\begin{enumerate}[label=\textup{(\roman*)}]
\item \(\eta_k\in U_k^\ell\);
\item \(\gamma_k\in(\F_{2^n}^\times)^\ell\);
\item \(1+\zeta\in(K^\times)^\ell\);
\item for every \(a\not\equiv b\pmod m\),
      \(\zeta^a+\zeta^b\in(K^\times)^\ell\).
\end{enumerate}

Thus, for an exact order-\(\ell\) character \(\chi\) of \(K^\times\),
failure is also equivalent to the extremal identity

\[
 \sum_{\substack{t\in\mu_m\\t\ne1}}\chi(1-t)=m-1.       \tag{\(\mathsf{CS}\)}
\]

It saturates every cyclotomic secant, not only the Frobenius orbit of
the distinguished Gaussian period.
\end{proposition}

\begin{proof}
The first equivalence follows from

\[
 \eta_k=\gamma_k^{2^{n/3}-1},
\]

because \(\ell\nmid2^{n/3}-1\).  Next,

\[
 \gamma_k=\zeta^{-1}(1+\zeta)^2.
\]

The order of \(\zeta\) is coprime to \(\ell\), and squaring is
invertible on the \(\ell\)-primary quotient.  Moreover the inclusion
\(\F_{2^n}^\times\subset K^\times\) preserves that quotient, since
\(\ell\nmid2^n+1\).  This proves the equivalence with (iii).

Write \(d=b-a\).  The factor \(\zeta^a\) is automatically an
\(\ell\)-th power.  If \(3\nmid d\), then (2) generates
\((\mathbb Z/m\mathbb Z)^\times\), so \(d\equiv2^j\pmod m\) and

\[
 1+\zeta^d=(1+\zeta)^{2^j}.
\]

If \(3\mid d\), then \(1+\zeta^d\in\F_{2^{2n/3}}^\times\).  This
group has order coprime to \(\ell\), because
\(\ord_\ell(2)=n\), so all its elements are \(\ell\)-th powers in
\(K\).  This proves (iii)\(\Rightarrow\)(iv); the converse is the
case \(a=0,b=1\).

Finally, all terms in \((\mathsf{CS})\) have complex absolute value
one.  Part (iv) makes them all one; conversely equality with the
positive real number \(m-1\) forces every term to be one by equality
in the triangle inequality.
\end{proof}

\begin{remark}[the natural secant Kummer cover is rational]
Over \(\F_4\), the Möbius function

\[
 M(T)=\frac{\omega+\omega^2T}{1+T}
\]

which expresses the cyclotomic secant has Kummer cover

\[
 Z^\ell=M(T),\qquad
 T=\frac{Z^\ell+\omega}{Z^\ell+\omega^2}.
\]

Its function field is therefore \(\F_4(Z)\), of genus zero.  A
positive-genus Weil estimate on this cover cannot rule out the
distinguished torsion fibre; the remaining information is its actual
Artin value.
\end{remark}

\begin{proposition}[Capelli form of the surviving lemma]
\label{prop:capelli}
\PROVED{} Put
\[
 A(X)=X^3+X+1,\qquad B(X)=X^2+X,\qquad
 R(X)=\frac{A(X)}{B(X)}.
\]
Let $Q_0(X)=X$ and, for $k\geq1$, define the monic binary polynomial
\[
 Q_k(X)=B(X)^{3^{k-1}}
        Q_{k-1}\!\left(R(X)\right).                         \tag{\(\star\)}
\]
Then $Q_k$ is the minimal polynomial of $\eta_k$ over $\F_2$ and is
irreducible of degree $3^k$.  Its first four values are
\begin{align*}
 Q_1(X)&=X^3+X+1,\\
 Q_2(X)&=X^9+X+1,\\
 Q_3(X)&=X^{27}+X^{25}+X^{24}+X^{19}+X^{18}+X^{16}
          +X^{10}+X^9+X^3+X+1,\\
 Q_4(X)&=X^{81}+X^{64}+X^{16}+X+1.
\end{align*}
For every prime $\ell\mid\Phi_{3^k}(2)$, the following are equivalent:
\begin{enumerate}[label=\textup{(\roman*)}]
\item $\chi_{k,\ell}(\eta_k)\neq1$;
\item $\eta_k$ is not an $\ell$-th power in
      $\F_{2^{3^k}}^{\times}$;
\item $X^{\ell}-\eta_k$ is irreducible over $\F_{2^{3^k}}$;
\item $Q_k(X^{\ell})$ is irreducible over $\F_2$.
\end{enumerate}
Thus the surviving character lemma is equivalently the assertion that
all the explicitly recursive compositions $Q_k(X^{\ell})$, for
$\ell\mid\Phi_{3^k}(2)$, are irreducible.  At the first composite
level $k=4$, the three targets are the monomial compositions of the
displayed pentanomial $Q_4$ by
$\ell\in\{2593,71119,97685839\}$.
\end{proposition}

\begin{proof}
Proposition~\ref{prop:eta-tower} gives $R(\eta_k)=\eta_{k-1}$ and an
irreducible cubic for $\eta_k$ over $\F_{2^{3^{k-1}}}$.  Since
$\eta_{k-1}\in\F_2(\eta_k)$, induction gives
$[\F_2(\eta_k):\F_2]=3^k$.  Clearing the denominator in the minimal
polynomial of $\eta_{k-1}$ gives \((\star)\), a monic polynomial of
degree $3^k$ having $\eta_k$ as a root: its leading term is
$A(X)^{3^{k-1}}=X^{3^k}+\cdots$.  It is therefore the minimal
polynomial.  Direct substitution in \((\star)\) gives the four displayed
polynomials.

Here $\Phi_{3^k}(2)=q_k^2+q_k+1\equiv1\pmod3$, so a prime
$\ell\mid\Phi_{3^k}(2)$ is not $3$.  The cyclotomic-divisor criterion
therefore gives $\ord_{\ell}(2)=3^k$, and hence
$\mu_{\ell}\subset\F_{2^{3^k}}$.  Since $\ell$ is prime, Kummer theory
says that $X^{\ell}-\eta_k$ either splits off a linear factor or is
irreducible, according as $\eta_k$ is or is not an $\ell$-th power.
This proves the equivalence of (ii) and (iii), while
Capelli's lemma gives the equivalence of (iii) and (iv).  The direct
product
$\F_{2^{3^k}}^{\times}\simeq
 \F_{q_k}^{\times}\times U_k$ and
$\ell\nmid(q_k-1)$ show that being an $\ell$-th power in the full
field is equivalent, for $\eta_k\in U_k$, to lying in $U_k^{\ell}$.
This is the equivalence with (i).
\end{proof}

\begin{remark}[Singer membership is not primitivity]
\label{rem:singer-not-primitive}
The difference-set identification cannot replace the character test.
Indeed, exponentiation
\[
 s_k:U_k\longrightarrow U_k,\qquad s_k(y)=y^{q_k+1},
\]
is an automorphism of order $6$: one has
$\gcd(q_k+1,|U_k|)=1$ and
$(q_k+1)^3\equiv-1\pmod{|U_k|}$, while
$(q_k+1)^2\equiv q_k\not\equiv1$.  The Singer equation is simply
$s_k(y)=y+1$.  Thus its consequence
$\ord(y+1)=\ord(y)$ is automatic order preservation by an
automorphism and supplies no order divisor by itself.
For example, take $q=32$ and let $y$ be a root of
$X^3+X^2+1$ over $\F_2$.  Then $y$ has order $7$ and
$q+1\equiv5\pmod7$; the defining cubic gives
$y^5+y+1=0$.  Hence $y\in\{x:x^{q+1}+x+1=0\}$, but $y$ does not
generate the norm-one group of order
$q^2+q+1=7\cdot151$.  What must be used in the Conway tower is the
recursive coefficient $\eta_{k-1}$ in \((*)\), not Singer membership
alone.
\end{remark}

\begin{proposition}[prime-step closure]\label{prop:prime-step}
\PROVED{} If $\Phi_{3^k}(2)$ is prime, then
\[
 C_k\iff C_{k-1}.
\]
Consequently $C_1,C_2,C_3$ hold analytically: the three new kernel
orders are respectively
\[
 7,\qquad73,\qquad262657,
\]
and all are prime.
\end{proposition}

\begin{proof}
Proposition~\ref{prop:eta-tower} proves that $\eta_k$ has degree $3$
over $\F_{q_k}$, so in particular $\eta_k\neq1$.  A nonidentity element
of a group of prime order generates it, and \((**)\) gives the first
claim.  The primality of $7$ and $73$ is immediate.  For completeness,
$262657-1=2^9\cdot3^3\cdot19$, and base $5$ gives the Pocklington
certificate
\[
 5^{262656}\equiv1\pmod{262657},\qquad
 \gcd\!\left(5^{262656/r}-1,262657\right)=1
 \quad(r=2,3,19),
\]
which proves the third primality assertion.
\end{proof}

\begin{conjecture}[$C_k$: $\gamma$-primitivity]\label{conj:ck}
$\gamma_k$ is primitive in $L^{\times}=\F_{2^{3^k}}^{\times}$.
\end{conjecture}

\begin{proposition}[equivalence]\label{prop:equiv}
\PROVED{} For primes $r$ with $f(r)\mid 3^k$ a $3$-power:
\[
  C_k
  \iff
  m_r=1\ \text{for every prime $r$ with } f(r)\in\{3,9,\dots,3^k\}.
\]
The family formula is the $0/1/4$ candidate of Section~\ref{sec:rule}
restricted to the $3$-power column.
\end{proposition}

\subsection{Certification status}

\PROVED{} Proposition~\ref{prop:prime-step} proves $C_k$ for
$k=1,2,3$ without an order computation.

\CERTIFIED{4\leq k\leq6} $C_k$ holds for $k=4,5,6$: $\gamma_k$ is primitive
in $\F_{2^{3^k}}^{\times}$, and
$\ord(\kappa_{3^k}+1)=3^{k+1}\,(2^{3^k}-1)$ is verified exactly --
product reconstruction, squarefreeness, primality by deterministic
Miller--Rabin below $3.3\times10^{24}$ and MR-64 probable-primality above,
an independent sieve re-deriving the small factors, cross-validation
against the independent term algebra for $k=1,2$ and against the
calculator rows $m_{2593}=1$, $m_{487}=1$
(\path{experiments/cyclotomic_3k_family.py},
\path{experiments/excess/}).

\CONSISTENT{} $k=7,8$: every known prime factor of $\Phi_{3^7}(2)$ and
$\Phi_{3^8}(2)$ divides $\ord(\gamma_k)$. Full certification is blocked
by the unfactored cofactors (FactorDB status composite-no-factor),
including their unknown prime multiplicities; no failure has been found.

\OPEN{} Whether ECM or GNFS can factor the $\Phi_{3^7}(2)$ and
$\Phi_{3^8}(2)$ cofactors within a realistic budget, converting $k=7,8$
from \textsc{consistent} to \textsc{certified}.

\subsection{Certified excess rows produced}

The following $m_r=1$ rows are certified by the order criterion
independently of the cofactor issue (analysis-level results; they are not
new Rust carries):
\begin{center}
\begin{tabular}{rl}
\toprule
$f$ & primes $r$ with $m_r=1$ \\
\midrule
$27$   & $262657$ \\
$81$   & $71119,\ 97685839$ \\
$243$  & $16753783618801,\ 192971705688577,\ 3712990163251158343$ \\
$729$  & $80191,\ 97687,\ 379081,\ \mathrm{P}42,\ \mathrm{P}90$ \\
$2187$ & $39367,\ 7606246033,\ 263196614521,\ 529063556041$ \\
$6561$ & $209953,\ 1299079,\ 70063267397606709277393$ \\
\bottomrule
\end{tabular}
\end{center}
($\mathrm{P}42$, $\mathrm{P}90$ denote the certified/probable primes of
$42$ and $90$ decimal digits dividing $\Phi_{3^6}(2)$.)

\section{The $0/1/4$ candidate}\label{sec:rule}

\begin{conjecture}[the $0/1/4$ rule]\label{conj:rule}
\CONSISTENT{} (not proved)
\[
  m_p=
  \begin{cases}
    0 & \text{if $\Qset(f(p))$ is not a singleton odd prime-power,}\\
    4 & \text{if $f(p)=2\cdot3^k$, $k\geq1$,}\\
    1 & \text{otherwise (singleton odd prime-power $\Qset(f(p))$).}
  \end{cases}
\]
The rule matches all $950$ calculator records with known $\Qset$-sets and
every OEIS-known row covered by those $\Qset$-sets.
\end{conjecture}

\subsection{Two maximal-order problems already contained in the rule}

The candidate contains strong explicit maximal-order assertions even
before the non-cyclotomic singleton chains are reached.  This prevents a
proof from following from field degree or irreducibility alone.

\begin{proposition}[the singleton-even arm is Conway--Fermat maximality]
\label{prop:fermat-arm}
\PROVED{} For $n\geq0$, put
\[
 h_n=2^{n+1},\qquad c_n=\kappa_{h_n}=2^{2^n},\qquad
 E_n=\F_{2^{h_n}},\qquad F_n=2^{2^n}+1.
\]
With $E_{-1}=\F_2$, let
\[
 \delta_n=\ord\bigl(c_nE_{n-1}^{\times}\bigr)
 \quad\text{in}\quad E_n^{\times}/E_{n-1}^{\times}.
\]
Then $1<\delta_n\mid F_n$, and the $m=0$ arm of
Conjecture~\ref{conj:rule} for every prime $p\mid F_n$ is equivalent to
\[
  \delta_n=F_n.
\]
\end{proposition}

\begin{proof}
The Conway finite-nimber tower satisfies
\[
 c_n^2+c_n=c_{n-1}c_{n-2}\cdots c_0
\]
(with the empty product $1$).  Thus $E_n/E_{n-1}$ is quadratic and
$c_n\notin E_{n-1}$, so the coset order is nontrivial.  Its ambient
quotient has order
\[
 \frac{2^{2^{n+1}}-1}{2^{2^n}-1}=F_n,
\]
which proves $1<\delta_n\mid F_n$; equivalently the relative norm is
$c_n^{F_n}=c_{n-1}\cdots c_0$.

If $p\mid F_n$, then $2^{2^n}\equiv-1\pmod p$, hence
$f(p)=2^{n+1}=h_n$.  Fermat numbers are pairwise coprime, so the entire
$p$-primary part of $|E_n^{\times}|$ lies in the quotient by
$E_{n-1}^{\times}$.  The full power criterion therefore gives
$m_p=0$ iff
$v_p(\delta_n)=v_p(F_n)$.  Requiring this for every prime divisor of
$F_n$ is exactly $\delta_n=F_n$.
\end{proof}

\begin{proposition}[Fibonacci-polynomial criterion]
\label{prop:fermat-fibotomic}
\PROVED{} Fix $n\geq0$, put $m=2^n$, $q=2^m$, and write
$K=E_{n-1}=\F_q$, $L=E_n=\F_{q^2}$ and
$a=a_{n-1}=\prod_{j<n}c_j$.  If
\[
 S_0(X)=0,\qquad S_1(X)=1,\qquad
 S_{r+2}(X)=S_{r+1}(X)+XS_r(X),
\]
then
\[
 \delta_n=\min\{r>0:S_r(a)=0\},
\]
and hence
\[
 \delta_n=F_n
 \iff
 S_{F_n/\ell}(a)\ne0
 \quad\hbox{for every prime }\ell\mid F_n.       \tag{\(\mathsf F_n\)}
\]

More precisely, for odd $d>1$ define the fibotomic polynomial
\[
 \mathfrak F_d(X)=
 \prod_{\substack{\{v,v^{-1}\}\\ \ord(v)=d}}
 \left(X-\frac{v}{(v+1)^2}\right),
\]
where the product takes one representative of each inverse-pair of
primitive $d$-th roots.  Then $\mathfrak F_d\in\F_2[X]$,
\[
 S_r(X)=\prod_{\substack{d\mid r\\d>1}}\mathfrak F_d(X)
 \quad(r\text{ odd}),                                      \tag{\(\mathsf B_r\)}
\]
and, if $A_{n-1}$ is the minimal polynomial of $a_{n-1}$ over $\F_2$,
\[
 \delta_n=d\iff A_{n-1}\mid\mathfrak F_d.
\]
Every irreducible factor of every $\mathfrak F_d$ with
$1<d\mid F_n$ has the same degree $m$.  Thus the Conway--Fermat
conjecture is exactly the assertion
$A_{n-1}\mid\mathfrak F_{F_n}$; field degree cannot distinguish that
stratum from any proper-divisor stratum.
\end{proposition}

\begin{proof}
Put $c=c_n$ and $u=(c+1)/c$.  Since $c^2+c=a$ and
$c^q=c+1$, the element $u$ lies in the norm-one group of order
$q+1=F_n$, and $\delta_n=\ord(u)$.  The two roots $c,c+1$ of
$T^2+T+a$ give the Binet identity
\[
 S_r(a)=c^r+(c+1)^r=\frac{1+u^r}{(1+u)^r}.                 \tag{\(\mathsf B\)}
\]
This proves the first criterion.  In a cyclic group of order $F_n$ an
element has full order precisely when its $F_n/\ell$-th power is
nontrivial for every prime $\ell\mid F_n$, proving
\((\mathsf F_n)\).

For odd $r$, the nontrivial $r$-th roots $v$ occur in inverse pairs,
and the common value $v/(v+1)^2$ of a pair is, by
\((\mathsf B)\), a root of $S_r$.  Frobenius permutes the pairs, so
the displayed factors have coefficients in $\F_2$.  Both sides of
\((\mathsf B_r)\) are monic of degree $(r-1)/2$, which proves the
factorization and the divisibility criterion.

Finally $2^m\equiv-1\pmod d$ for every nontrivial $d\mid F_n$.
Because $m$ is a power of two, $\ord_d(2)=2m$, and the least $j>0$
with $2^j\equiv\pm1\pmod d$ is $m$.  This is exactly the Frobenius
orbit length of $v/(v+1)^2$: indeed
\[
 \frac{v}{(v+1)^2}=\frac{w}{(w+1)^2}
 \iff (v+w)(vw+1)=0,
\]
so only the inverse pair $\{v,v^{-1}\}$ has the same image.  This
proves the equal-degree assertion.
\end{proof}

\begin{proposition}[coarse top-step data cannot force maximality]
\label{prop:fermat-countermodels}
\PROVED{} For every divisor $d>1$ of $F_n$ there are elements
$a_d\in\F_q$ and $c_d\in\F_{q^2}$ such that
\[
 [\F_2(a_d):\F_2]=m,\qquad
 [\F_2(c_d):\F_2]=2m,\qquad
 c_d^q=c_d+1,\qquad c_d^2+c_d=a_d,
\]
\[
 \operatorname{Tr}_{\F_q/\F_2}(a_d)=1,\qquad
 \Norm_{\F_{q^2}/\F_q}(c_d)=a_d,
\]
but the order of $c_d\F_q^\times$ is exactly $d$.  In particular all
proper-divisor strata obey the same degree, trace, Artin--Schreier and
relative-norm data as the Conway top step.
\end{proposition}

\begin{proof}
Choose $v$ of order $d$ in the norm-one subgroup of
$\F_{q^2}^{\times}$ and put
\[
 c_d=\frac1{v+1},\qquad a_d=\frac{v}{(v+1)^2}.
\]
Then $v^q=v^{-1}$, so $c_d^q=c_d+1$; the remaining displayed
identities follow directly.  The orbit calculation in the proof of
Proposition~\ref{prop:fermat-fibotomic} gives the two field degrees.
Consequently $T^2+T+a_d$ is irreducible over $\F_q$, which is
equivalent to the stated absolute trace, and
$(c_d+1)/c_d=v$ gives quotient order $d$.
\end{proof}

\begin{proposition}[the exact recursively selected factor]
\label{prop:fermat-resultant}
\PROVED{} Set $a_{-1}=1$, $a_i=\prod_{j\leq i}c_j$, and let $A_i$ be
its minimal polynomial over $\F_2$, with $A_{-1}(Y)=Y+1$.  Then
\[
 a_i^2+a_{i-1}a_i+a_{i-1}^3=0
\]
and
\[
 A_i(X)=\operatorname{Res}_Y
 \bigl(A_{i-1}(Y),X^2+YX+Y^3\bigr).                       \tag{\(\mathsf R_i\)}
\]
Hence the unresolved Conway--Fermat lemma is the precise recursive
claim that the factor selected by \((\mathsf R_i)\) lies in the
full-order fibotomic stratum, rather than in any equal-degree
$\mathfrak F_d$ with $d\mid F_{i+1}$ and $d<F_{i+1}$.
\end{proposition}

\begin{proof}
Multiplying $c_i^2+c_i=a_{i-1}$ by $a_{i-1}^2$ gives the first
identity.  We prove the degree assertion and the resultant formula
inductively from $a_{-1}=1$.  The resultant in \((\mathsf R_i)\) is monic of degree
$2\deg A_{i-1}$ and vanishes at $a_i$.  Since
$c_i=a_i/a_{i-1}\notin E_{i-1}$, one has $a_i\notin E_{i-1}$.
But $E_i/\F_2$ has degree $2^{i+1}$, and every proper subfield of
$E_i$ is contained in its unique index-two subfield $E_{i-1}$.
Therefore $a_i$ has full degree $2^{i+1}=2\deg A_{i-1}$, so the
resultant is its minimal polynomial.
\end{proof}

\begin{proposition}[one-step maximality is not an inductive invariant]
\label{prop:fermat-one-step-no-go}
\PROVED{} Let

\[
 C(Y)=Y^{16}+Y^{15}+Y^{12}+Y^{11}+Y^{10}+Y^8+1
\]

and apply the exact Conway resultant operator of
Proposition~\ref{prop:fermat-resultant}:

\[
\begin{split}
 B(X)&=\operatorname{Res}_Y(C(Y),X^2+YX+Y^3)\\
 &=X^{32}+X^{31}+X^{30}+X^{27}+X^{26}+X^{24}
   +X^{23}+X^{22}\\
 &\qquad{}+X^{19}+X^{15}+X^{14}+X^{11}+X^8+X^7+1.
\end{split}                                                    \tag{\(\mathsf{NI}_1\)}
\]

Then \(C\) is an irreducible factor of the full-order fibotomic
stratum \(\mathfrak F_{F_4}\), where \(F_4=65537\), whereas \(B\) is
an irreducible factor of the proper-order stratum

\[
 \mathfrak F_{6700417},\qquad
 6700417=F_5/641.
                                                               \tag{\(\mathsf{NI}_2\)}
\]

Thus the exact one-step resultant can send a full-order predecessor
factor to a proper-order target factor, already at the first composite
Fermat level.  In particular, predecessor maximality plus the transition
formula alone is not an induction theorem.  This is not a counterexample
to the distinguished Conway path: \(C\) is an alternative factor at the
\(A_3\) level and need not have the earlier ancestry of the actual
\(A_3\).
\end{proposition}

\begin{proof}
First, exact modular squaring gives

\[
 Y^{2^{16}}\equiv Y\pmod C,
\]

and

\[
 Y^{256}+Y\equiv
 R:=Y^{15}+Y^{14}+Y^{13}+Y^9+Y^8+Y^2+Y\pmod C.
\]

The Bézout identity \(UC+VR=1\) holds for

\[
\begin{split}
 U={}&Y^{14}+Y^{13}+Y^{12}+Y^{11}+Y^7+Y^6+Y^5+Y^4\\
    &\quad{}+Y^3+Y+1,\\
 V={}&Y^{15}+Y^{14}+Y^{12}+Y^{11}+Y^9+Y^8+Y^5+Y^3+Y+1.
\end{split}
\]

Rabin's criterion, for which two is the only prime divisor of sixteen,
therefore proves that \(C\) is irreducible.

The Fibonacci polynomials satisfy, in characteristic two,

\[
 S_{2r}=S_r^2,\qquad
 S_{2r+1}=S_{r+1}^2+YS_r^2,
\]

and hence

\[
 S_{2^m+1}(Y)=1+\sum_{j=0}^{m-1}Y^{2^j}.
                                                               \tag{\(\mathsf{NI}_3\)}
\]

For a root \(y\) of \(C\), the coefficient of \(Y^{15}\) gives
\(\operatorname{Tr}_{\F_{2^{16}}/\F_2}(y)=1\).  Equation
\((\mathsf{NI}_3)\) therefore gives
\(S_{65537}(y)=0\).  Since \(65537\) is prime, \(y\) belongs to the
full-order stratum \(\mathfrak F_{65537}\).

Now choose \(z\) with \(z^2+z=y\) and put \(x=yz\).
The trace just computed makes \(Z^2+Z+y\) irreducible over
\(\F_{2^{16}}\).  Thus \(x\notin\F_{2^{16}}\); because every proper
subfield of \(\F_{2^{32}}\) lies in \(\F_{2^{16}}\), \(x\) has degree
thirty-two over \(\F_2\).  Moreover

\[
 x^2+yx+y^3=y^2(z^2+z+y)=0.                          \tag{\(\mathsf{NI}_4\)}
\]

Modulo the relation in \((\mathsf{NI}_4)\), direct reduction gives

\[
 B(X)=C(Y)\bigl(Q_1(Y)X+Q_0(Y)\bigr),
                                                               \tag{\(\mathsf{NI}_5\)}
\]

where

\[
\begin{split}
 Q_1={}&Y^{29}+Y^{28}+Y^{25}+Y^{24}+Y^{21}+Y^{18}
       +Y^{17}+Y^{16}\\
      &\quad{}+Y^{15}+Y^{14}+Y^{11}+Y^{10}+Y^9+Y^6,\\
 Q_0={}&Y^{32}+Y^{31}+Y^{29}+Y^{28}+Y^{27}+Y^{26}
       +Y^{25}+Y^{23}\\
      &\quad{}+Y^{20}+Y^{14}+Y^{12}+Y^{11}+Y^{10}+Y^9+1.
\end{split}
\]

Thus \(B(x)=0\).  Its displayed degree is thirty-two, so it is the
minimal polynomial of \(x\), and consequently the stated resultant.

It remains to identify the target stratum.  From a pair
\((f,g)=(S_r,S_{r+1})\), put

\[
 A=f^2,\qquad D=g^2+Xf^2.
\]

Reading one further binary digit sends the pair to \((A,D)\) for a
zero and to \((D,g^2)\) for a one.  Repeated reduction modulo \(B\)
along the binary word

\[
 6700417=(11001100011110110000001)_2
\]

ends at the exact residue pair

\[
 \bigl(S_{6700417},S_{6700418}\bigr)
 \equiv
 \bigl(\mathtt{00000000},\mathtt{1875a343}\bigr)
 \pmod B,                                           \tag{\(\mathsf{NI}_6\)}
\]

where a hexadecimal bit in position \(j\) is the coefficient of \(X^j\).
This is an \(O(\log 6700417)\) polynomial-identity certificate, not
field sampling.  Hence \(S_{6700417}(x)=0\).  The integer \(6700417\)
is prime, so \(x\) belongs to the exact order-\(6700417\) fibotomic
stratum.  Finally, if \(v=(z+1)/z\), then

\[
 y=\frac{v}{(v+1)^2},\qquad
 x=yz=\frac{v}{(v+1)^3},
\]

showing directly that this is the same fibotomic/elliptic transition as
the Conway resultant.
\end{proof}

\begin{proposition}[complete ancestry and exact trace-fibre exhaustion]
\label{prop:fermat-ancestry-exhaustion}
\PROVED{} For \(i\geq0\) and \(-1\leq j<i\), the distinguished
elements satisfy
\[
 \operatorname{Tr}_{E_i/E_j}(a_i)=a_j,\qquad
 \operatorname{Norm}_{E_i/E_j}(a_i)=a_j^{3^{i-j}}.
                                                               \tag{\(\mathsf{AN}_1\)}
\]
Conversely, fix \(i\geq0\) and \(x_i\in E_i\).  Put
\[
 x_j=\operatorname{Tr}_{E_i/E_j}(x_i)\qquad(-1\leq j<i).
\]
If \(x_{-1}=1\) and
\[
 \operatorname{Norm}_{E_j/E_{j-1}}(x_j)=x_{j-1}^3
 \qquad(0\leq j\leq i),                                \tag{\(\mathsf{AN}_2\)}
\]
then \(x_i\) is an \(\mathbb F_2\)-Frobenius conjugate of \(a_i\).
Thus \((\mathsf{AN}_1)\)--\((\mathsf{AN}_2)\) are a
necessary-and-sufficient ancestry invariant for the actual resultant
factor.

Now fix \(n\geq1\), put
\[
 F=E_{n-1}=\mathbb F_q,\qquad K=E_n=\mathbb F_{q^2},
\]
and let \(\ell\mid q+1\).  Set
\[
 H=(K^\times)^\ell,\qquad
 U=\ker\operatorname{Norm}_{K/F}.
\]
For every \(t\in F^\times\),
\[
 \#\{h\in H:\operatorname{Tr}_{K/F}(h)=t\}
   =\frac{q+1}{\ell}-1.                                \tag{\(\mathsf{AN}_3\)}
\]
If \(t=a_{n-1}\), every element counted in
\((\mathsf{AN}_3)\) has absolute trace one and is a normal, hence
full-degree, element of \(K/\mathbb F_2\).  Thus, whenever \(q+1\)
is composite, full degree, normality, the complete Frobenius-orbit
length, and the actual relative-trace value all occur among
\(\ell\)-th powers.

The normalized relative norms of the elements in
\((\mathsf{AN}_3)\) form exactly
\[
 \mathcal B_\ell=
 \left\{\frac{u}{(u+1)^2}:u\in U^\ell,\ u\ne1\right\},
 \qquad
 |\mathcal B_\ell|
 =\frac12\left(\frac{q+1}{\ell}-1\right).               \tag{\(\mathsf{AN}_4\)}
\]
More precisely, an element \(h\) of relative trace \(t\) is an
\(\ell\)-th power and has relative norm \(t^3\) if and only if
\(t\in\mathcal B_\ell\).  For \(t=a_{n-1}\) this is equivalent to
\[
 S_{(q+1)/\ell}(a_{n-1})=0.                            \tag{\(\mathsf{AN}_5\)}
\]
Finally \([a_n]=[c_n]\) in \(K^\times/K^{\times\ell}\).
Consequently the entire earlier ancestry is recoverable exactly, but
after the last quadratic step its power-residue content is precisely
the original fibotomic membership \((\mathsf{AN}_5)\).  Trace,
separate norm, normality, and Frobenius-orbit invariants give no
additional obstruction.
\end{proposition}

\begin{proof}
The relation
\[
 a_i^2+a_{i-1}a_i+a_{i-1}^3=0
\]
is the minimal polynomial of \(a_i\) over \(E_{i-1}\).  It gives
immediate trace \(a_{i-1}\) and norm \(a_{i-1}^3\); transitivity gives
\((\mathsf{AN}_1)\).

For the converse, transitivity of trace gives
\(\operatorname{Tr}_{E_j/E_{j-1}}(x_j)=x_{j-1}\).
Together with \((\mathsf{AN}_2)\) this says that \(x_j\) is a root of
\[
 X^2+x_{j-1}X+x_{j-1}^3.
\]
Starting from \(x_{-1}=a_{-1}=1\), induct on \(j\).  If
\(x_{j-1}\) is a Frobenius conjugate of \(a_{j-1}\), the displayed
polynomial is the corresponding Frobenius conjugate of the polynomial
for \(a_j\); its two roots are the two lifts of that conjugacy to
\(E_j\).  Hence \(x_j\), and finally \(x_i\), is a Frobenius conjugate
of the distinguished element.

Since \(\gcd(\ell,q-1)=1\), one has \(F^\times\subset H\).
Multiplication by any \(t\in F^\times\) bijects the trace-one and
trace-\(t\) fibres in \(H\).  In characteristic two,
\(\operatorname{Tr}_{K/F}(h)=0\) is equivalent to \(h^q=h\), so the
nonzero trace-zero elements are exactly \(F^\times\), all of which lie
in \(H\).  If the common nonzero-fibre size is \(M\), then
\[
 \frac{q^2-1}{\ell}=|H|=(q-1)+(q-1)M,
\]
which gives \((\mathsf{AN}_3)\).

The distinguished \(a_{n-1}\) has absolute trace one by
\((\mathsf{AN}_1)\).  Thus every \(h\) in its relative-trace fibre
also has absolute trace one.  To justify normality, let
\(N=[K:\mathbb F_2]\), a power of two, and let \(\rho\) generate its
cyclic Galois group.  After choosing one normal generator, \(K\) is
the regular module over
\[
 \mathbb F_2[\langle\rho\rangle]
 \simeq\mathbb F_2[X]/((X+1)^N),
\]
a local ring.  An element is a normal generator exactly when its
group-ring coordinate is a unit, equivalently when multiplication by
\((\rho+1)^{N-1}\) does not kill it.  Since
\[
 (\rho+1)^{N-1}=1+\rho+\cdots+\rho^{N-1},
\]
this last value is its absolute trace.  Hence absolute trace one is
equivalent to normality.

Write \(h=tz\) in a nonzero trace fibre.  Then
\(\operatorname{Tr}_{K/F}(z)=1\), so \(z^q=z+1\) and
\[
 \operatorname{Norm}_{K/F}(h)=t^2z(z+1).
\]
Put \(u=(z+1)/z\in U\); then \(z=1/(u+1)\) and
\[
 z(z+1)=\frac{u}{(u+1)^2}.                            \tag{\(\mathsf{AN}_6\)}
\]
We claim \(z\in H\) if and only if \(u\in U^\ell\).  The forward
implication follows from \(u=z^{q-1}\).  Conversely, if
\(u\in U^\ell\), then \(u\in H\).  The right side of
\((\mathsf{AN}_6)\) lies in \(F^\times\subset H\), so in
\(K^\times/H\) one has
\[
 [u]=2[u+1].
\]
Because \(\ell\) is odd, multiplication by two is invertible; hence
\([u]=0\) forces \([u+1]=0\), and
\(z=(u+1)^{-1}\in H\).  The map
\(u\mapsto u/(u+1)^2\) identifies precisely the inverse pair
\(\{u,u^{-1}\}\), proving \((\mathsf{AN}_4)\).

The equality \(\operatorname{Norm}(h)=t^3\) is, after division by
\(t^2\), exactly \(t=z(z+1)\).  This proves the first assertion
following \((\mathsf{AN}_4)\).  The set \(U^\ell\) has order
\((q+1)/\ell\), and the Fibonacci-polynomial factorization identifies
its nontrivial inverse pairs with the roots of
\(S_{(q+1)/\ell}\), proving \((\mathsf{AN}_5)\).  Finally
\(a_n=a_{n-1}c_n\) and \(a_{n-1}\in F^\times\subset H\), so
\([a_n]=[c_n]\).
\end{proof}

\begin{proposition}[supersingular elliptic form of the Conway transition]
\label{prop:fermat-elliptic}
\PROVED{} For $i\geq0$ put
\[
 u_i:=\frac{c_i+1}{c_i}.
\]
Fix $n\geq1$, set $Q=2^{2^{n-1}}$, $v=u_{n-1}$ and $u=u_n$.
Then the exact Conway recursion is
\[
 \frac{u}{(u+1)^2}=\frac{v}{(v+1)^3},                    \tag{\(\mathsf E\)}
\]
or equivalently
\[
 u(v+1)^3=v(u+1)^2,
 \qquad
 u+u^{-1}=v^2+v+1+v^{-1}.
\]
The normalization of this correspondence is the fixed supersingular
elliptic curve over $\F_2$
\[
 \mathcal E:\qquad y^2+y=x^3+x^2.                       \tag{\(\mathsf W\)}
\]
The maps to the $v$- and $u$-lines have degrees two and three,
respectively; in particular the transition has no rational section
recovering $u$ from $v$.

On the original $(u,v)$ correspondence, the rational map
\[
 \mathcal A(u,v)=\left(\frac{u+v}{uv+1},v^{-1}\right)
\]
extends to an automorphism of order four, with
\[
 \mathcal A^2(u,v)=(u^{-1},v).
\]
For the distinguished Conway point $P_n=(u_n,u_{n-1})$ one has
\[
 \Frob_Q(P_n)=\mathcal A(P_n)
 \quad\hbox{or}\quad
 \mathcal A^{-1}(P_n).                                   \tag{\(\mathsf Q\)}
\]
Thus, at \(P_n\), relative Frobenius is one of the two quarter-turns
\(\mathcal A^{\pm1}\) on the elliptic correspondence.
\end{proposition}

\begin{proof}
Write $a_i=\prod_{j\leq i}c_j$.  The identities
\[
 a_{n-1}=\frac{u}{(u+1)^2},\qquad
 a_{n-1}=a_{n-2}c_{n-1}
 =\frac{v}{(v+1)^2}\frac1{v+1}
\]
give \((\mathsf E)\); inversion gives the second displayed form.

Put $U=u+1$, $V=v+1$, and blow up the affine singularity
\((U,V)=(0,0)\) by
$X=U/V$, $Y=V$.  After division by $V^2$ the equation becomes
\[
 XY^2+(X^2+1)Y+X^2=0.
\]
Its homogenization
\[
 H=XY^2+(X^2+Z^2)Y+X^2Z
\]
has partial derivatives $(Y^2,X^2+Z^2,X^2)$ and is therefore smooth.
The point $[1:1:0]$ is a flex.  The linear change
\[
 (x,y,z)=(Z,X,X+Y+Z)
\]
takes $H=0$ to
$y^2z+yz^2=x^3+x^2z$, which is \((\mathsf W)\).  This curve has
five \(\F_2\)-points, hence trace \(-2\), and is supersingular.  The original
equation has degrees two in $u$ and three in $v$, proving the degrees
of the two projections.  A rational section of the nontrivial
degree-two finite cover would induce a retraction of its degree-two
function-field extension, impossible for a proper finite field
extension.

Direct substitution in \((\mathsf E)\) gives
\[
 \frac{\mathcal A_u}{(\mathcal A_u+1)^2}
 =v\frac{u}{(u+1)^2}
 =\frac{v^{-1}}{(v^{-1}+1)^3},
\]
so $\mathcal A$ preserves the correspondence.  A second substitution
gives $\mathcal A^2(u,v)=(u^{-1},v)$, proving the order assertion.
Thus this birational self-map and its birational inverse extend over
the smooth projective normalization.

Finally $v^Q=v^{-1}$ and
\[
 \left(\frac{u}{(u+1)^2}\right)^Q
 =\left(\frac{v}{(v+1)^3}\right)^Q
 =v\frac{u}{(u+1)^2}.
\]
The fibres of $z\mapsto z/(z+1)^2$ are the inverse-pairs
$\{z,z^{-1}\}$.  The first coordinate of $\Frob_Q(P_n)$ is therefore
either $\mathcal A_u(P_n)$ or its inverse, while the second coordinate
is $v^{-1}$, proving \((\mathsf Q)\).
\end{proof}

\begin{proposition}[new-prime character neutrality and the exact Kummer target]
\label{prop:fermat-kummer}
\PROVED{} Retain Proposition~\ref{prop:fermat-elliptic} and let
$\ell$ be a prime divisor of $Q^2+1=F_n$.  If
$\chi:E_n^\times=\F_{Q^4}^\times\to\mu_\ell$ has exact order
$\ell$, and
\[
 \alpha=\chi(u),\qquad \beta=\chi(u+1),
\]
then
\[
 \alpha=\beta^2,\qquad \beta^{Q^2}=\beta^{-1}.           \tag{\(\mathsf K\)}
\]
The second identity is tautological because
$Q^2\equiv-1\pmod\ell$; in particular the exact lower transition and
the quarter-turn Frobenius symmetry do not by themselves exclude
$\beta=1$.

The following are equivalent:
\begin{enumerate}[label=\textup{(\roman*)}]
\item the order $\delta_n$ contains the full $\ell$-primary part of
      $F_n$;
\item $u_n\notin U_n^\ell$, where
      $U_n=\ker\Norm_{E_n/E_{n-1}}$;
\item $c_n\notin(E_n^\times)^\ell$;
\item $X^\ell-c_n$ is irreducible over $E_n$.
\end{enumerate}
Consequently the all-level Conway--Fermat theorem is exactly the
nontriviality, for every $\ell\mid F_n$, of this new Kummer symbol at
the distinguished point of \((\mathsf W)\).  No lower-field
multiplicative character can supply its value.
\end{proposition}

\begin{proof}
The lower field has order $Q^2$, and
$\gcd(\ell,Q^2-1)=1$, so $\chi$ is trivial on all of
$E_{n-1}^\times$, in particular on every expression in $v$.  Applying
$\chi$ to \((\mathsf E)\) gives $\alpha=\beta^2$.  Since
$u^{Q^2}=u^{-1}$,
\[
 (u+1)^{Q^2}=u^{-1}+1=\frac{u+1}{u},
\]
which gives the other identity in \((\mathsf K)\).  It reduces to
$\beta^{-1}=\beta^{-1}$ modulo $\ell$.

In the cyclic group $U_n$ of order $F_n$, nonmembership in
$U_n^\ell$ is equivalent to having the full $\ell$-primary order.
Here
\[
 U_n^\ell=U_n\cap(E_n^\times)^\ell.                    \tag{\(\mathsf K_0\)}
\]
Indeed, if \(u=x^\ell\in U_n\), then
\(\Norm_{E_n/E_{n-1}}(x)^\ell=1\); because
\(\gcd(\ell,Q^2-1)=1\), this forces \(\Norm(x)=1\), so \(x\in U_n\).
The reverse inclusion is immediate.
Moreover \((\mathsf E)\) lies in the lower field, hence is an
$\ell$-th power in $E_n$.  In
$E_n^\times/(E_n^\times)^\ell$ it therefore says
\[
 [u]=2[u+1].
\]
Since $\ell$ is odd, $u$ is an $\ell$-th power exactly when $u+1$ is.
But $u+1=c_n^{-1}$, proving the equivalence of the first three
assertions.  Finally $\mu_\ell\subset E_n$, so the prime-degree Kummer
criterion proves the equivalence with irreducibility of
$X^\ell-c_n$.
\end{proof}

\begin{proposition}[four-point saturation of a Fermat failure]
\label{prop:fermat-four-plane}
\PROVED{} Fix \(n\geq2\), put

\[
 c=c_n,\qquad d=c_{n-1},\qquad b=a_{n-2},\qquad
 Q=|E_{n-2}|=2^{2^{n-1}}.
\]

Then, for one \(\epsilon\in\F_2\),

\[
 c^Q=c+d+\epsilon,\qquad c^{Q^2}=c+1.              \tag{\(\mathsf{FP}_1\)}
\]

Consequently, if \(\ell\mid F_n\) and \(c\) is an \(\ell\)-th power
in \(E_n\), the complete affine plane

\[
 c+\F_2+\F_2d=\{c,c+1,c+d,c+d+1\}                 \tag{\(\mathsf{FP}_2\)}
\]

consists of \(\ell\)-th powers.  Its product is nevertheless the
old-field element

\[
 \Norm_{E_n/E_{n-2}}(c)=b^3,                        \tag{\(\mathsf{FP}_3\)}
\]

which is automatically an \(\ell\)-th power.  Thus the exact four-point
saturation and its norm are compatible with failure; they do not supply
the missing nonresidue.
\end{proposition}

\begin{proof}
The tower equations give

\[
 d^2+d=b,\qquad c^2+c=bd=d^3+d^2,\qquad d^Q=d+1.
\]

For \(s=c^Q+c\), subtraction of the equation for \(c\) from its
\(Q\)-th power gives \(s^2+s=b\).  Hence
\(s\in\{d,d+1\}\), proving the first identity in
\((\mathsf{FP}_1)\); applying \(Q\)-Frobenius once more gives the
second.  Frobenius preserves the subgroup of \(\ell\)-th powers, so
one power value forces all four values in \((\mathsf{FP}_2)\).
Finally,

\[
 c(c+1)=bd,\qquad
 (c+d)(c+d+1)=b(d+1),
\]

and their product is \(b^2d(d+1)=b^3\).  Since
\(\ell\mid Q^2+1\) is coprime to \(|E_{n-2}^\times|=Q-1\), every
old-field element is an \(\ell\)-th power in \(E_n\).
\end{proof}

\begin{proposition}[CM form and the Kummer-cover obstruction]
\label{prop:fermat-elliptic-kummer}
\PROVED{} On \(\mathcal E:y^2+y=x^3+x^2\), let \(O\) be the origin
and \(T=(0,0)\).  Then
\(\mathcal E(\F_2)=\langle T\rangle\cong\mathbb Z/5\mathbb Z\), with
\[
 -T=(0,1),\qquad 2T=(1,1),\qquad -2T=(1,0).
\]
On the normalized Conway correspondence, the function
\(\mathfrak c=1/(u+1)\) is
\[
 \mathfrak c=\frac{x}{x^2+x+y},\qquad
 \operatorname{div}(\mathfrak c)=[-T]+2[O]-3[-2T].     \tag{\(\mathsf{EC}_1\)}
\]
The automorphism \(\mathcal A\) satisfies
\[
 \mathcal A^2(P)=T-P,\qquad
 \mathfrak c(\mathcal A^2P)=\mathfrak c(P)+1.          \tag{\(\mathsf{EC}_2\)}
\]
If \(\pi\) is \(2\)-power Frobenius and \(i=\pi+1\), then
\(i^2=-1\), \(i(T)=2T\), and, after possibly replacing
\(\mathcal A\) by its inverse,
\[
 \mathcal A=\tau_{2T}\circ i.
\]
Both orientations fix \(R_0=-2T\).

For \(n\geq4\), set
\[
 Q=2^{2^{n-1}},\qquad R=\sqrt Q=2^{2^{n-2}},
\]
and let \(P_n\) be the distinguished Conway point.  If
\(X_n=P_n-R_0\), then
\[
 \Frob_Q(X_n)=\pm iX_n,\qquad
 X_n\in\ker([R]\mp i).                                 \tag{\(\mathsf{EC}_3\)}
\]
Each kernel is cyclic of order \(R^2+1=Q+1\).

For every prime \(\ell\mid Q^2+1\), the Kummer cover
\[
 \mathcal C_\ell:\qquad z^\ell=\mathfrak c
\]
is geometrically connected, ramified exactly above
\(-T,O,-2T\), and has genus
\[
 g(\mathcal C_\ell)=\frac{3\ell-1}{2}.                 \tag{\(\mathsf{EC}_4\)}
\]
The automorphism \(\mathcal A\) does not lift to
\(\mathcal C_\ell\).  Hence the quarter-turn on the base supplies no
action on the Kummer fibre above \(P_n\); the surviving Conway
condition is exactly
\[
 \mathfrak c(P_n)=c_n\notin
 \bigl(\F_{Q^4}^{\times}\bigr)^\ell.                   \tag{\(\mathsf{EC}_5\)}
\]
It is the Artin-symbol problem for one distinguished fibre of this
ramified cover, not a consequence of the CM action on \(\mathcal E\).
\end{proposition}

\begin{proof}
In the blow-up coordinates of
Proposition~\ref{prop:fermat-elliptic}, \(U=u+1\), \(V=v+1\),
\(X=U/V\), \(Y=V\), and
\((x:y:z)=(Z:X:X+Y+Z)\).  On \(z=1\), homogeneous representatives give
\[
 V=\frac{1+x+y}{x},\qquad U=\frac{y(1+x+y)}{x^2}.
\]
Since \(y(1+x+y)=xy+x^3+x^2=x(x^2+x+y)\), this proves the formula for
\(\mathfrak c=1/U\).  With \(d=x^2+x+y\), direct intersection gives
\[
 \operatorname{div}(x)=[T]+[-T]-2[O],\qquad
 \operatorname{div}(d)=[T]+3[-2T]-4[O].
\]
Indeed \(d=0\) gives \(y=x^2+x\), whence the curve equation gives
\(x(x+1)^3=0\), and \(d\) has pole order four at \(O\).  Subtraction
proves \((\mathsf{EC}_1)\).

The function \(v=(y+1)/x\) has pole divisor \([O]+[T]\), so the deck
involution of its degree-two map is \(J(P)=T-P\).
Since \(\mathcal A^2\) fixes \(v\) and is nontrivial, it equals \(J\).
Also \(\mathcal A^2\) sends \(u\) to \(u^{-1}\), and therefore
\[
 \mathfrak c(\mathcal A^2P)
 =\frac1{u(P)^{-1}+1}
 =\frac{u(P)}{u(P)+1}
 =\mathfrak c(P)+1,
\]
proving \((\mathsf{EC}_2)\).

The Frobenius polynomial of \(\mathcal E/\F_2\) is
\(\pi^2+2\pi+2=0\).  Thus \(i=\pi+1\) satisfies \(i^2=-1\), and it
acts on \(\mathcal E(\F_2)\) as multiplication by two.
Write \(\mathcal A=\tau_S\circ\alpha\), with \(\alpha(O)=O\).
As \(\mathcal A\) is defined over \(\F_2\), \(\alpha\) commutes with
\(\pi\).  The centralizer of \(\pi\) in the supersingular endomorphism
algebra is \(\mathbb Q(i)\), whose integral units are
\(\{\pm1,\pm i\}\).  The linear part of \(\mathcal A^2\) is
\([-1]\), so \(\alpha=\pm i\).  Comparing translation parts gives
\[
 \mathcal A=\tau_{2T}\circ i
 \quad\text{or}\quad
 \mathcal A=\tau_{-T}\circ(-i);
\]
these are inverse and both fix \(R_0=-2T\).

For \(n\geq4\), \(2^{n-1}\) is divisible by eight.  Since
\(\pi^8=[16]\),
\[
 \pi^{2^{n-1}}=[2^{2^{n-2}}]=[R].
\]
Translation of \(R_0\) to the origin makes
\(\mathcal A^{\pm1}\) the linear maps \(\pm i\), proving
\((\mathsf{EC}_3)\).  The separable endomorphisms \([R]\mp i\) have
degree \(R^2+1\).  For every odd prime \(p\), the Tate module
\(T_p(\mathcal E)\) is free of rank one over \(\mathbb Z_p[i]\)
(the polynomial \(X^2+1\) is separable away from \(2\)).  Hence the
\(p\)-primary kernel is the additive quotient
\(\mathbb Z_p[i]/(R\mp i)\).  The matrix of \(R\mp i\) over
\(\mathbb Z_p\) has entries with gcd one and determinant \(R^2+1\),
so its first Smith invariant is one.  Combining the odd primary parts
shows that the geometric kernel is cyclic of order \(R^2+1\).

Every prime \(\ell\mid Q^2+1\) is coprime to \(15\).  Thus the
valuations \(1,2,-3\) in \((\mathsf{EC}_1)\) are nonzero modulo
\(\ell\).  Kummer theory makes \(z^\ell=\mathfrak c\) geometrically
connected and totally ramified at exactly those three points.
Riemann--Hurwitz gives
\[
 2g(\mathcal C_\ell)-2=3(\ell-1),
\]
which proves \((\mathsf{EC}_4)\).

It remains to prove non-lifting.  Choose
\(\mathcal A=\tau_{2T}\circ i\), put
\(B_j=\mathcal A^j(O)\), and retain \(R_0=-2T\).  Then
\[
 (B_0,B_1,B_2,B_3)=(O,2T,T,-T).
\]
For \(D_j=\operatorname{div}(\mathfrak c\circ\mathcal A^{-j})\),
pullback gives
\[
 D_j=2[B_j]+[B_{j+3}]-3[R_0].                          \tag{\(\mathsf{EC}_6\)}
\]
If \(\sum_{j=0}^3r_jD_j\equiv0\pmod\ell\), comparison at the four
\(B_j\)'s gives \(2r_j+r_{j+1}=0\) cyclically.  Iteration gives
\(15r_0=0\), and \(\ell\nmid15\) forces every \(r_j=0\).  Hence
\[
 \mathfrak c,\quad\mathfrak c\circ\mathcal A,\quad
 \mathfrak c\circ\mathcal A^2,\quad
 \mathfrak c\circ\mathcal A^3
\]
are linearly independent in
\[
 \overline{\F}_2(\mathcal E)^\times/
 \overline{\F}_2(\mathcal E)^{\times\ell}.
\]
A lift of \(\mathcal A\) would normalize the cyclic deck group, hence
send \(z\) to \(gz^r\) for
\(r\in(\mathbb Z/\ell\mathbb Z)^\times\), and would require
\(\mathfrak c\circ\mathcal A=g^\ell\mathfrak c^r\), contradicting
this independence.

Finally, set
\[
 \chi(t)=t^{(Q^4-1)/\ell},\qquad
 \tau=\chi(\mathfrak c(P_n)).
\]
Since \(\Frob_Q^2(P_n)=\mathcal A^2(P_n)\),
\((\mathsf{EC}_2)\) gives
\(\mathfrak c(P_n)^{Q^2}=\mathfrak c(P_n)+1\).
Because \(Q^2\equiv-1\pmod\ell\), the four CM/Frobenius character
values are only
\[
 \tau,\qquad\tau^{\pm Q},\qquad\tau^{-1},\qquad
 \tau^{\mp Q}.
\]
All norm identities among them are tautological.  This proves neither
\(\tau=1\) nor \(\tau\ne1\): the desired assertion remains
\((\mathsf{EC}_5)\).
\end{proof}

\begin{proposition}[Miller-unit and generalized-Jacobian location]
\label{prop:fermat-generalized-jacobian}
\PROVED{} Retain Proposition~\ref{prop:fermat-elliptic-kummer}, put
\(S=-2T\), \(q=Q^2\), and \(K_n=\F_{Q^4}=\F_{q^2}\).  Then

\[
 \mathfrak c^{-1}=f_{3,S}
\]

up to a nonzero scalar, where \(f_{3,S}\) is the Miller tripling
function with divisor

\[
 \operatorname{div}(f_{3,S})=3[S]-[3S]-2[O].          \tag{\(\mathsf{GJ}_1\)}
\]

For every prime \(\ell\mid q+1\),

\[
 \#\mathcal E(K_n)=(q-1)^2,\qquad
 \mathcal E(K_n)[\ell]=0,\qquad
 [\ell]:\mathcal E(K_n)\xrightarrow{\sim}\mathcal E(K_n). \tag{\(\mathsf{GJ}_2\)}
\]

Let \(B=\{-T,O,-2T\}\) and \(U=\mathcal E\setminus B\).  The Kummer
class of \(\mathfrak c\) on \(U\) has nonzero tame residue vector

\[
 (1,2,-3)\pmod\ell.                                  \tag{\(\mathsf{GJ}_3\)}
\]

It is therefore ramified and does not come from an unramified
\(\mu_\ell\)-class on \(\mathcal E\).  In the generalized Jacobian

\[
 1\longrightarrow\mathbb G_m^2\longrightarrow J_B
 \longrightarrow\mathcal E\longrightarrow0,          \tag{\(\mathsf{GJ}_4\)}
\]

Hilbert 90 and \((\mathsf{GJ}_2)\) give

\[
 J_B(K_n)/\ell J_B(K_n)
 \cong(K_n^\times/K_n^{\times\ell})^2.                \tag{\(\mathsf{GJ}_5\)}
\]

Thus the residue vector has a genuinely ramified toric projection;
the ordinary \(K_n\)-rational elliptic quotient cannot detect that
projection.  Ordinary \(K_n\)-rational Tate--Lichtenbaum and formal-
group data cannot by themselves decide \((\mathsf{EC}_5)\).  This
does not assert that the ramified class evaluates nontrivially at the
distinguished point, nor does it exclude a calculation using
non-rational \(\mathcal E[\ell]\) or the full extension structure of
\(J_B\).  An elliptic completion must evaluate that
generalized-Jacobian Miller-unit residue.
\end{proposition}

\begin{proof}
The point \(T\) has order five, so \(3S=3(-2T)=-T\).  Equation
\((\mathsf{EC}_1)\) therefore gives

\[
 \operatorname{div}(\mathfrak c^{-1})
 =3[S]-[3S]-2[O],
\]

which proves \((\mathsf{GJ}_1)\).

For \(n\geq4\), the Frobenius relation used in
Proposition~\ref{prop:fermat-elliptic-kummer} gives

\[
 \pi^{2^{n+1}}=[q].
\]

This is Frobenius on \(K_n=\F_{2^{2^{n+1}}}\), so

\[
 \#\mathcal E(K_n)=\deg([q]-[1])=(q-1)^2.
\]

As \(\ell\mid q+1\) is odd, it is coprime to \(q-1\).  This proves
\((\mathsf{GJ}_2)\); multiplication by \(\ell\) is likewise an
automorphism of the characteristic-two formal group because its
linear coefficient is a unit.

The divisor \((\mathsf{EC}_1)\) gives the residue vector
\((\mathsf{GJ}_3)\), which is nonzero because
\(\ell\nmid6\).  The standard generalized-Jacobian sequence for the
three-point reduced rational modulus is \((\mathsf{GJ}_4)\).  Its
torus is split; Hilbert 90 makes the map on \(K_n\)-points surjective,
and the kernel created on passing modulo \(\ell\) vanishes because
\(\mathcal E(K_n)[\ell]=0\).  This proves \((\mathsf{GJ}_5)\) and the
stated scope of the toric obstruction.
\end{proof}

\begin{proposition}[exact extension character and cohomological exhaustion]
\label{prop:fermat-gj-character}
\PROVED{} Retain Proposition~\ref{prop:fermat-generalized-jacobian} and put
\[
 D_1=[-T]-[-2T],\qquad D_2=[O]-[-2T]
\]
in \(\operatorname{Div}^0(B)\).  The character lattice of the split
torus in \(J_B\) is
\[
 X^*(\mathbb G_m^2)=\operatorname{Div}^0(B)
 =\mathbb ZD_1\oplus\mathbb ZD_2,
\]
and the extension is classified, under
\(\mathcal E\cong\operatorname{Pic}^0(\mathcal E)\), by
\[
 \lambda(aD_1+bD_2)=(a+2b)T.                         \tag{\(\mathsf{GJ}_6\)}
\]
Thus its entire extension class is rational five-torsion.  Moreover
\[
 r:=\operatorname{div}(\mathfrak c)=D_1+2D_2,
 \qquad\lambda(r)=5T=O,
\]
so \(r\) defines a character \(\varphi_r:J_B\to\mathbb G_m\).  If
\(j_T:U\to J_B\) is the generalized Abel map normalized by
\(j_T(T)=0\), then
\[
 \varphi_r(j_T(P))=\mathfrak c(P).                    \tag{\(\mathsf{GJ}_7\)}
\]
For every prime \(\ell\mid q+1\), even the non-rational geometric
\(\mathcal E[\ell]\) supplies no further rational Kummer class:
\[
 H^1(K_n,\mu_\ell^2)\xrightarrow{\sim}
 H^1(K_n,J_B[\ell]),
 \qquad
 J_B(K_n)/\ell J_B(K_n)\cong
 (K_n^\times/K_n^{\times\ell})^2.                    \tag{\(\mathsf{GJ}_8\)}
\]
In torus coordinates dual to \((D_1,D_2)\), the character is
\((x,y)\mapsto xy^2\), hence the nonzero linear functional
\((x,y)\mapsto x+2y\) modulo \(\ell\).  Naturality of Kummer theory gives
\[
 \varphi_{r,*}\!\left(\delta_J(j_T(P_n))\right)
   =P_n^*[\mathfrak c]=[\mathfrak c(P_n)]
   \in K_n^\times/K_n^{\times\ell}.                  \tag{\(\mathsf{GJ}_9\)}
\]
Consequently \((\mathsf{EC}_5)\) is exactly the assertion that the
normalized Abel class of \(P_n\) avoids the hyperplane \(x+2y=0\).
Torus translation above any fixed elliptic point realizes every value
of this functional.  Hence the CM point, the full extension class, and
non-rational elliptic \(\ell\)-torsion do not force
\((\mathsf{EC}_5)\); the remaining datum is the explicit toric coordinate
of this one distinguished normalized Abel lift.
\end{proposition}

\begin{proof}
For a reduced three-point rational modulus, the generalized-Jacobian
torus is \((\mathbb G_m^3)/\mathbb G_m\), with character lattice
\(\operatorname{Div}^0(B)\).  The standard extension-class map sends a
boundary divisor to its class in \(\operatorname{Pic}^0(\mathcal E)\).
Here
\[
 (-T)-(-2T)=T,\qquad O-(-2T)=2T,
\]
which proves \((\mathsf{GJ}_6)\).  Equation \((\mathsf{EC}_1)\) gives
\(r=D_1+2D_2\), and its image is \(T+4T=O\).  Characters of \(J_B\)
are precisely the torus characters whose pushout extension class
vanishes, so \(r\) yields \(\varphi_r\).  The generalized-Albanese
universal property gives \((\mathsf{GJ}_7)\) after normalization at
\(T\).  Indeed, \(T\notin B\) and
\(\mathfrak c(T)\in\F_2^\times=\{1\}\).

The \(K_n\)-Frobenius on \(\mathcal E\) is \([q]\), and
\(q\equiv-1\pmod\ell\).  It therefore acts as \([-1]\) on
\(\mathcal E[\ell]\).  Since \(\ell\) is odd,
\[
 \mathcal E[\ell](K_n)=0,
 \qquad
 H^1(K_n,\mathcal E[\ell])
 =\mathcal E[\ell]/(\Frob-1)\mathcal E[\ell]=0.
\]
The cohomology sequence of
\[
 0\longrightarrow\mu_\ell^2\longrightarrow J_B[\ell]
 \longrightarrow\mathcal E[\ell]\longrightarrow0
\]
then proves the first isomorphism in \((\mathsf{GJ}_8)\).  Lang's
theorem and Kummer theory prove the second.  The natural Kummer square
for \(\varphi_r\), together with \((\mathsf{GJ}_7)\), gives
\((\mathsf{GJ}_9)\).  Finally
\(\mathbb G_m^2(K_n)/\ell\to J_B(K_n)/\ell\) is an isomorphism:
surjectivity uses bijectivity of \([\ell]\) on \(\mathcal E(K_n)\),
and injectivity uses \(\mathcal E(K_n)[\ell]=0\).  The character
\((x,y)\mapsto xy^2\) is surjective on the \(\ell\)-quotient, proving
the fixed-fibre statement and the asserted exact boundary.
\end{proof}

\begin{proposition}[explicit Rosenlicht coordinate and its reciprocity boundary]
\label{prop:fermat-rosenlicht}
\PROVED{} Retain Proposition~\ref{prop:fermat-gj-character}, write
\(S=-2T\), \(P=P_n\), and put \(X=P-S\).  For
\(\ell\mid q+1\), let \(Z\in\mathcal E(K_n)\) be the unique point with

\[
 [\ell]Z=P-T,
 \qquad W=T+Z.                                      \tag{\(\mathsf{RC}_1\)}
\]

If \(e\ell\equiv1\pmod{5(Q+1)}\), then

\[
 Z=[e]X+[2e]T,
 \qquad W=[e]X+[2e+1]T,                             \tag{\(\mathsf{RC}_2\)}
\]

so the division point remains an explicit translate of the same cyclic
CM kernel as \(X\).  Choose a Miller function \(F=F_{\ell,Z}\) with

\[
 \operatorname{div}(F)
 =\ell[Z]-[[\ell]Z]-(\ell-1)[O].
\]

In the Rosenlicht coordinates dual to
\(D_1=[-T]-[S]\), \(D_2=[O]-[S]\), the normalized Abel Kummer class of
\(P\) is

\[
 \rho_1=\frac{F(2T)}{F(S)},
 \qquad
 \rho_2=\frac{F(2T)}{F(-T)},                       \tag{\(\mathsf{RC}_3\)}
\]

and its Miller-unit character is exactly

\[
 [\mathfrak c(P)]=[\rho_1\rho_2^2]
 \quad\hbox{in }K_n^\times/K_n^{\times\ell}.       \tag{\(\mathsf{RC}_4\)}
\]

There is also a formula avoiding the \(\ell\)-division point.  On
\(\mathcal E:y^2+y=x^3+x^2\), set

\[
 A_5=1+x+x^2+y+xy,
 \qquad B_5=1+x+xy,
 \qquad g_1=\frac{A_5}{B_5},
 \qquad g_2=\frac1{B_5}.
\]

Then

\[
 \operatorname{div}(g_1)=5D_1,
 \qquad \operatorname{div}(g_2)=5D_2,
 \qquad \mathfrak c^5=g_1g_2^2.                    \tag{\(\mathsf{RC}_5\)}
\]

If \(5\bar5\equiv1\pmod\ell\), the full normalized Abel class in
\(J_B(K_n)/\ell J_B(K_n)\) has coordinates

\[
 \bigl([g_1(P)]^{\bar5},[g_2(P)]^{\bar5}\bigr).
                                                               \tag{\(\mathsf{RC}_6\)}
\]

More explicitly, for the consecutive Conway variables
\(u=u_n\), \(v=u_{n-1}\), put

\[
 U=u+1,\qquad V=v+1,\qquad D=V^2+V+U.
\]

Then the normalized point has \(x=V/D\), \(y=U/D\), and

\[
 g_1(P)=
 \frac{D^2+D(U+V)+V^2+UV}{D^2+VD+UV},
 \qquad
 g_2(P)=\frac{D^2}{D^2+VD+UV},                       \tag{\(\mathsf{RC}_7\)}
\]

with \(g_1(P)g_2(P)^2=U^{-5}\).

These formulas compute the missing toric coordinate exactly, but do not
force it to be nonzero modulo \(\ell\).  Indeed, Weil reciprocity gives

\[
 \mathfrak c(P)
 =\mathfrak c(W)^\ell
   \frac{F(2T)^3}{F(S)F(-T)^2}.                      \tag{\(\mathsf{RC}_8\)}
\]

Applying reciprocity instead to
\(G(Y)=\mathfrak c(Y+T)\), whose divisor is
\([S]+2[-T]-3[2T]\), returns precisely the inverse identity

\[
 \frac{F(S)F(-T)^2}{F(2T)^3}
 =\frac{\mathfrak c(W)^\ell}{\mathfrak c(P)}.
                                                               \tag{\(\mathsf{RC}_9\)}
\]

Thus the boundary Miller values and the rational five-torsion splitting
exhaust the Rosenlicht coordinate, but collapsing them by the evident
reciprocity law is exactly equivalent to the original Kummer value
\([\mathfrak c(P)]\); it supplies no independent nonresidue constraint.
\end{proposition}

\begin{proof}
Since \(P=X+3T\), one has \(P-T=X+2T\).  Multiplication by \(\ell\)
is invertible both on the cyclic kernel containing \(X\) and on
\(\langle T\rangle\), which proves \((\mathsf{RC}_2)\).

Put \(h(Y)=F(Y-T)\).  Translation of the displayed divisor of \(F\)
gives

\[
 \operatorname{div}(h)=\ell[W]-[P]-(\ell-1)[T].
\]

The Rosenlicht evaluation of this principal relation against \(D_1,D_2\)
is

\[
 \frac{h(S)}{h(-T)}=\frac{F(2T)}{F(S)},
 \qquad
 \frac{h(S)}{h(O)}=\frac{F(2T)}{F(-T)},
\]

which proves \((\mathsf{RC}_3)\).  Weil reciprocity between \(h\) and
\(\mathfrak c\), using

\[
 \operatorname{div}(\mathfrak c)=[-T]+2[O]-3[S],
 \qquad \mathfrak c(T)=1,
\]

gives \((\mathsf{RC}_8)\), and hence
\((\mathsf{RC}_4)\).

Direct intersection with the Weierstrass equation, including the local
orders at the indicated five-torsion points, gives

\[
 \operatorname{div}(A_5)=5[-T]-5[O],
 \qquad
 \operatorname{div}(B_5)=5[S]-5[O].
\]

Thus \(g_1,g_2\) have the divisors asserted in
\((\mathsf{RC}_5)\).  Both sides of
\(\mathfrak c^5=g_1g_2^2\) have the same divisor and value one at
\(T\), proving the exact identity.  Since five is invertible modulo
\(\ell\), this proves \((\mathsf{RC}_6)\).  Substitution of
\(x=V/D\), \(y=U/D\) gives \((\mathsf{RC}_7)\); here
\(\mathfrak c(P)=1/U\), so its last identity is also immediate from
\((\mathsf{RC}_5)\).

Finally, the divisor of \(G\) follows by translating that of
\(\mathfrak c\).  Weil reciprocity between \(F\) and \(G\) is
\((\mathsf{RC}_9)\), proving the exact neutrality statement.
\end{proof}

When $F_n$ is prime, irreducibility and $\delta_n>1$ force
$\delta_n=F_n$, recovering Lenstra's Fermat-prime rule.  When $F_n$ is
composite, irreducibility forces only one nontrivial divisor: it does not
force every prime-power factor.  The condition $\delta_n=F_n$ is exactly
the maximal-coset-order hypothesis studied for Conway's quadratic tower
by Popovych~\cite{Popovych18}.  It is strictly weaker than primitivity of
$c_n$, which already fails at $n=2$ since
$\ord(c_2)=5\cdot17=85<255$; Popovych uses the all-level
hypothesis instead to prove that products such as $c_0c_n$ are primitive.
Wiedemann's related iterated quadratic tower
poses the analogous full-Fermat-order problem~\cite{Wiedemann88}.  Thus
the apparently easy ``singleton even implies zero'' clause already
contains a longstanding explicit maximal-coset-order problem.
Popovych verified the Conway condition only through $n=11$; the recent
analysis of Cagliero--Herman--Szechtman still retains the quotient order
as an undetermined divisor~\cite{CHS25}.  No all-level proof or
counterexample is presently known.

\begin{proposition}[the $3$-power arm is Gaussian-period primitivity]
\label{prop:gauss-arm}
\PROVED{} Fix $k\geq1$.  The $m=1$ arm of
Conjecture~\ref{conj:rule} for every prime $p$ with
$f(p)\in\{3,9,\ldots,3^k\}$ is equivalent to $C_k$: the Gaussian period
\[
  \gamma_k=\zeta+\zeta^{-1}\in\F_{2^{3^k}}^{\times},
  \qquad \zeta^{3^{k+1}}=1,
\]
is primitive.
\end{proposition}

\begin{proof}
Propositions~\ref{prop:halfangle} and~\ref{prop:tower} show that the old
prime-power factors of $2^{3^k}-1$ propagate from lower levels.  For a
prime $p$ with $f(p)=3^j$, the full power criterion says $m_p=1$ iff
$v_p(\ord\gamma_j)=v_p(2^{3^j}-1)$.  Combining these statements for all
$j\leq k$ gives every prime-power factor of $2^{3^k}-1$, which is
equivalent to primitivity; the converse is immediate.
\end{proof}

The two arms are instances of the same projective statement.  The
quadratic arm asks that the class of $c_n$ generate
$E_n^\times/E_{n-1}^\times$.  In the cubic arm the natural map
\[
 U_k\longrightarrow \F_{q_k^3}^{\times}/\F_{q_k}^{\times}
\]
is an isomorphism, and the increment from $C_{k-1}$ to $C_k$ asks that
the class of $\gamma_k$
generate this quotient.  Equivalently, the projective class of the
multiplication-by-$\gamma_k$ matrix generates a Singer cycle in
$\operatorname{PGL}_3(q_k)$.  The $0/1/4$ rule therefore contains a
quadratic and a cubic maximal-projective-order conjecture, rather than
two unrelated order coincidences.

\begin{proposition}[cyclotomic-type residue reduction]
\label{prop:unit-residue}
\PROVED{} Let $h=3^k$ and let $\ell$ be any prime divisor of
$2^h-1$.  Then
\[
 \gamma_k=\zeta^{-1}(1+\zeta)^2
\]
is an $\ell$-th power in $\F_{2^h}$ if and only if $1+\zeta$ is an
$\ell$-th power in $\F_{2^{2h}}$.  The same equivalence holds at every
power $\ell^a\mid2^h-1$; however, non-$\ell$-power already implies that
$\gamma_k$ has the full $\ell$-primary part of $2^h-1$ in its order.
\end{proposition}

\begin{proof}
Because $\ell\mid2^h-1$, one has
$\ell\nmid2^h+1$.  Extension from $\F_{2^h}^{\times}$ to
$\F_{2^{2h}}^{\times}$ therefore preserves the $\ell$-primary quotient.
Also $\ell\neq3$, so exponentiation by $\ell$ is an automorphism of the
$3^{k+1}$-group generated by $\zeta$; hence $\zeta^{-1}$ is an
$\ell$-th power.  Finally, squaring is an automorphism of every
$\ell$-primary quotient for odd $\ell$.  The displayed factorization
then proves both assertions.
\end{proof}

Even in the standard prime-conductor case, general results on
multiplicative orders of Gaussian periods provide large lower bounds but
not maximal order~\cite{ASV10}.  For the present prime-power conductor,
the norm recursion $\gamma_k^3+\gamma_k=\gamma_{k-1}$ propagates the old
factors only.  At level $k$, a proof must still show that every new
prime-power factor of $\Phi_{3^k}(2)$ divides $\ord(\gamma_k)$.  That is
the precise new-prime obstruction; degree $3^k$ guarantees at most that
some exact-order factor occurs, not that all of them do.

The Dickson formulation cleanly separates what the existing theory does
and does not supply.  Mullin--Mahalanobis give the corresponding
irreducibility criterion for $D_{3^k}(X,1)+1$~\cite{MM05}; irreducibility
is not primitivity.  For prime conductor, Gao--Vanstone's analogous
Gaussian-period primitivity statement remains conjectural~\cite{GV95},
while Breuer shows that small divisors of the period index can reflect
units and class numbers in an intermediate real quadratic
field~\cite{Breuer21}.  The real $3$-power cyclotomic tower here has only
odd-degree intermediate fields, so that quadratic norm argument does not
apply.  Proposition~\ref{prop:eta-tower} supplies the appropriate cubic
replacement and places the root on a Singer difference set, but neither
the three symmetric coefficients nor the difference-set property proves
that this distinguished member generates $U_k$; the explicit
$q=32$ example in Remark~\ref{rem:singer-not-primitive} rules out that
shortcut.

\subsection{The component set alone is insufficient}

\PROVED{} $\Qset(f(p))$ alone does not determine $m_p$:
\begin{center}
\begin{tabular}{ccc}
\toprule
$\Qset$ & $m=4$ & $m=1$ \\
\midrule
$\{9\}$   & $p=19$,   $f=18$  & $p=73$,   $f=9$ \\
$\{81\}$  & $p=163$,  $f=162$ & $p=2593$, $f=81$ \\
$\{243\}$ & $p=1459$, $f=486$ & $p=487$,  $f=243$ \\
\bottomrule
\end{tabular}
\end{center}

\PROVED{} The split for $\Qset=\{9\}$ is explained by the order
criterion: $\ord(\kappa_9+1)=3^3\cdot(2^9-1)=27\cdot511=13797$ and
$511=7\cdot73$, so $73\mid\ord(\kappa_9+1)$ but
$19\nmid\ord(\kappa_9+1)$; adding $m=4$ changes the order and introduces
the factor $19$.

\PROVED{} The lower bound of the exception column is unconditional:
$m_p\geq4$ for all $p$ with $f(p)=2\cdot3^k$
(Proposition~\ref{prop:exception}).

\CONJECTURED{} The matching upper bound $m_p\leq4$ on that column;
its exact current-level formulation is $D'_k$ in
Conjecture~\ref{conj:dk}, not the Gaussian-period conjecture $C_k$.

\subsection{The $m=4$ upper bound: settled at every visible prime}
\label{sec:m4}

The $m=4$ translate $\kappa_{3^k}+4$ involves the nimber $4\in\F_{16}$,
which does not lie in $\F_2(\zeta)$ (as $4\nmid2\cdot3^k$); the
translate lives in the degree-$4\cdot3^k$ compositum $\F_{2^{4h}}$,
$h=3^k$. (June 2026 pass; probe
\path{experiments/exception_column_m4.py}.)

\begin{proposition}[the corrected norm]\label{prop:m4norm}
\PROVED{} Let $\sigma=\Frob^{2h}$ generate
$\mathrm{Gal}(\F_{2^{4h}}/\F_{2^{2h}})$. Then $\sigma(4)=5$, and with
$\omega:=\zeta^{h}=\kappa_2$,
\[
  N:=\Norm_{\F_{2^{4h}}/\F_{2^{2h}}}(\kappa_{3^k}+4)
   =(\kappa+4)(\kappa+5)
   =\kappa^2+\kappa+\omega .
\]
\end{proposition}

\begin{proof}
$\F_{16}\cap\F_{2^{2h}}=\F_4$ because $\gcd(4,2\cdot3^k)=2$, so
$\sigma|_{\F_{16}}$ is the nontrivial element of
$\mathrm{Gal}(\F_{16}/\F_4)$, which swaps the two roots $4,5$ of the
Artin--Schreier minimal polynomial $y^2+y+\omega$ of $4$ over $\F_4$.
(Equivalently: $2h\equiv2\pmod4$ and the Frobenius orbit of $4$ is
$4\to6\to5\to7$, so $\sigma(4)=4^{2^2}=5$.) Then
$(\kappa+4)(\kappa+5)=\kappa^2+\kappa+(4\nimmul5)$ and
$4\nimmul5=4^2+4=\omega$ by the same Artin--Schreier relation.
\end{proof}

\begin{remark}[erratum]\label{rem:erratum}
The 2026-06-10 draft of this note stated the norm as
$(\kappa+4)(\kappa+\Frob^2(4))=(\kappa+4)(\kappa+6)$. The conjugate
$6=4^2$ is $\Frob^1(4)$, one squaring short of $\Frob^2(4)=5$. The slip
is visible without computation: $4\nimmul6=14\notin\F_4$, so
$(\kappa+4)(\kappa+6)\notin\F_{2^{2h}}$ and cannot be a norm.
\end{remark}

The corrected norm collapses the $m=4$ test into the \emph{same} sparse
trinomial model $\F_2[x]/(x^{2h}+x^h+1)$ as the rest of the $3$-power
analysis: no compositum arithmetic is needed.

\begin{corollary}[in-field criterion]\label{cor:m4crit}
\PROVED{} For a prime $p$ with $f(p)=2\cdot3^k$ and $v_p(2^{2h}-1)=1$,
write $\overline{N}=N^{2^h}$ and $M_k:=\overline{N}/N=N^{2^h-1}\in U$.
Then
\[
  m_p=4
  \iff
  p\mid\ord(M_k)
  \iff
  M_k^{(2^h+1)/p}\neq1 .
\]
\end{corollary}

\begin{proof}
$m_p\geq4$ is Proposition~\ref{prop:exception}, so $m_p=4$ iff
$\kappa+4$ has no $p$-th root. By Propositions~\ref{prop:power}
and~\ref{prop:norm} with $E=4h$, $f=2h$ (and
$v_p(2^{4h}-1)=v_p(2^{2h}-1)=1$ by lifting the exponent, since
$p\nmid2$), that holds iff $N^{(2^{2h}-1)/p}\neq1$. Since
$\F_{2^{2h}}^{\times}=L^{\times}\times U$ with coprime orders,
$p\mid2^h+1$, and $(2^{2h}-1)/p=(2^h-1)\cdot\bigl((2^h+1)/p\bigr)$, the
test equals $M_k^{(2^h+1)/p}\neq1$.
\end{proof}

\begin{proposition}[exact current-level form of the exception conjecture]
\label{prop:dk-current-level}
\PROVED{} Let $k\geq1$.  Retain the notation above, put
$q=2^h$, $n=3^{k+1}=3h$,
and write
\[
 \zeta:=\kappa_{3^k}\in K:=\F_{2^{2h}},\qquad
 \widetilde K:=\F_{2^{4h}},\qquad \beta:=\zeta+4.
\]
Thus $\zeta$ is a primitive $n$-th root of unity.  Define
\[
 N:=\Norm_{\widetilde K/K}(\beta)
   =\zeta^2+\zeta+\zeta^h,
 \qquad M_k:=N^{q-1},
 \qquad \Psi_k:=\frac{\Phi_{2\cdot3^k}(2)}3,
\]
and
\[
 R_k(Y):=\operatorname{Res}_X
 \left(X^{2h}+X^h+1,\,Y+X^h+X^2+X\right)\in\F_2[Y].
\]
Then $R_k$ is the irreducible minimal polynomial of $N$ over $\F_2$
and has degree $2h$.  The prime divisors of $\Psi_k$ are exactly the
primes $\ell\ne3$ with $\ord_\ell(2)=2\cdot3^k$, with full valuations
\[
 v_\ell(\Psi_k)=v_\ell(2^{2h}-1)=v_\ell(2^h+1).
\]
Moreover the following are equivalent:
\begin{enumerate}[label=\textup{(\alph*)}]
\item $m_\ell=4$ for every prime $\ell$ with
      $f(\ell)=2\cdot3^k$;
\item $\Psi_k\mid\ord(M_k)$;
\item $N\notin(K^\times)^\ell$ for every prime $\ell\mid\Psi_k$;
\item $R_k(X^\ell)$ is irreducible over $\F_2$ for every prime
      $\ell\mid\Psi_k$.
\end{enumerate}
\end{proposition}

\begin{proof}
The order identity $\ord_{3^{k+1}}(2)=2h$ makes
$X^{2h}+X^h+1$ the irreducible minimal polynomial of $\zeta$.
We first prove that $N=\zeta^h+\zeta^2+\zeta$ also has degree $2h$.
Suppose $N^{2^j}=N$, and put $a\equiv2^j\pmod{3h}$.  The resulting
vanishing polynomial, of degree less than $3h$, has support equal to
the symmetric difference of
\[
 \{ah,a,2a\}\quad\hbox{and}\quad\{h,1,2\}.                \tag{1}
\]
Every multiple of $X^{2h}+X^h+1$ of degree less than $3h$ has support
a disjoint union of triples $\{r,r+h,r+2h\}$.
If $a\equiv2\pmod3$, the support in (1) contains $h$ and $2h$ but not
$0$, a contradiction.  If $a\equiv1\pmod3$, the $h$-terms cancel and
each residue class modulo $h$ occurs at most twice; a nonempty union of
such triples is again impossible.  Hence the support is empty, so
$\{a,2a\}=\{1,2\}$ modulo $3h$.  Since $h\geq3$, this forces $a=1$.
Thus the Frobenius orbit of $N$ has length $2h$, and the displayed
resultant is its irreducible minimal polynomial.

Writing $A=2^{3^{k-1}}$ gives
$\Phi_{2\cdot3^k}(2)=A^2-A+1\equiv3\pmod9$, so division by $3$
removes exactly the prime dividing the cyclotomic index.  The standard
cyclotomic-divisor criterion now says that the prime divisors of
$\Psi_k$ are precisely the $\ell$ with $\ord_\ell(2)=2h$.  Their
entire primary contribution to $2^{2h}-1$ lies in this cyclotomic
factor and, since $\ell\nmid2^h-1$, in $2^h+1$.  This proves the
valuation statement.

Fix $\ell\mid\Psi_k$.  Propositions~\ref{prop:translates}
and~\ref{prop:exception} make the translates
$\zeta,\zeta+1,\zeta+2,\zeta+3$ into $\ell$-th powers, so
$m_\ell\geq4$.  The corrected norm and the exact power criterion give
\[
 \beta^{(2^{4h}-1)/\ell}
 =N^{(2^{2h}-1)/\ell},
\]
and therefore
\[
 m_\ell=4\iff N\notin(K^\times)^\ell.                    \tag{2}
\]
This uses the full $\ell$-primary valuation and has no squarefreeness
or non-Wieferich hypothesis.  Since $q-1\equiv-2\pmod\ell$, raising
to $q-1$ is invertible on $K^\times/(K^\times)^\ell$, so
\[
 N\in(K^\times)^\ell\iff M_k\in(K^\times)^\ell.           \tag{3}
\]
The entire $\ell$-primary part of $K^\times$ lies in its norm-one
factor of order $q+1$.  In a cyclic group, nonmembership in the
subgroup of $\ell$-th powers is equivalent to the order containing the
full $\ell$-primary part.  Equations (2)--(3) prove the equivalence of
the first three assertions.

Finally $\mu_\ell\subset K$.  Kummer theory says that
$X^\ell-N$ is irreducible over $K$ exactly when
$N\notin(K^\times)^\ell$.  Since $K=\F_2(N)$, Capelli's lemma says in
turn that this is equivalent to irreducibility of $R_k(X^\ell)$ over
$\F_2$, completing the proof.
\end{proof}

\begin{conjecture}[$D'_k$: the exact current exception factor]
\label{conj:dk}
\CONJECTURED{} One has
\[
 D'_k:\qquad \Psi_k\mid\ord(M_k).
\]
By Proposition~\ref{prop:dk-current-level}, this is exactly the
$m=4$ upper bound for the primes at level $f=2\cdot3^k$.
\end{conjecture}

\begin{proposition}[exact current-layer projection and dominant-factor sieve]
\label{prop:dk-projection}
\PROVED{} Let \(k\geq2\), put
\[
 s=3^{k-1},\qquad h=3s,\qquad a=2^s,\qquad q=a^3=2^h,
 \qquad \Lambda_k=a^2-a+1=3\Psi_k,
\]
and retain \(N,M_k,\zeta,\omega=\zeta^h\) from
Proposition~\ref{prop:dk-current-level}.  Define
\[
 P_k:=M_k^{a+1}.
\]
Then \(P_k^{\Lambda_k}=1\), and for every prime
\(\ell\mid\Psi_k\),
\[
 v_\ell(\ord P_k)=v_\ell(\ord M_k).                    \tag{\(\mathsf D_1\)}
\]
Moreover
\[
 P_k=\frac{\zeta^4+\zeta+\omega}
           {\omega^2\zeta^4+\zeta^3+1}.               \tag{\(\mathsf D_2\)}
\]
\[
 [\F_2(P_k):\F_2]=2h.                                 \tag{\(\mathsf D_3\)}
\]
and
\[
 v_3(\ord P_k)=
 \begin{cases}1,&k\text{ odd},\\0,&k\text{ even}.
 \end{cases}                                           \tag{\(\mathsf D_4\)}
\]
Set \(c_k=3\) for odd \(k\) and \(c_k=1\) for even \(k\).  Then
\[
 \ord(P_k)\geq c_k(2h+1).                              \tag{\(\mathsf D_5\)}
\]
Consequently, if \(\ell\mid\Psi_k\) and
\[
 \ell>\frac{\Lambda_k}{c_k(2h+1)},                    \tag{\(\mathsf D_6\)}
\]
then \(P_k\), and hence \(M_k\), contains the full
\(\ell\)-primary part of \(\Lambda_k\) in its order.  In particular
\(D'_k\) follows whenever \(\Psi_k\) is prime; this proves \(D'_2\)
and \(D'_3\) analytically, since \(\Psi_2=19\) and
\(\Psi_3=87211\).
\end{proposition}

\begin{proof}
The factorization
\[
 q+1=(a+1)(a^2-a+1)=(a+1)\Lambda_k
\]
and \(M_k^{q+1}=1\) give \(P_k^{\Lambda_k}=1\).  Also
\(\gcd(\Lambda_k,a+1)=3\).  Since \(k\geq2\), one has
\(a\equiv-1\pmod9\), so
\(\Lambda_k\equiv3\pmod9\) and
\(\gcd(\Psi_k,a+1)=1\).  Exponentiation by \(a+1\) therefore preserves
all the primary valuations belonging to \(\Psi_k\), proving
\((\mathsf D_1)\).

The lifted congruence
\[
 2^{3^{k-1}}\equiv3^k-1\pmod{3^{k+1}}
\]
gives, modulo \(3h\),
\[
 a\equiv h-1,\quad a^2\equiv h+1,\quad
 a^3\equiv-1,\quad a^4\equiv2h+1.                     \tag{\(\mathsf D_7\)}
\]
As \(M_k=N^{a^3-1}\),
\[
 P_k=\frac{N^{a^4}N^{a^3}}{N^aN}.
\]
Using \(N=\zeta^2+\zeta+\omega\) and \((\mathsf D_7)\), direct Frobenius
expansion gives
\begin{align*}
 N^a&=\omega^2\zeta^{-2}+\omega\zeta^{-1}+\omega^2,\\
 N^{a^3}&=\zeta^{-2}+\zeta^{-1}+\omega^2,\\
 N^{a^4}&=\omega\zeta^2+\omega^2\zeta+\omega.
\end{align*}
Multiplying \(N^{a^3}\) and \(N^a\) by the same cancelling factor
\(\zeta^2\), and using \(1+\omega+\omega^2=0\), yields
\begin{align*}
 N^{a^4}(\zeta^2N^{a^3})&=\zeta^4+\zeta+\omega,\\
 (\zeta^2N^a)N&=\omega^2\zeta^4+\zeta^3+1,
\end{align*}
which proves \((\mathsf D_2)\).

The minimal polynomial of \(\zeta\) over \(\F_2\) is
\(X^{2h}+X^h+1\).  Write the numerator and denominator in
\((\mathsf D_2)\) as
\(C,D\).  Since \(P_k^q=P_k^{-1}\), membership in \(\F_q\) would force
\(P_k=1\).  But \(C=D\), after reducing
\(\zeta^{2h}=\zeta^h+1\), would give
\[
 \zeta^{h+4}+\zeta^h+\zeta^3+\zeta+1=0,               \tag{\(\mathsf D_8\)}
\]
a nonzero polynomial of degree \(h+4<2h\).

If \(P_k\in\F_{a^2}\), its order divides
\(\gcd(\Lambda_k,a^2-1)=3\), so
\(P_k\in\{1,\omega,\omega^2\}\).  The first case was excluded; the
other two respectively give
\[
 \zeta^{h+3}+\zeta=0
\]
and
\[
 \zeta^{h+4}+\zeta^{h+3}+\zeta^4+\zeta^3+\zeta+1=0.
\]
Both are again nonzero polynomials of degree below \(2h\).  Every proper
divisor of \(2h=2\cdot3^k\) divides either \(h\) or \(2h/3\); hence every
proper subfield of \(\F_{2^{2h}}\) lies in either \(\F_q\) or
\(\F_{a^2}\).  This proves \((\mathsf D_3)\).

For the \(3\)-part, observe that
\[
 P_k^{\Lambda_k/3}
 =M_k^{(q+1)/3}
 =N^{(q^2-1)/3}
 =\Norm_{\F_{2^{2h}}/\F_4}(N).                         \tag{\(\mathsf D_9\)}
\]
Choose \(u\in\F_{16}\) with \(u^2+u=\omega\).  Then
\(u^4=u+1\), \(u^5=\omega\), and \(u\) has order \(15\).  The
extensions \(\F_{2^{2h}}/\F_4\) and \(\F_{16}/\F_4\) are linearly
disjoint, and over \(\F_{16}\) the minimal polynomial of \(\zeta\)
is \(X^h+\omega\), irreducible because
\(\ord_{3h}(16)=h\).  Since
\[
 N=(\zeta+u)(\zeta+u+1),
\]
base change in \((\mathsf D_9)\) gives
\[
 \Norm_{\F_{2^{2h}}/\F_4}(N)
 =(u^h+\omega)((u+1)^h+\omega).
\]
The residues of \(h=3^k\) modulo \(15\) are \(3,9,12,6\) as
\(k\equiv1,2,3,0\pmod4\); reduction with \(u^4=u+1\) gives
\[
 P_k^{\Lambda_k/3}=
 \begin{cases}\omega^2,&k\text{ odd},\\1,&k\text{ even}.
 \end{cases}
\]
Because \(v_3(\Lambda_k)=1\), this proves \((\mathsf D_4)\).

Let \(e=\ord(P_k)\).  From \((\mathsf D_3)\),
\(\ord_e(2)=2h\).  Since
\(e\mid\Lambda_k\) and \(e\notin\{1,3\}\), some prime
\(r\mid\Psi_k\) divides \(e\); every such \(r\) has
\(\ord_r(2)=2h\), so \(r\equiv1\pmod{2h}\) and
\(r\geq2h+1\).  When \(k\) is odd, \((\mathsf D_4)\) also gives
\(3r\mid e\), which proves \((\mathsf D_5)\).

If the full \(\ell\)-primary part of \(\Lambda_k\) were missing from
\(e\), then \(e\mid\Lambda_k/\ell\).  This contradicts
\((\mathsf D_5)\) under \((\mathsf D_6)\).
Equation \((\mathsf D_1)\) transfers the conclusion back to \(M_k\), and the
last assertions follow.
\end{proof}

\begin{proposition}[fibotomic form of the current power-residue test]
\label{prop:dk-fibotomic}
\PROVED{} Retain Proposition~\ref{prop:dk-projection}, put
\[
 U=\{x\in\F_{q^2}^{\times}:x^{q+1}=1\},\qquad
 \pi(x)=\frac{x}{(x+1)^2}=\frac1{x+x^{-1}},
\]
and set \(Z=\pi(M_k)\).  The map \(\pi\) is two-to-one from
\(U\setminus\{1\}\) onto
\[
 \mathcal T_q=\{z\in\F_q:
   \operatorname{Tr}_{\F_q/\F_2}(z)=1\},
\]
with fibres \(\{x,x^{-1}\}\).

For an odd prime \(\ell\), define
\[
 H_\ell(W)=\sum_{j=0}^{(\ell-1)/2}
   \binom{\ell-j-1}{j}W^j\in\F_2[W].
\]
If \(\ell\mid\Psi_k\), then
\[
 M_k\in U^\ell
 \iff
 \text{there is }W\in\mathcal T_q\text{ with }
 W^\ell+Z H_\ell(W)^2=0.                               \tag{\(\mathsf F\)}
\]
The polynomial
\[
 \mathcal F_{\ell,Z}(W)=W^\ell+ZH_\ell(W)^2
\]
either has no root in \(\F_q\), or splits there into \(\ell\)
distinct linear factors.  Finally \(H_\ell\) is squarefree and is a
product of exactly \((\ell-1)/(2h)\) irreducible polynomials over
\(\F_2\), each of degree \(h\).  Thus \(D'_k\) is equivalently the
assertion that \(\mathcal F_{\ell,Z}\) is rootless over \(\F_q\) for
every prime \(\ell\mid\Psi_k\).
\end{proposition}

\begin{proof}
For \(x\in U\), one has \(x^q=x^{-1}\), so \(\pi(x)\in\F_q\).
The equation \(\pi(x)=z\) is
\[
 X^2+z^{-1}X+1=0.
\]
The characteristic-two quadratic criterion says that this polynomial
is irreducible over \(\F_q\) exactly when
\(\operatorname{Tr}(z)=1\).  Its roots are inverse and have norm one.
Counting, or the same quadratic argument in reverse, proves the stated
two-to-one description.

The defining polynomial identity is
\[
 \frac{X^\ell+1}{(X+1)^\ell}
 =H_\ell\!\left(\frac{X}{(X+1)^2}\right).              \tag{\(\mathsf H\)}
\]
If \(M_k=T^\ell\) with \(T\in U\), put \(W=\pi(T)\).  Identity
\((\mathsf H)\) gives
\[
 Z=\pi(T^\ell)=\frac{W^\ell}{H_\ell(W)^2},
\]
so \(W\) satisfies \((\mathsf F)\).  Conversely, lift a trace-one
solution \(W\) to \(T\in U\).  Then \((\mathsf H)\) gives
\(\pi(T^\ell)=\pi(M_k)\), hence \(T^\ell=M_k\) or \(M_k^{-1}\).
Replacing \(T\) by \(T^{-1}\) in the second case proves
\(M_k\in U^\ell\).

In fact every \(\F_q\)-root of \(\mathcal F_{\ell,Z}\) has trace one.
Indeed, \(H_\ell(W)\ne0\) at a root, since otherwise the equation
would force \(W=0\), whereas \(H_\ell(0)=1\).  Lift a nonzero root
through \(\pi(T)=W\).  If \(\operatorname{Tr}(W)=0\), the quadratic
criterion puts \(T\) in \(\F_q\).  Identity \((\mathsf H)\) then gives
\(\pi(T^\ell)=Z\), but for \(y\in\F_q\setminus\{1\}\),
\[
 \pi(y)=\left(\frac{y}{y+1}\right)^2+\frac{y}{y+1}
\]
has absolute trace zero, whereas \(Z=\pi(M_k)\) has trace one.  This
is a contradiction.

If one root exists, therefore lift it to \(T\in U\) and orient \(T\)
so that \(T^\ell=M_k\).  Then the
\(\ell\) elements
\[
 \pi(T\xi),\qquad \xi\in\mu_\ell,
\]
are roots in \(\F_q\).  They are distinct: an inverse collision would
give \(T^2\in\mu_\ell\), hence \(M_k^2=1\), whereas \(M_k\ne1\)
because \(N\notin\F_q\).  Since \(\mathcal F_{\ell,Z}\) has degree
\(\ell\), it splits completely.  Its derivative is \(W^{\ell-1}\),
and trace-one roots are nonzero, so they are simple.

Finally, the roots of \(H_\ell\) are
\[
 \pi(\xi),\qquad \xi\in\mu_\ell\setminus\{1\},
\]
with \(\xi\) and \(\xi^{-1}\) identified.  Their Frobenius-orbit
length is the least \(d>0\) with \(2^d\equiv\pm1\pmod\ell\).  Since
\(\ord_\ell(2)=2h\), this least \(d\) is \(h\).  The number and degrees
of the factors follow, and the roots are distinct, proving
squarefreeness.
\end{proof}

\begin{proposition}[translated-binomial form of $D'_k$]
\label{prop:dk-translated-binomial}
\PROVED{} Retain \(k\geq2\), \(h=3^k\), and choose
\(u\in\F_{16}\) with \(u^2+u=\omega\).  Put

\[
 L=\F_{16^h}=\F_{2^{4h}},\qquad \beta=\zeta+u.
\]

Then the minimal polynomial of \(\beta\) over \(\F_{16}\) is the
irreducible degree-\(h\) polynomial

\[
 B_k(Y)=(Y+u)^h+\omega.
\]

For every prime \(\ell\mid\Psi_k\), the following are equivalent:
\begin{enumerate}[label=\textup{(\roman*)}]
\item the full \(\ell\)-primary part occurs in \(\ord(M_k)\);
\item \(\beta\notin(L^\times)^\ell\);
\item \(B_k(X^\ell)=(X^\ell+u)^h+\omega\) is irreducible over
      \(\F_{16}\).
\end{enumerate}
More precisely, \(B_k(X^\ell)\) is either irreducible of degree
\(h\ell\), or is a product of exactly \(\ell\) irreducibles of degree
\(h\).
\end{proposition}

\begin{proof}
Since

\[
 \ord_{3h}(16)
 =\frac{\ord_{3h}(2)}{\gcd(\ord_{3h}(2),4)}=h,
\]

the element \(\zeta\) has degree \(h\) over \(\F_{16}\).  As
\(\zeta^h=\omega\), its minimal polynomial there is
\(X^h+\omega\); translation by \(u\) gives \(B_k\).

Let \(K=\F_{2^{2h}}\).  The corrected norm is

\[
 N=\Norm_{L/K}(\beta)=\beta^{2^{2h}+1}.
\]

Because \(2^{2h}\equiv1\pmod\ell\), this norm induces multiplication
by two on \(L^\times/(L^\times)^\ell\).  Hence it first gives

\[
 \beta\in(L^\times)^\ell
 \iff N\in(L^\times)^\ell.
\]

The natural map

\[
 K^\times/(K^\times)^\ell
 \longrightarrow L^\times/(L^\times)^\ell
\]

is injective: restriction followed by the degree-two norm is
multiplication by two, which is invertible modulo the odd prime
\(\ell\).  Since \(N\in K\), this sharpens the equivalence to

\[
 \beta\in(L^\times)^\ell
 \iff N\in(K^\times)^\ell,
\]

and Proposition~\ref{prop:dk-current-level} identifies the negation
with the full \(\ell\)-primary assertion.  Kummer theory and Capelli's
lemma give the irreducibility equivalence.

If \(\beta=\alpha^\ell\) with \(\alpha\in L\), every root of
\(B_k(X^\ell)\) lies in \(L\), because \(\mu_\ell\subset L\).  Its
\(\ell\)-th power is a conjugate of \(\beta\), of degree \(h\) over
\(\F_{16}\); the root therefore has degree both at least and at most
\(h\).  All irreducible factors consequently have degree \(h\), and
there are exactly \(\ell\) of them.
\end{proof}

\begin{corollary}[short translated cyclotomic period]
\label{cor:dk-short-period}
\PROVED{} Let \(\chi:L^\times\to\mu_\ell\) have exact order
\(\ell\), realize its values as complex roots of unity, and write
\(\chi(\beta)=\xi_\ell^r\).  Then

\[
 \mathcal J_{k,\ell}
 :=\sum_{\theta^h=\omega}\chi(\theta+u)
 =\sum_{j=0}^{h-1}\xi_\ell^{\,r16^j},
\]

and

\[
 \beta\in(L^\times)^\ell
 \iff \mathcal J_{k,\ell}=h.
\]
\end{corollary}

\begin{proof}
The roots of \(X^h+\omega\) are \(\zeta^{16^j}\), \(0\leq j<h\),
and \(u\in\F_{16}\).  Thus

\[
 \chi(\zeta^{16^j}+u)
 =\chi((\zeta+u)^{16^j})
 =\xi_\ell^{\,r16^j}.
\]

If \(r=0\), every summand is one.  Conversely, equality with the
positive real number \(h\) forces equality in the complex triangle
inequality and hence forces every summand, in particular the first,
to be one.
\end{proof}

\begin{proposition}[full \(\F_{16}\)-line saturation]
\label{prop:dk-line-saturation}
\PROVED{} Retain Proposition~\ref{prop:dk-translated-binomial}, put
\(s=3^{k-1}\), and choose \(u\) so that it has order fifteen and
\(\omega=u^5\).  For a prime \(\ell\mid\Psi_k\), write

\[
 \mathcal C=L^\times/(L^\times)^\ell,
 \qquad v(c)=[\zeta+c]\quad(c\in\F_{16}),
 \qquad A=2^s\pmod\ell,
\]

using additive notation for \(\mathcal C\).  Then \(A\) has order six,
\(A^2-A+1=0\), and \(A^3=-1\).  The four base-field classes vanish,

\[
 v(c)=0\qquad(c\in\F_4),                              \tag{\(\mathsf{L}_{16,1}\)}
\]

while for \(c\in\F_{16}^\times\),

\[
 v(\omega c^{-2^s})=A v(c),\qquad
 v(c^{-2^h})=-v(c).                                  \tag{\(\mathsf{L}_{16,2}\)}
\]

For the twelve exponents not divisible by five one has

\[
\begin{array}{c|c}
 j & \{m:v(u^m)=A^jv(u)\}\\ \hline
 0&\{1,4\}\\
 1&\{3,12\}\\
 2&\{11,14\}\\
 3&\{7,13\}\\
 4&\{6,9\}\\
 5&\{2,8\}.
\end{array}                                             \tag{\(\mathsf{L}_{16,3}\)}
\]

Consequently

\[
 \zeta+u\in(L^\times)^\ell
 \iff \zeta+c\in(L^\times)^\ell
      \quad\hbox{for every }c\in\F_{16}.              \tag{\(\mathsf{L}_{16,4}\)}
\]

If \(D'_k\) holds for \(\ell\), none of the twelve translates with
\(c\notin\F_4\) is an \(\ell\)-th power.  Equivalently, for an
exact-order-\(\ell\) complex character \(\chi\) of \(L^\times\) and
\(t=\chi(\zeta+u)\),

\[
 \sum_{c\in\F_{16}}\chi(\zeta+c)
 =4+2\sum_{j=0}^5t^{A^j},                              \tag{\(\mathsf{L}_{16}\)}
\]

and failure of \(D'_k\) is exactly the extremal value sixteen.
Thus Frobenius, conjugation, and all \(\F_4\)-translates generate only
scalar multiples of the one unresolved Kummer class \(v(u)\).
\end{proposition}

\begin{proof}
One has \(\ord_\ell(2)=2h=6s\), proving the assertions about \(A\).
For \(c=0\), the element \(\zeta\) has \(\ell\)-coprime order \(3h\).
For \(c\in\{1,\omega,\omega^2\}\),
Propositions~\ref{prop:translates} and~\ref{prop:exception} give order
dividing \(3h(2^h-1)\), again coprime to \(\ell\).  This proves
\((\mathsf{L}_{16,1})\); constants in \(\F_{16}^\times\) and
\(\zeta\) also have zero class.

The congruence \(2^s\equiv h-1\pmod{3h}\) gives
\(\zeta^{2^s}=\omega\zeta^{-1}\).  Discarding scalar factors with
zero class,

\[
 A v(c)=[(\zeta+c)^{2^s}]
 =[\omega\zeta^{-1}+c^{2^s}]
 =[\zeta+\omega c^{-2^s}]
 =v(\omega c^{-2^s}).
\]

Likewise \(2^h\equiv-1\pmod{3h}\) and modulo \(\ell\), so

\[
 -v(c)=[(\zeta+c)^{2^h}]
 =[\zeta^{-1}+c^{2^h}]
 =[\zeta+c^{-2^h}].
\]

On exponents modulo fifteen the two transformations are

\[
 R(m)=5-(2^s\bmod15)m,\qquad
 S(m)=-(2^h\bmod15)m.
\]

For odd \(k\) they are \((5-2m,7m)\), and for even \(k\) they are
\((5-8m,13m)\).  Iterating either pair on the twelve residues outside
\(5\mathbb Z/15\mathbb Z\) gives the table
\((\mathsf{L}_{16,3})\).  Since every \(A^j\) is nonzero modulo
\(\ell\), all twelve classes vanish exactly when \(v(u)\) vanishes,
proving \((\mathsf{L}_{16,4})\) and its nonvanishing counterpart.

Applying \(\chi\) to the table gives each \(t^{A^j}\) twice, while
the four \(\F_4\)-terms are one.  If \(t=1\), the sum is sixteen;
conversely equality with the positive real number sixteen forces every
unit-modulus summand to be one by equality in the triangle inequality.
\end{proof}

\begin{proposition}[semiprimitive exhaustion and the saturated planes]
\label{prop:dk-semiprimitive}
\PROVED{} Retain Proposition~\ref{prop:dk-line-saturation}, put

\[
 h=3^k,\qquad q=2^h,\qquad
 K=\F_{q^2},\qquad L=\F_{q^4},
\]

and, for a prime \(\ell\mid\Psi_k\), set

\[
 H_K=(K^\times)^\ell,\qquad H_L=(L^\times)^\ell.
\]

Every nontrivial order-\(\ell\) character \(\chi\) on \(K^\times\)
has, for the canonical additive character,

\[
 G_K(\chi^j)=q\qquad(1\leq j<\ell).                 \tag{\(\mathsf{SP}_1\)}
\]

Every order-\(\ell\) character on \(L^\times\) is the norm lift of
one on \(K^\times\), and quadratic Hasse--Davenport gives

\[
 G_L(\widetilde\chi^j)=-q^2\qquad(1\leq j<\ell).    \tag{\(\mathsf{SP}_2\)}
\]

Consequently, whenever \(a,b,a+b\not\equiv0\pmod\ell\),

\[
 J_K(\chi^a,\chi^b)=q,\qquad
 J_L(\widetilde\chi^a,\widetilde\chi^b)=-q^2,      \tag{\(\mathsf{SP}_3\)}
\]

while \(J(\chi^a,\chi^{-a})=-1\) in either field.

Failure of \(D'_k\) at \(\ell\) has the following two equivalent
additive forms:

\[
 W:=\F_{16}+\F_{16}\zeta,\qquad
 W\setminus\{0\}\subset H_L,                       \tag{\(\mathsf{SP}_4\)}
\]

and

\[
 V:=\F_4+\F_4\zeta+\F_4\zeta^2,\qquad
 V\setminus\{0\}\subset H_K.                       \tag{\(\mathsf{SP}_5\)}
\]

Thus a bad translate saturates an eight-dimensional \(\F_2\)-space in
\(L\), and its norm-down form saturates a six-dimensional
\(\F_2\)-space in \(K\), not merely the original affine line.

The exact semiprimitive Fourier equations nevertheless permit both
planes.  If \(W^\perp\) is the additive annihilator of \(W\), put

\[
 A_L=|(W^\perp\setminus\{0\})\cap H_L|,\qquad R=q^2,\qquad w=256.
\]

Then orthogonality and \((\mathsf{SP}_2)\) force

\[
 A_L=\frac{(R-1)(R+w)-\ell R(w-1)}{\ell w},
 \qquad
 |W^\perp\setminus(H_L\cup\{0\})|
 =\frac{(R-1)(R+w)(\ell-1)}{\ell w}.                \tag{\(\mathsf{SP}_6\)}
\]

Likewise, for \(w_0=64\),

\[
 A_K:=|(V^\perp\setminus\{0\})\cap H_K|
 =\frac{(q+1)(q-w_0)+\ell q(w_0-1)}{\ell w_0},       \tag{\(\mathsf{SP}_7\)}
\]

and

\[
 |V^\perp\setminus(H_K\cup\{0\})|
 =\frac{(\ell-1)(q+1)(q-w_0)}{\ell w_0}.            \tag{\(\mathsf{SP}_8\)}
\]

All four quantities are nonnegative integers with the required total
sizes.  Hence the complete semiprimitive Gauss and Jacobi table, even
together with the dimensions of the saturated spaces, gives no
contradiction.  A proof of \(D'_k\) must use the distinguished
cyclotomic placement of \(W\) or \(V\), not only the semiprimitive
character table or additive-subspace incidence.
\end{proposition}

\begin{proof}
Since \(\ord_\ell(2)=2h\), one has \(q\equiv-1\pmod\ell\).
Thus every order-\(\ell\) character of \(K^\times\) is trivial on
\(\F_q^\times\).  Decompose \(K^\times\) into its \(q+1\)
projective \(\F_q\)-lines.  On the trace-zero line \(\F_q\), the
additive sum is \(q-1\); on every other line it is \(-1\).  The induced
nontrivial projective character sums to zero, giving
\(G_K(\chi)=q\).  The same argument applies to every nontrivial power.
The norm induces an isomorphism on the order-\(\ell\) character groups,
so Hasse--Davenport proves \((\mathsf{SP}_2)\).  The standard quotient
formula for Jacobi sums proves \((\mathsf{SP}_3)\), with the inverse
case evaluated separately.

If \(D'_k\) fails, Proposition~\ref{prop:dk-line-saturation} puts every
\(\zeta+c\), \(c\in\F_{16}\), in \(H_L\).  Since
\(\F_{16}^\times\subset H_L\), every nonzero
\(a\zeta+b\) is an \(\ell\)-th power, proving
\((\mathsf{SP}_4)\); the converse follows by taking \(a=1,b=u\).

For \((\mathsf{SP}_5)\), every polynomial of degree at most two over
\(\F_4\) is a product of linear factors or is an irreducible quadratic.
In the first case its value at \(\zeta\) is an \(\ell\)-th power by
\((\mathsf{SP}_4)\).  In the second, choose a root
\(c\in\F_{16}\setminus\F_4\); its value is

\[
 (\zeta+c)(\zeta+c^4)
 =\Norm_{L/K}(\zeta+c),
\]

again an \(\ell\)-th power.  Conversely the corrected norm
\(N=\zeta^2+\zeta+\omega\) belongs to \(V\); membership in \(H_K\)
is exactly failure of \(D'_k\).

Finally the two additive Fourier values of \(H_L\) are

\[
 \frac{-1-R(\ell-1)}\ell\quad\hbox{on }H_L,
 \qquad
 \frac{R-1}\ell\quad\hbox{off }H_L,
\]

by \((\mathsf{SP}_2)\).  Fourier inversion applied to
\(|W\cap H_L|=w-1\) gives \((\mathsf{SP}_6)\).  The corresponding
values for \(H_K\) are

\[
 \frac{-1+q(\ell-1)}\ell\quad\hbox{on }H_K,
 \qquad
 \frac{-1-q}\ell\quad\hbox{off }H_K,
\]

and \(|V\cap H_K|=w_0-1\) gives
\((\mathsf{SP}_7)\)--\((\mathsf{SP}_8)\).  Integrality follows from
\(\ell\mid q+1\), \(w\mid q^2\), and \(w_0\mid q\); nonnegativity is
immediate from \(\ell\leq q+1\) and \(q\geq2^9\).
\end{proof}

\begin{proposition}[cyclotomic Dirichlet coefficient and subspace-polynomial exhaustion]
\label{prop:dk-dirichlet-exhaustion}
\PROVED{} Retain \(k\geq2\), put \(h=3^k\), and write
\[
 K=\mathbb F_{4^h}=\mathbb F_4(\zeta),\qquad
 \zeta^h=\omega.
\]
Let \(\ell\mid\Psi_k\), let
\(\chi:K^\times\to\mu_\ell\) have exact order \(\ell\), realized
complexly, and let
\(\widetilde\chi=\chi\circ\operatorname{Norm}_{\mathbb F_{16^h}/K}\).
Put
\[
 s=h/3,\qquad A=2^s\pmod\ell,
\]
choose \(u\in\mathbb F_{16}\) of order fifteen with \(u^5=\omega\),
and set \(t=\widetilde\chi(\zeta+u)\).  For \(0\leq r<h\), define
\[
 \mathcal A_r(\chi)=
 \sum_{\substack{f\in\mathbb F_4[X]\ {\rm monic}\\\deg f=r}}
       \chi(f(\zeta)).
\]
Then
\[
 \mathcal A_0=1,\qquad \mathcal A_1=4,\qquad
 \mathcal A_2=10+\sum_{j=0}^5t^{A^j}.                  \tag{\(\mathsf{DP}_1\)}
\]
Failure of \(D'_k\) at \(\ell\) is equivalent to the single extremal
coefficient
\[
 \mathcal A_2=16.                                      \tag{\(\mathsf{DP}_2\)}
\]
If \(\mathcal A'_1\) is the degree-one coefficient after the quadratic
constant-field extension to \(\mathbb F_{16}\), then
\[
 \mathcal A'_1=2\mathcal A_2-\mathcal A_1^2
               =2\mathcal A_2-16.                     \tag{\(\mathsf{DP}_3\)}
\]
Thus the sixteen-term \(\mathbb F_{16}\)-line extremum and the
six-term irreducible-quadratic extremum are the same constraint under
constant-field base change, not two independent conditions.

More precisely, let
\[
 L_\chi(T)=\sum_{r=0}^{h-1}\mathcal A_rT^r
          =(1-T)L_\chi^*(T),\qquad
 L_\chi^*(T)=\sum_{r=0}^{d}b_rT^r,\qquad d=h-2,
\]
so \(b_r=\sum_{j=0}^r\mathcal A_j\).  The functional equation has root
number \(+1\) and gives
\[
 b_{d-r}=2^{d-2r}\overline{b_r}.                       \tag{\(\mathsf{DP}_4\)}
\]
If \(D'_k\) fails, then
\[
 (b_0,b_1,b_2)=(1,5,21),\qquad
 b_{d-2}=21\cdot2^{h-6}.                               \tag{\(\mathsf{DP}_5\)}
\]
For
\[
 V=\langle1,\zeta,\zeta^2\rangle_{\mathbb F_4},
\]
the distinguished trace annihilator is
\[
 V^\perp
 =\zeta\{f(\zeta):f\in\mathbb F_4[X],\ \deg f\leq h-4\}.
                                                               \tag{\(\mathsf{DP}_6\)}
\]
Consequently, in the bad case,
\[
 \sum_{x\in V^\perp}\chi(x)
 =3b_{d-2}=63\cdot2^{h-6},                             \tag{\(\mathsf{DP}_7\)}
\]
exactly the value forced independently by additive Poisson summation
and \(G_K(\chi)=2^h\).  Thus the first two Dirichlet coefficients,
the exact cyclotomic trace-dual placement, the semiprimitive Gauss
sign, the functional equation, and the Riemann hypothesis are mutually
compatible with failure.

Finally put
\[
 D=\zeta^4+\zeta,\qquad E=D^5(D^3+1).
\]
The \(\mathbb F_4\)-subspace polynomial of \(V\) is
\[
\begin{split}
 P_V(Z)={}&Z^{64}
 +\bigl((1+D^3)^4+E^3\bigr)Z^{16}\\
 &+\bigl(D^{12}+E^3(1+D^3)\bigr)Z^4
 +E^3D^3Z.
\end{split}                                             \tag{\(\mathsf{DP}_8\)}
\]
Both \(D\) and \(E\) are \(\ell\)-th powers unconditionally.  Hence
the derivative coefficient \(E^3D^3\) supplies no missing nonresidue;
a subspace-polynomial proof would have to use an additive relation
among the \(Z^{16}\)- and \(Z^4\)-coefficients rather than merely
products of the saturated elements.
\end{proposition}

\begin{proof}
The four monic linears \(X+c\), \(c\in\mathbb F_4\), evaluate to
\(\ell\)-th powers, so \(\mathcal A_1=4\).  Of the sixteen monic
quadratics over \(\mathbb F_4\), ten are reducible and therefore also
have character value one.  The six irreducibles are
\[
 f_c(X)=(X+c)(X+c^4),
\]
indexed by the six Frobenius pairs in
\(\mathbb F_{16}\setminus\mathbb F_4\).  Since
\[
 \widetilde\chi(\zeta+c)=\chi(f_c(\zeta)),
\]
the orbit table of Proposition~\ref{prop:dk-line-saturation} gives
their six values exactly \(t^{A^j}\), \(0\leq j<6\).
This proves \((\mathsf{DP}_1)\).  The six summands have modulus one,
so \(\mathcal A_2=16\) if and only if all are one; the \(j=0\) term
then gives \(t=1\), which is exactly failure of \(D'_k\).

The reciprocal roots of \(L_\chi\) are squared under quadratic
constant-field extension.  Newton's identity therefore gives
\[
 \mathcal A'_1=2\mathcal A_2-\mathcal A_1^2,
\]
proving \((\mathsf{DP}_3)\).

The character is primitive modulo \(X^h+\omega\) and even because it
is trivial on \(\mathbb F_4^\times\).  Standard function-field
Dirichlet theory gives \(L_\chi=(1-T)L_\chi^*\), with
\(\deg L_\chi^*=d=h-2\), and the coefficient relation in
\((\mathsf{DP}_4)\) up to its root number.  We determine that sign
directly.  Let \(U\) be the \(\mathbb F_4\)-space of polynomials in
\(\zeta\) of degree at most \(h-2\).  Newton sums for
\(X^h+\omega\) give
\[
 \operatorname{Tr}(\zeta^r)=0\qquad(1\leq r<h),
\]
so \(U^\perp=\mathbb F_4\zeta\).  Additive Poisson summation and
\(\chi(\mathbb F_4^\times\zeta)=1\) give
\[
 3b_d=\sum_{x\in U}\chi(x)=\frac34G_K(\chi).
\]
Proposition~\ref{prop:dk-semiprimitive} gives
\(G_K(\chi)=2^h\), so
\[
 b_d=2^{h-2}=2^d.
\]
The root number is therefore \(+1\).  Under failure,
\(b_2=1+4+16=21\), and \((\mathsf{DP}_4)\) at \(r=2\)
gives \((\mathsf{DP}_5)\).

For \(x=\sum_{j=0}^{h-1}x_j\zeta^j\), the same Newton-sum calculation
gives
\[
 \operatorname{Tr}(x)=x_0,\qquad
 \operatorname{Tr}(x\zeta)=\omega x_{h-1},\qquad
 \operatorname{Tr}(x\zeta^2)=\omega x_{h-2}.
\]
This proves \((\mathsf{DP}_6)\).  Every nonzero polynomial of degree
at most \(h-4\) is a unique \(\mathbb F_4^\times\)-multiple of a
monic one, while \(\chi(\zeta)=1\).  Hence
\[
 \sum_{x\in V^\perp}\chi(x)=3b_{h-4}=3b_{d-2}.
\]
Direct Poisson summation under
\(V\setminus\{0\}\subset\ker\chi\) gives
\((63/64)G_K(\chi)\), the same value
\(63\cdot2^{h-6}\).  This proves \((\mathsf{DP}_7)\).

It remains to derive the subspace polynomial.  First put
\(U_1=\langle1,\zeta\rangle_{\mathbb F_4}\).  Since
\(\prod_{c\in\mathbb F_4}(Z+c)=Z^4+Z\),
\[
\begin{split}
 P_{U_1}(Z)
 &=(Z^4+Z)^4+D^3(Z^4+Z)\\
 &=Z^{16}+(1+D^3)Z^4+D^3Z.
\end{split}
\]
Evaluating at \(\zeta^2\) gives
\[
 E=P_{U_1}(\zeta^2)=D^8+D^5=D^5(D^3+1).
\]
Because \(P_{U_1}\) is \(\mathbb F_4\)-linearized,
\[
 P_V(Z)=P_{U_1}(Z)^4+E^3P_{U_1}(Z),
\]
which expands to \((\mathsf{DP}_8)\).

Finally,
\[
 D=\prod_{c\in\mathbb F_4}(\zeta+c),
\]
so \(\chi(D)=1\).  Also \(E\) is the product of all sixteen
monic-quadratic values.  The ten reducible values have character one,
and the other six multiply to
\[
 t^{\sum_{j=0}^5A^j}=1
\]
because \(A\) has order six modulo \(\ell\).  Thus \(\chi(E)=1\)
without assuming failure, which finishes the proof.

\end{proof}

\begin{proposition}[relative Dickson form and its degree ceiling]
\label{prop:dk-relative-dickson}
\PROVED{} Retain

\[
 s=3^{k-1},\quad a=2^s,\quad E=\F_a,\quad K=\F_{a^6},
 \quad \Lambda_k=a^2-a+1,
\]

and set \(T_k=\{x\in K^\times:x^{\Lambda_k}=1\}\).  The element
\(P_k\in T_k\) has degree six over \(E\).  Put

\[
 Y_k=P_k+P_k^{-1}\in\F_{a^3},\qquad
 G_k(T)=\prod_{j=0}^2(T+Y_k^{a^j})\in E[T].
\]

Then \(G_k\) is irreducible of degree three and
\(X^3G_k(X+X^{-1})\) is the irreducible reciprocal minimal polynomial
of \(P_k\) over \(E\).  If \(\mathcal D_\ell\) is characterized by

\[
 \mathcal D_\ell(v+v^{-1})=v^\ell+v^{-\ell},
\]

then, for every prime \(\ell\mid\Psi_k\),

\[
 P_k\notin T_k^\ell
 \iff G_k(\mathcal D_\ell(X))
      \text{ is irreducible over }E.                 \tag{\(\mathsf{RD}\)}
\]

In the bad case the composition splits over \(\F_{a^3}\), and all
its irreducible factors over \(E\) have degree dividing three.
Moreover

\[
 T_k\cap\F_{a^2}^\times=\mu_3,\qquad
 T_k\cap\F_{a^3}^\times=\{1\}.
\]

Thus the elements of \(T_k\) of relative degree below six are exactly
\(\mu_3\), while the bad subgroup \(T_k^\ell\) contains exactly

\[
 \frac{\Lambda_k}{\ell}-3
\]

elements of exact degree six.  Except in the already settled case
\(\Lambda_k/\ell=3\), full relative degree and reciprocal-torus
structure cannot prove \((D'_k)\).
\end{proposition}

\begin{proof}
Proposition~\ref{prop:dk-projection} gives
\([\F_2(P_k):\F_2]=6s\), so \(E(P_k)=K\).  Since
\(a^3+1=(a+1)\Lambda_k\), one has
\(P_k^{a^3}=P_k^{-1}\), whence \(Y_k\in\F_{a^3}\).  The element
\(P_k\) is quadratic over \(E(Y_k)\); its relative degree six forces
\([E(Y_k):E]=3\), proving the minimal-polynomial assertions.

Suppose \(P_k\notin T_k^\ell\).  The whole current
\(\ell\)-primary part of \(K^\times\) lies in \(T_k\), so this is
equivalent to \(P_k\notin(K^\times)^\ell\).  Kummer theory and
Capelli make

\[
 X^{3\ell}G_k(X^\ell+X^{-\ell})
\]

irreducible of degree \(6\ell\).  If
\(\alpha^\ell=P_k\) and \(t=\alpha+\alpha^{-1}\), then
\(\mathcal D_\ell(t)=Y_k\).  Since \(\alpha\) is at most quadratic
over \(E(t)\), the degree of \(t\) is at least \(3\ell\); it is a root
of the degree-\(3\ell\) polynomial
\(G_k(\mathcal D_\ell(X))\), proving the forward implication in
\((\mathsf{RD})\).

Conversely, if \(P_k=\alpha^\ell\) with \(\alpha\in T_k\), then
\(\alpha^{a^3}=\alpha^{-1}\), so
\(t\in\F_{a^3}\).  The same holds after multiplying \(\alpha\) by
any \(\ell\)-th root of unity and after applying the three
\(E\)-conjugates.  Hence every root of the Dickson composition lies in
\(\F_{a^3}\), proving the bad-case factorization statement and the
reverse implication.

Finally the two displayed intersections have orders

\[
 \gcd(\Lambda_k,a^2-1)=3,
 \qquad \gcd(\Lambda_k,a^3-1)=1.
\]

Every proper intermediate field of \(K/E\) lies in \(\F_{a^2}\) or
\(\F_{a^3}\), so only \(\mu_3\) has relative degree below six.  The
cyclic group \(T_k^\ell\) has order \(\Lambda_k/\ell\) and contains
\(\mu_3\), because \(\ell\ne3\).  Removing those three elements gives
the asserted exact count.
\end{proof}

\begin{proposition}[trace-degree ceiling for the fibotomic criterion]
\label{prop:dk-trace-degree-ceiling}
\PROVED{} Retain Proposition~\ref{prop:dk-fibotomic}, and let
\[
 \mathcal I_\ell:=\pi\bigl(U^\ell\setminus\{1\}\bigr)\subseteq\F_q
 \qquad(\ell\text{ a prime divisor of }\Psi_k).
\]
Every element of \(\mathcal I_\ell\) has absolute trace one, and
\[
 |\mathcal I_\ell|
 =\frac{(q+1)/\ell-1}{2}
 \geq\frac{3a+2}{2}>a.                                \tag{\(\mathsf T\)}
\]
At least \(a/2+1\) elements of \(\mathcal I_\ell\) have exact degree
\(h\) over \(\F_2\).  Consequently, absolute trace one together with
full degree \(h\) cannot prove that the distinguished
\(Z\) lies outside \(\pi(U^\ell)\).
\end{proposition}

\begin{proof}
The two-to-one trace-one parametrization in
Proposition~\ref{prop:dk-fibotomic} restricts to \(U^\ell\).  This
cyclic subgroup has order \((q+1)/\ell\), which is odd, so its only
self-inverse element is \(1\).  This proves the equality in
\((\mathsf T)\).  Since \(\ell\leq\Psi_k=\Lambda_k/3\) and
\(q+1=(a+1)\Lambda_k\), it also gives
\[
 |\mathcal I_\ell|\geq\frac{3(a+1)-1}{2}
 =\frac{3a+2}{2}.
\]
Every proper subfield of \(\F_q\) lies in
\(\F_{2^{h/3}}=\F_a\), because \(h=3^k\).  Hence at most \(a\)
elements of \(\mathcal I_\ell\) have degree smaller than \(h\), leaving
at least \((3a+2)/2-a=a/2+1\) of exact degree \(h\).

For completeness, the distinguished \(Z\) itself has exact degree
\(h\).  Indeed, \((\mathsf D_3)\) and \(P_k\in\F_2(M_k)\) force
\([\F_2(M_k):\F_2]=2h\).  The element \(M_k\) is quadratic over
\(\F_2(Z)\), satisfying
\[
 X^2+Z^{-1}X+1=0,
\]
while \(Z\in\F_q\).  Therefore \([\F_2(Z):\F_2]=h\).
\end{proof}

\begin{proposition}[all proper-subfield norms are blind]
\label{prop:dk-all-norms-blind}
\PROVED{} Retain \(k\geq2\) and the notation of
Proposition~\ref{prop:dk-current-level}.  Let \(\ell\) be a prime
divisor of \(\Psi_k\), and let \(d\) be any proper divisor of \(2h\).
Then
\[
 \Norm_{\F_{2^{2h}}/\F_{2^d}}
\]
annihilates the entire \(\ell\)-Sylow subgroup of
\(\F_{2^{2h}}^\times\).  Equivalently, multiplying an element by any
current \(\ell\)-primary component leaves every multiplicative norm to
a proper subfield unchanged.
\end{proposition}

\begin{proof}
The norm exponent is
\[
 E_d=\frac{2^{2h}-1}{2^d-1}.
\]
Since \(\ord_\ell(2)=2h\) and \(d<2h\), one has
\(\ell\nmid2^d-1\).  Therefore
\[
 v_\ell(E_d)=v_\ell(2^{2h}-1),
\]
so exponentiation by \(E_d\) kills the full \(\ell\)-Sylow subgroup.
\end{proof}

\begin{remark}[erratum: current-level versus cumulative $D_k$]
\label{rem:dk-erratum}
\PROVED{} One has
\[
 \frac{2^{3^k}+1}{3^{k+1}}
 =\prod_{j=1}^k\frac{\Phi_{2\cdot3^j}(2)}3
 =\prod_{j=1}^k\Psi_j.
\]
The earlier version of this note called divisibility by the left-hand
side $D_k$ and claimed it equivalent to the current-level $m=4$
statement.  It additionally requires every \emph{old} exception factor
to divide $\ord(M_k)$, so it is strictly stronger unless an old-factor
propagation theorem is supplied.  Proposition~\ref{prop:twisted} gives
no such theorem.  The exact conjecture is $D'_k$ above.
\end{remark}

For $k=1$, $\Psi_1=1$ and $D'_1$ is vacuous.
\CERTIFIED{k\leq6} $D'_k$ holds for $k=2,\dots,6$: every prime factor of
the fully factored $\Phi_{2\cdot3^k}(2)$ divides $\ord(M_k)$. Hence
$m_p=4$ exactly --- not merely $\geq4$ --- for every prime $p$ with
$f(p)\in\{18,54,162,486,1458\}$:
\begin{center}
\begin{tabular}{rl}
\toprule
$f$ & primes $p$ with $m_p=4$ \\
\midrule
$18$   & $19$ \\
$54$   & $87211$ \\
$162$  & $163,\ 135433,\ 272010961$ \\
$486$  & $1459,\ 139483,\ 10429407431911334611,$ \\
       & $918125051602568899753$ \\
$1458$ & $227862073,\ 3110690934667,\ 216892513252489863991753,$ \\
       & $1102099161075964924744009,\ \mathrm{P}78$ \\
\bottomrule
\end{tabular}
\end{center}
($\mathrm{P}78$ is the certified $78$-digit prime of
$\Phi_{2\cdot3^6}(2)$; PRP-local here, factordb marks it proven.) The
anchor rows $m_{19}=m_{163}=m_{1459}=4$ (DiMuro / calculator) are
reproduced, never assumed; the remaining eleven rows are new.
Independent cross-checks: the term algebra of
\path{experiments/ordinal_excess_probe.py} re-derives the $p=19$ root
pattern ($m=0,\dots,3$ root, $m=4$ no root) and $m_{87211}=4$; an
explicit Artin--Schreier compositum model $\F[y]/(y^2+y+\omega)$
verifies $\sigma(4)=5$, the norm identity, and the full direct power
test $(\kappa+4)^{(2^{4h}-1)/p}\neq1$ on the smallest prime of each
level $k\leq6$.

\CONSISTENT{} $k=7,8$: every \emph{known} prime factor of
$\Phi_{2\cdot3^7}(2)$ and $\Phi_{2\cdot3^8}(2)$ --- seven and five
primes respectively; factordb CF entries verified locally by
divisibility, primality, and order tests, the small ones re-derived by
sieve, and the list cross-audited against the raw factordb API after an
adversarial review caught a dropped $43$-digit $k=7$ prime --- also has
$m_p=4$. An $m_p\geq5$ example, if one exists, now hides strictly
inside the unfactored cofactors; their prime multiplicities are likewise
unknown until factorization.

\begin{proposition}[twisted norm tower]\label{prop:twisted}
\PROVED{} With $\eta=\zeta^3=\kappa_{3^{k-1}}$,
\[
  \Norm_{\F_{2^{2\cdot3^k}}/\F_{2^{2\cdot3^{k-1}}}}(N_k)
  =\eta^2+\omega^2\eta+1 ,
\]
which is neither $N_{k-1}=\eta^2+\eta+\omega$, nor one of its Frobenius
conjugates (their exponent pattern is $(2^{j+1},2^j)$), nor a scaling
$N_{k-1}(\lambda\eta)$ with $\lambda\in\mu_3$.
\end{proposition}

\begin{proof}
The generator $\tau$ of the degree-$3$ step sends $\zeta\mapsto
\zeta^{c}$ with $c=2^{2\cdot3^{k-1}}$ a primitive cube root of unity
modulo $3^{k+1}$; the cube roots are $1+a3^k$, so
$\tau(\zeta)=\omega^{a}\zeta$ ($a\neq0$) while $\tau(\omega)=\omega$.
Hence the norm is $\prod_{r^3=\eta}(r^2+r+\omega)$ over the three cube
roots $r$ of $\eta$. Writing $\alpha,\beta$ for the two roots of
$T^2+T+\omega$ (so $\alpha+\beta=1$, $\alpha\beta=\omega$), each factor
is $(r+\alpha)(r+\beta)$, and the product collapses by the resolvent
\[
  \prod_{r^3=\eta}(r^2+r+\omega)
  =(\alpha^3+\eta)(\beta^3+\eta)
  =\eta^2+(\alpha^3+\beta^3)\eta+(\alpha\beta)^3
  =\eta^2+\omega^2\eta+1,
\]
since $\alpha^3+\beta^3=(\alpha+\beta)^3+\alpha\beta(\alpha+\beta)
=1+\omega=\omega^2$ and $(\alpha\beta)^3=\omega^3=1$. Machine-checked at
$k=2,3,4$.
\end{proof}

\begin{proposition}[trace form and norm blindness]
\label{prop:dk-trace}
\PROVED{} Retain the notation of
Proposition~\ref{prop:dk-current-level}, and put
\[
 \gamma_k:=\zeta+\zeta^{-1},\qquad
 A_k:=\gamma_k^2+\gamma_k+1,\qquad
 Q:=2^{2\cdot3^{k-1}}=2^{2h/3}.
\]
Set $\gamma_0=A_0=1$.
Then, with $\overline N=N^{2^h}$,
\[
 N+\overline N=A_k,\qquad N\overline N=A_k^Q,
 \qquad M_k+M_k^{-1}=A_k^{2-Q},
 \qquad Z=A_k^{Q-2}.                                    \tag{4}
\]
Moreover
\[
 \Norm_{\F_{2^{3^k}}/\F_{2^{3^{k-1}}}}(A_k)=A_{k-1},
 \qquad
 \Norm_{\F_{2^{3^k}}/\F_{2^{3^{k-1}}}}(Z)=A_{k-1}^{-1}. \tag{5}
\]
For every prime $\ell\mid\Psi_k$, however, this relative norm is
blind to the required $\ell$-power residue class.
\end{proposition}

\begin{proof}
Since $2^h\equiv-1\pmod{3h}$ and
$\zeta^h+\zeta^{2h}=1$, direct expansion gives
\[
 N+\overline N
 =\zeta^2+\zeta^{-2}+\zeta+\zeta^{-1}+1=A_k.
\]
The congruence $Q\equiv h+1\pmod{3h}$ follows inductively by cubing
$1+3^r$ modulo $3^{r+2}$.  Hence
$\gamma_k^Q=\zeta^{h+1}+\zeta^{-(h+1)}$, and another direct
expansion yields $N\overline N=A_k^Q$.  Therefore
\[
 M_k+M_k^{-1}
 =\frac{(N+\overline N)^2}{N\overline N}=A_k^{2-Q},
\]
and \(Z=\pi(M_k)=1/(M_k+M_k^{-1})\), which proves (4).

Over $B=\F_{2^{h/3}}$, the element $\gamma_k$ has minimal polynomial
$X^3+X+\gamma_{k-1}$.  Base-change this cubic extension by
$B_2=B\F_4=\F_{2^{2h/3}}$; its compositum is
$\F_{2^{2h}}$, and the degree remains three.  With
$\omega=\zeta^h\in\F_4$ one has
$A_k=(\gamma_k+\omega)(\gamma_k+\omega^2)$.  Evaluating that minimal
polynomial at $\omega$ and $\omega^2$ gives
\[
 \Norm_{\F_{2^{2h}}/B_2}(\gamma_k+\omega)
   =\gamma_{k-1}+\omega^2,
 \qquad
 \Norm_{\F_{2^{2h}}/B_2}(\gamma_k+\omega^2)
   =\gamma_{k-1}+\omega.
\]
Their product is $A_{k-1}$.  Since norm commutes with this scalar
extension, this is precisely
$\Norm_{\F_{2^h}/B}(A_k)$, proving the first identity in (5).
Furthermore \(Z=A_k^{Q-2}\), so
\[
 \Norm_{\F_{2^h}/B}(Z)=A_{k-1}^{Q-2}=A_{k-1}^{-1},
\]
because \(Q-2\equiv-1\pmod{2^{h/3}-1}\).
This proves the second identity.

If $\ell\mid\Psi_k$, then $\ell\nmid2^h-1$ and
$\ell\nmid2^{h/3}-1$.  Thus both $A_k$ and $A_{k-1}$ lie in
multiplicative groups on which raising to the $\ell$-th power is an
automorphism, so (5) is tautological on $\ell$-power classes.  Equivalently,
$\ord_\ell(Q)=3$ and $1+Q+Q^2\equiv0\pmod\ell$: the degree-three norm
annihilates the relevant character.  The surviving current-level
condition is the Capelli
test $R_k(X^\ell)$ irreducible from
Proposition~\ref{prop:dk-current-level}.
\end{proof}

So, unlike the $\gamma$ tower of Proposition~\ref{prop:tower},
old-prime full-order parts do \emph{not propagate by that argument}, and
$D'_k$ is genuinely per-level for the present method: the certification at
each $k$ stands on its own.  Structurally, at the degree-$3$ step with
$Q=2^{2\cdot3^{k-1}}$, the relative norm kernel has order
\[
 Q^2+Q+1=\Phi_{3^k}(2)\Phi_{2\cdot3^k}(2).
\]
Its gcd with $Q-1$ is $3$, so there is no direct-product splitting at
the $3$-part as in Proposition~\ref{prop:norm-kernel-seam}; nevertheless,
for every $\ell\ne3$ the new $\ell$-Sylow subgroups of both columns lie
entirely in the norm kernel.  This is the
prime-to-$3$ analogue of Proposition~\ref{prop:norm-kernel-seam} and
explains why a lower-level norm identity cannot determine $D'_k$.

\begin{remark}[the $3$-part]\label{rem:threepart}
\PROVED{} Proposition~\ref{prop:dk-projection} and
\(v_3(a+1)=k\) imply
\[
 v_3(\ord M_k)=k+1\quad(k\text{ odd}),\qquad
 v_3(\ord M_k)\leq k\quad(k\text{ even}).
\]
Thus the odd-level full \(3\)-part is now explained analytically.
\CONSISTENT{} At even levels \(2\leq k\leq6\), the sharper observed
value is \(v_3(\ord M_k)=k-1\); the remaining one-step deficiency is
not determined by the projection theorem.
\end{remark}

\section{Boundedness}\label{sec:bounded}

\CONJECTURED{} (Lenstra~\cite{Lenstra77}) $m_p$ is globally bounded.
Lenstra gave unconditional lower-bound rules (singleton odd $\Qset$
forces positive excess; $f(p)=2\cdot3^k$ forces excess at least $4$) but
left absolute boundedness open.

\PROVED{} If the $0/1/4$ candidate holds, then $m_p\leq4$ for all $p$;
proving Conjecture~\ref{conj:rule} would settle boundedness.

\PROVED{} No proof for arbitrary finite-field elements can come from a
fixed initial translate window: Proposition~\ref{prop:translate-no-go}
constructs elements of arbitrarily large exact degree for which every
translate in any prescribed finite window is a $p$-th power.  A boundedness
proof must therefore use the distinguished Conway-tower identities.

\PROVED{} The $3$-power column cannot produce $m\in\{2,3\}$
(Proposition~\ref{prop:column}); any $C_k$ failure jumps directly to
$m_r\geq4$.

\PROVED{} (conditional on $C_k$) For every prime $r$ with $f(r)=3^j$,
$j\leq k$: $C_k$ implies $m_r=1$. With $C_k$ certified for $k\leq6$, all
primes in the $3$-power column with $f(r)\leq6561$ have $m_r=1$
unconditionally.

\CERTIFIED{k\leq6} (June 2026 pass) On the $2\cdot3^k$ exception column
the matching upper bound is now certified at every visible prime:
$m_p=4$ exactly for every prime with $f(p)=2\cdot3^k$, $k\leq6$, and for
every known prime at $k=7,8$ (Section~\ref{sec:m4}). Boundedness on the
exception column therefore holds unconditionally through $f=1458$ and is
equivalent at each new level to $D'_k$
(Conjecture~\ref{conj:dk}) beyond.

\OPEN{} Boundedness outside the $3$-power and $2\cdot3^k$ columns. The
non-cyclotomic singleton chains (e.g.\ the $11$-chain and the $23$, $29$,
$47$ components) have $m=1$ rows in the calculator data and OEIS, but no
order-criterion proof covers them the way $C_k$ covers the $3$-power
column.

\OPEN{} Whether any prime $p$ has $m_p\geq5$. No counterexample is known;
the question is not known to be decidable by any current structural
argument outside the $3$-adic columns.

\PROVED{} A proof of the full $0/1/4$ candidate would in particular prove
both the Conway--Fermat maximality statement
$\delta_n=F_n$ at every quadratic level
(Proposition~\ref{prop:fermat-arm}) and the Gaussian-period primitivity
statement $C_k$ at every cubic level
(Proposition~\ref{prop:gauss-arm}).  These are independent new-factor
problems; neither follows from the existing norm recursions.

\section{Ranked next moves}\label{sec:moves}

\begin{enumerate}[leftmargin=*]
\item \emph{Conway--Fermat quotient lemma (highest analytical priority).}
Prove or refute $\delta_n=F_n$ in Proposition~\ref{prop:fermat-arm}.
Proposition~\ref{prop:fermat-fibotomic} reduces it to the explicit
nonvanishing tests $S_{F_n/\ell}(a_{n-1})\ne0$, while
Proposition~\ref{prop:fermat-countermodels} proves that degree, trace,
norm, and the top Artin--Schreier equation cannot suffice.  The missing
step is to show that the resultant recursion of
Proposition~\ref{prop:fermat-resultant} always selects the full-order
fibotomic stratum.  Proposition~\ref{prop:fermat-one-step-no-go} proves
that predecessor maximality and the exact one-step operator are not enough:
an alternative full-order \(A_3\)-factor maps to the proper
\(F_5/641\)-stratum.  Proposition~\ref{prop:fermat-ancestry-exhaustion}
now gives a necessary-and-sufficient trace/norm ladder for the
distinguished factor's entire earlier ancestry.  It also proves that
every actual trace fibre contains many normal full-degree
\(\ell\)-th powers, and that imposing the actual relative norm compresses
exactly to the original test
\(S_{(q+1)/\ell}(a_{n-1})=0\).  Thus the missing input is not another
coarse ancestry invariant.  Equivalently, Proposition~\ref{prop:fermat-kummer}
asks for the nontrivial Kummer symbol of \(c_n\) at the distinguished
point of the supersingular correspondence; its quarter-turn Frobenius
symmetry is proved character-neutral and, by
Proposition~\ref{prop:fermat-elliptic-kummer}, does not lift to the
ramified Kummer cover.  Proposition~\ref{prop:fermat-generalized-jacobian}
places the surviving Miller-unit class in the toric part of the three-point
generalized Jacobian.  Proposition~\ref{prop:fermat-gj-character} exhausts
the extension class and non-rational elliptic torsion, while
Proposition~\ref{prop:fermat-rosenlicht} gives both its exact
Rosenlicht/Miller coordinate and a rational five-torsion splitting.  Their
evident reciprocity collapse is proved circular.  The next step is therefore
to prove that this now-explicit distinguished value of \(x+2y\) is nonzero,
not to seek another coordinate model or ordinary elliptic pairing.
\item \emph{New-prime lemma for the $3$-power spine.}  For every
$\ell^a\mid\Phi_{3^k}(2)$, prove
$v_\ell(\ord\gamma_k)=a$.  The norm tower already carries all old
factors, so this single per-level statement is exactly $C_k$.  Equivalently,
prove that the distinguished root $\eta_k$ of
$X^3+\eta_{k-1}X^2+(\eta_{k-1}+1)X+1$ generates the norm-one torus
$U_k$.  The sharp analytical languages are therefore either the
$\ell$-power residue symbol of the cyclotomic unit $1+\zeta$ or the
character theory of this explicit cubic trace/norm fibre, not another
order census.  Proposition~\ref{prop:cubic-partitions} removes every
prime-index subgroup below its partition bound, while
Proposition~\ref{prop:cubic-short-period} identifies failure with the
extremal value of one short cyclotomic period; a completion must control
the remaining characters.  Proposition~\ref{prop:cubic-secant-saturation}
shows that failure saturates all cyclotomic secants, while the associated
Kummer cover is rational; the missing input is therefore the exact Artin
value of the distinguished cyclotomic unit rather than a generic Weil bound.
Proposition~\ref{prop:cubic-ray-extension} identifies the conditional
Galois obstruction as a locally admissible
\(C_\ell\rtimes_2C_{3^k}\) ray extension, so local reciprocity alone cannot
exclude it.  Proposition~\ref{prop:cubic-reciprocity-no-go} sharpens this
to an exact exceptional-zero statement: norm operators, the plus
Stickelberger element, and \(2-\sigma\) all vanish on the one-dimensional
\(2\)-eigenspace, while Hilbert reciprocity multiplies the complete
conjugate orbit of the desired local symbol to an automatic identity.
The next arithmetic object is therefore a first-order unit regulator in
that simple exceptional eigenspace.
\item \emph{General transverse reciprocity.}  Evaluate the power-residue
symbol of $\Theta_{q,s}$ from
Proposition~\ref{prop:singleton-transverse}.  A theorem forcing its full
exact-order support would prove the singleton-one arm simultaneously for
the $11$-chain and the $23$, $29$, $47$, $179$, \ldots{} components.
The easy Kummer norm is provably the wrong object.
\item \emph{Exception-column current-factor lemma.}  Prove $D'_k$,
equivalently prove $R_k(X^\ell)$ irreducible for every
$\ell\mid\Psi_k$ in Proposition~\ref{prop:dk-current-level}, or prove
the polynomial in Proposition~\ref{prop:dk-fibotomic} rootless for its
distinguished \(Z\).  Proposition~\ref{prop:dk-projection} already
forces one current prime and every factor above an explicit threshold.  The
twisted norm formula in Proposition~\ref{prop:twisted} and the norm
blindness in Proposition~\ref{prop:dk-trace} rule out direct induction
on $k$; a residue-symbol argument for
$N_k=\zeta^2+\zeta+\omega$ must account for the twist.
Proposition~\ref{prop:dk-translated-binomial} and
Proposition~\ref{prop:dk-relative-dickson} give equivalent explicit
translated-binomial and relative Dickson irreducibility tests, and prove
that full relative degree is still too weak.
Proposition~\ref{prop:dk-line-saturation} further shows that failure
saturates the complete \(\F_{16}\)-affine line; a completion must make its
explicit sixteen-term character sum non-extremal by information beyond
Frobenius and base-field translation.
Proposition~\ref{prop:dk-semiprimitive} strengthens the bad case to complete
six- and eight-dimensional additive spaces and proves that the exact
semiprimitive Gauss--Jacobi/Fourier incidence table permits them.  Generic
semiprimitive estimates therefore cannot supply the missing identity.
Proposition~\ref{prop:dk-dirichlet-exhaustion} uses the distinguished
cyclotomic placement and identifies failure with the extremal second
Dirichlet coefficient \(\mathcal A_2=16\).  Its functional equation,
trace dual, root number, constant-field base change, and multiplicative
subspace-polynomial coefficient all reproduce identities compatible with
that extremum.  A completion must evaluate this explicit six-term period
non-extremally or introduce an independent relation involving higher
Dirichlet coefficients.
\item \emph{Numerical work only as falsification and audit.}  The
$m_{359}$/$m_{719}$ transverse norm, complete factorizations at $k=7,8$,
and exact-order censuses remain useful counterexample searches and
checks of any proposed lemma, but they no longer lead the proof program.
\end{enumerate}

\section{Provenance and validation}\label{sec:provenance}

\subsection{Independent oracle}

\path{experiments/ordinal_excess_probe.py} is a small local term-algebra
oracle -- not a replacement for CGSuite~\cite{CGSuite} or the C++
calculator~\cite{Peeters}, but a way to verify the first subtle cases
without using the Rust production tower as its own oracle. It implements
the impartial term algebra used by the calculators, a
multiplicative-order test for small component fields, and a fixed-base
exponentiation path (ported from the calculator's strategy) for targeted
root tests where full order factorization is unnecessary. Current probe
output:

\begin{center}
\begin{tabular}{rrlc}
\toprule
$p$ & $m$ & $\Qset$ & root? \\
\midrule
$7$  & $0$ & $\{3\}$  & yes \\
$7$  & $1$ & $\{3\}$  & no \\
$19$ & $1$ & $\{9\}$  & yes \\
$19$ & $4$ & $\{9\}$  & no \\
$73$ & $1$ & $\{9\}$  & no \\
$89$ & $1$ & $\{11\}$ & no \\
$47$ & $1$ & $\{23\}$ & no \\
$179$ & $1$ & $\{89\}$ & no \quad(\texttt{--deep}) \\
\bottomrule
\end{tabular}
\end{center}

The $p=47$ line certifies $m_{47}=1$ using only lower verified rows; since
$f(47)=23$ and $\Qset(23)=\{23\}$, the carry is
$\alpha_{47}=\kappa_{23}+1=\omega^{(\omega^7)}+1$, the value now
implemented in \path{tower.rs}. The $p=179$ line is the first locally
certified dependency row beyond the quick probe path; it still uses only
the already-certified lower row $m_{89}=1$.

\subsection{The $\alpha_{47}$ promotion}

\emph{Superseded 2026-06-13:} this pass moved the boundary to
$\alpha_{53}$ by promoting a single locally verified row; the finite
excess table is now the full OEIS A380496 b-file ($126$ rows, odd primes
$3 \le p \le 709$), so the refusal boundary sits at $\alpha_{719}$ and the
boundary test is \path{boundary_returns_none_past_prime_709}. The original
pass is recorded below for provenance.

The shipped state at the time: \path{src/scalar/big/ordinal/tower.rs} carries
$\alpha_{47}$ with the landmark test
\path{locally_verified_alpha_47_landmark} and moves the refusal
boundary to $\alpha_{53}$; the documented boundaries in
\path{nim.rs}, \path{mod.rs}, \path{src/games/nimber_game.rs},
\texttt{README.md}, and \texttt{docs/OPEN.md} were updated in step. The full
gate ran clean: \texttt{cargo fmt --check}, \texttt{cargo test},
\texttt{cargo check}/\texttt{clippy} with and without the
\texttt{python} feature (warnings denied), and the probe under
\texttt{py\_compile}.

\subsection{Supporting probes}

The $3^k$-family certification chain is
\path{experiments/cyclotomic_3k_family.py} (committed, maintained),
joined 2026-06-12 by \path{experiments/exception_column_m4.py} (the
$2\cdot3^k$ exception column, Section~\ref{sec:m4}; committed,
maintained, stdlib-only, $\sim$2 minutes), with the research-run probes
under \path{experiments/excess/} (machine-generated, rescued from
ephemeral storage; triage before citing). The analysis in
Sections~\ref{sec:norm}--\ref{sec:bounded} synthesizes the 2026-06-10
parallel research run plus the 2026-06-12 exception-column pass; each
claim above carries its own tag, and no claim depends on an unverified
solver.

\section*{Acknowledgements}
The \texttt{ogdoad} library, the experiments, and this draft were
developed with LLM assistance; the analysis consolidates a structured
parallel research run (2026-06-10) and an independent adversarial proof
audit (2026-07-20). Mathematical direction, framing, and responsibility
for claims remain the author's.

\begin{thebibliography}{99}
\bibitem{ASV10} O.~Ahmadi, I.~E.~Shparlinski, J.~F.~Voloch,
\emph{Multiplicative order of Gauss periods}, Int.~J.~Number Theory
\textbf{6} (2010), 877--882; arXiv:0707.4034.
\bibitem{Breuer21} F.~Breuer, \emph{Multiplicative orders of Gauss
periods and the arithmetic of real quadratic fields}, Finite Fields
Appl.\ \textbf{73} (2021), 101848.
\bibitem{Carlitz53} L.~Carlitz, \emph{Distribution of primitive roots in
a finite field}, Quart.~J.~Math. Oxford (2) \textbf{4} (1953), 4--10.
\bibitem{CHS25} L.~Cagliero, A.~Herman, and F.~Szechtman,
\emph{Artin--Schreier towers of finite fields}, Finite Fields Appl.\
\textbf{106} (2025), 102606; arXiv:2405.10159.
\bibitem{Conway76} J.~H.~Conway, \emph{On Numbers and Games}, Academic
Press, 1976.
\bibitem{DFRR19} L.~De~Feo, H.~Randriam, \'E.~Rousseau, \emph{Standard
lattices of compatibly embedded finite fields}, Proc.\ ISSAC 2019;
arXiv:1906.00870.
\bibitem{DiMuro} J.~DiMuro, \emph{On $\mathrm{On}_p$}, arXiv:1108.0962.
\bibitem{GV95} S.~Gao, S.~A.~Vanstone, \emph{On orders of optimal basis
generators}, Math.\ Comp.\ \textbf{64} (1995), 1227--1233.
\bibitem{Iwasawa68} K.~Iwasawa, \emph{On explicit formulas for the norm
residue symbol}, J.~Math. Soc. Japan \textbf{20} (1968), 151--165.
\bibitem{KS98} E.~Kaltofen, V.~Shoup, \emph{Subquadratic-time factoring of
polynomials over finite fields}, Math.\ Comp.\ \textbf{67} (1998),
1179--1197.
\bibitem{Lenstra77} H.~W.~Lenstra, Jr., \emph{On the algebraic closure of
two}, Indag.\ Math.\ \textbf{39} (1977), 389--396;
\url{https://pub.math.leidenuniv.nl/~lenstrahw/PUBLICATIONS/1977e/art.pdf}.
\bibitem{Lenstra78} H.~W.~Lenstra, Jr., \emph{Nim multiplication},
S\'eminaire de Th\'eorie des Nombres de Bordeaux, 1977--1978.
\bibitem{MM05} R.~C.~Mullin, A.~Mahalanobis, \emph{Dickson bases and
finite fields}, CACR Technical Report 2005--03, University of Waterloo,
2005.
\bibitem{MRS20} T.~M\'esz\'aros, L.~R\'onyai, T.~Szab\'o,
\emph{Singer difference sets and the projective norm graph},
arXiv:1908.05591 (2019).
\bibitem{Meinardus54} G.~Meinardus,
\emph{Asymptotische Aussagen \"uber Partitionen},
Math.\ Z.\ \textbf{59} (1954), 388--398.
\bibitem{OEIS} OEIS Foundation Inc., entry A380496,
\emph{The On-Line Encyclopedia of Integer Sequences},
\url{https://oeis.org/A380496}.
\bibitem{CGSuite} A.~N.~Siegel, \emph{CGSuite} (software),
\texttt{NimFieldCalculator},
\url{https://github.com/aaron-siegel/cgsuite}.
\bibitem{Peeters} D.~Peeters, \emph{transfinite-nim-calculator}
(software),
\url{https://github.com/DjangoPeeters/transfinite-nim-calculator}.
\bibitem{Popovych18} R.~Popovych, \emph{Multiplicative orders of
elements in Conway's towers of finite fields}, Algebra Discrete Math.
\textbf{25} (2018), 137--146; arXiv:1509.01958.
\bibitem{vzGS92} J.~von~zur~Gathen, V.~Shoup, \emph{Computing Frobenius
maps and factoring polynomials}, Comput.\ Complexity \textbf{2} (1992),
187--224.
\bibitem{Wiedemann88} D.~Wiedemann, \emph{An iterated quadratic
extension of $GF(2)$}, Fibonacci Quart.\ \textbf{26} (1988), 290--295.
\bibitem{Washington97} L.~C.~Washington, \emph{Introduction to
Cyclotomic Fields}, 2nd ed., Graduate Texts in Mathematics 83,
Springer, 1997.
\end{thebibliography}

\end{document}
